\documentclass[11pt]{article}
\usepackage[margin=1in]{geometry}
\usepackage[T1]{fontenc}\usepackage[utf8]{inputenc}\usepackage{lmodern}
\usepackage{amsmath,amssymb,amsthm,mathtools,booktabs,microtype}
\usepackage[colorlinks=true,linkcolor=blue,urlcolor=blue,citecolor=blue,hypertexnames=false]{hyperref}
\hypersetup{pdftitle={Learning narrow approximations with Gaussian Bayesian neural networks},pdfauthor={ChatGPT-6 Astra; Claude Opus 5.5; Kaarel Hänni; Dmitry Vaintrob}}
\newtheorem{theorem}{Theorem}[section]\newtheorem{lemma}[theorem]{Lemma}
\newtheorem{corollary}[theorem]{Corollary}\newtheorem{proposition}[theorem]{Proposition}
\theoremstyle{remark}\newtheorem{remark}[theorem]{Remark}
\newcommand{\E}{\mathbb E}\newcommand{\R}{\mathbb R}
\newcommand{\KL}{\operatorname{KL}}\newcommand{\clip}{\operatorname{clip}}
\title{Learning narrow approximations with\\Gaussian Bayesian neural networks}
\author{ChatGPT-6 Astra \and Claude Opus 5.5 \and Kaarel H\"anni \and
Dmitry Vaintrob\thanks{Authors are listed in alphabetical order.}}
\date{17 November 2026}
\begin{document}\maketitle
\begin{abstract}
We study whether a large Bayesian neural network predicts as well as
smaller neural approximations. The learner has one prescribed finite
architecture, independent centered Gaussian weights, a raw Gaussian
likelihood, and the ordinary final posterior mean using all observations.
For shallow ReLU networks and for bounded monotone activations such as
tanh, sigmoid, and hardtanh, any input law on the cube and a finite
label second moment suffice for a high-probability oracle with exact
leading comparator coefficient one. At fixed dimension, coefficient
cap, and confidence, the remainder is $\widetilde O(C/\sqrt n)$
against width-$C$ comparison networks, without an output cap.
The finite-dimensional resource is an explicit path mass or quadratic
parameter energy. Both families cover every learner width that can hold the comparison
network; a constant neuron implementing an output bias is counted in
that width. Deterministic schedules independent of the unknown teacher
compete with width-$n^{1/4}$ comparisons at excess
$\widetilde O(n^{-1/4})$, already at comparable learner widths,
simultaneously over all mean-field-or-stronger Gaussian scales.
The finite-variance conclusions pay a polynomial confidence cost and
make no expected-risk claim under arbitrary heavy-tailed labels.
Expected-risk versions under marginally sub-Gaussian labels are
also proved for ReLU at every eligible width and for the bounded-activation
variant with $N\lambda\ge c$.

Separate results remove every original comparison-coefficient bound.
At suitable fixed learner depth, tanh learners compete with
polynomial-width ReLU networks on Boolean inputs with arbitrary input
law, or on continuous inputs under density or slab regularity.
In one dimension, a shallow ReLU learner with $N\ge n$ and input
density bounded above and below has excess
$\widetilde O(C^{4/7}n^{-4/7}+C^{4/5}n^{-3/5})$ for a declared width
budget. Schedules independent of $C$ also give polynomial growing-width
oracles. These coefficient-free results allow marginally sub-Gaussian
labels. All conclusions use specified likelihood calibrations; a
constant-input example shows that fixed physical observation variance
can fail as width grows. There is no architecture mixture, prior support
restriction, or clipping of actual observations or predictions.
The sufficient constants may be large, and efficient posterior
computation is not established.
\end{abstract}

\section{The learning rule and the comparison question}
\label{sec:model}
An approximation guarantee asks whether a large Bayesian network
learns about as well as the best much smaller neural network, even when
the comparison network does not fit the data exactly. We prove several
such guarantees for the ordinary raw posterior mean under independent
Gaussian weights. The simplest shallow result allows any input law and
finite-variance labels, with coefficient magnitudes counted explicitly.
Other results remove all original comparison coefficients under the
stated input regularity assumptions. The same Gaussian prior family is
used throughout; no heavy-tailed or hierarchical sparsity prior is
substituted.

For the strongest coefficient-controlled bounded-activation theorem,
read Section~\ref{sec:all-width-finite-variance}; for a direct ReLU
learner on arbitrary input laws, read Section~\ref{sec:all-width-relu}.
The preceding sections provide their comparison and capacity lemmas.
The fixed-depth coefficient-free construction is stated in the next
sections, and the one-dimensional ReLU analysis gives sharper rates
under its two-sided density assumption. Each theorem specifies its
activation, likelihood calibration, and learner-width requirement.

Fix hidden depth $H$ and hidden width $N$. On $[-1,1]^d$, with
$\bar x=(1,x)$, write
\[
 h^1=\tanh(W^1\bar x),\quad
 h^\ell=\tanh(W^\ell h^{\ell-1})\ (2\le\ell\le H),\quad
 z_w=v^\top h^H.
\]
There are no internal biases. A constant hidden wire in each compiler
implements affine offsets. For any $0<\nu\le1$, all weight coordinates
are independent centered Gaussians with variances
\[
 \frac\nu d\quad\hbox{in the first layer},\qquad
 \frac\nu N\quad\hbox{in hidden-to-hidden layers},\qquad
 \frac\nu{N^2}\quad\hbox{in the readout}.
\]
The prior is denoted by $\mu_{N,\nu}$. Canonical scaling exponent
$\gamma\ge1$ means $\nu=N^{1-\gamma}$; $\nu=1$ is mean field.
The architecture and prior are prescribed before observing the sample.

For IID observations $S=((X_i,Y_i))_{i=1}^n$, the actual posterior is
\begin{equation}
 d\Pi_S(w)=Z_S^{-1}
 \exp\left[-\frac N{2\lambda\nu n}
                  \sum_{i=1}^n(z_w(X_i)-Y_i)^2\right]d\mu_{N,\nu}(w).
 \label{eq:actual-posterior}
\end{equation}
Its Gaussian observation variance is $n\lambda\nu/N$. The general fixed-depth theorem selects $\lambda$ by the evidence
calibration below; the shallow results specify deterministic scales. The returned
prediction is $\bar z_S(x)=\int z_w(x)d\Pi_S(w)$, and
$R(f)=\E(f(X)-Y)^2$ denotes squared test risk. No output is clipped
in the likelihood or returned prediction.

\begin{proposition}[Why the physical likelihood scale matters]
\label{prop:fixed-variance-obstruction}
Keep the tanh model above, with no output bias, but use a fixed physical
Gaussian observation variance $\sigma^2>0$. If $X=x_0$ and $Y=1$
almost surely, then after $n$ observations its ordinary posterior mean
satisfies
\[
 R(\bar z_S)\ \ge\
 \left(\frac{\sigma^2N}{\sigma^2N+n\nu}\right)^2.
\]
The posterior-average risk obeys the same lower bound. A width-one
bounded-coefficient tanh network realizes the target at every depth.
\end{proposition}
\begin{proof}
Conditional on the hidden weights, the output at $x_0$ is a centered
Gaussian of variance
$s^2=(\nu/N^2)\sum_j h_j^H(x_0)^2\le\nu/N$. Gaussian conditioning
on the repeated label one gives conditional posterior mean
$n s^2/(\sigma^2+n s^2)\le n\nu/(\sigma^2N+n\nu)$.
Averaging over the hidden posterior preserves this bound; Jensen's
inequality also gives the posterior-average risk statement.
For the teacher, set $c=\tanh(1)$, take first-layer input weights zero
and bias one, each subsequent scalar hidden weight $1/c$, and readout
$1/c$. Its output is exactly one and all coefficients are less than two.
\end{proof}
Thus a fixed physical variance fails when $N/n\to\infty$, even for
this constant target and $\nu\le1$. This is a statement about that
likelihood schedule, not about all calibrations. The factor
$n\lambda\nu/N$ in the positive results is material.

Our comparison class consists of ordinary ReLU networks with fixed
hidden depth $h$, arbitrary finite real coefficients and biases, and
width at most $C$. For continuous inputs a regularity assumption is
needed to remove all coefficient magnitudes from a quantitative bound:
a huge coefficient can otherwise separate two arbitrarily close input
points. Boolean inputs need no input-distribution restriction.

The general fixed-depth argument uses an existing Gaussian raw-output
oracle, stated with its full learning rule in the next section.
The compiler constructions remove the original comparison coefficient
cap while preserving raw outputs and fixed learner depth. The shallow
one-dimensional results provide a direct statistical analysis using
localized variation bounds and empirical midpoint comparisons.
A second direct argument applies in arbitrary dimension under any
input law on the cube. For arbitrary shallow teachers with a coefficient cap it gives
$\widetilde O(\sqrt{(d+1)A_n/n})$, with an explicit variation-sensitive
resource $A_n$. Randomly replacing a subset of the posterior neurons
by comparison neurons supplies the empirical midpoint inequality;
bounded ridge-feature capacity then transfers it to test risk.
A weighted second-moment argument permits arbitrary finite-variance
labels in this confidence theorem.
A direct cloned comparison and a variable statistical tolerance then
extend the coefficient-controlled rate to every $N\ge C$, without
a condition on $N\lambda$. Its explicit quadratic teacher resource
differs from the variation-weighted resource of the wider-learner proof.
The earlier wide-learner hardtanh instance also competes with
coefficient-capped ReLU teachers at the same depth. A separate appendix uses one-dimensional smoothness to
sharpen the tanh rate. A first ReLU argument compresses empirical midpoint values into at
most $n$ neurons. A direct balanced comparison and a joint
empirical-loss/energy cutoff subsequently remove the requirement
$N\ge n$: the coefficient-controlled oracle only requires the teacher
to fit in the learner. Neither argument needs input regularity.
The sharper one-dimensional tanh appendix retains a teacher-output
bound; the arbitrary-input-law result removes it by a risk-weighted
multiplier estimate. Neither result restricts prior or posterior outputs.

\section{The finite Gaussian raw-output oracle used here}
\label{sec:statistical-interface}
We state the statistical input precisely, so the learning rule and the
scope of the reduction do not depend on conventions from another paper.
The theorem is proved in~\cite{gaussian}; its extension to unbounded
labels is proved in~\cite{tails}. Neither result uses an architecture
mixture or a support restriction.

For $H\ge2$, put
\[
 v_H=\max\{4H,H^2+2H-2\},\qquad \rho_H=2H,\qquad r_H=v_H/\rho_H.
\]
Let $\omega=N/\nu$, and let
$M_n=\max\{1,s_Y\sqrt{8\log(8n/\delta)}\}$, where
$P(|Y|>t)\le2e^{-t^2/(2s_Y^2)}$.
For the normalized observable $z_w/M_n$ and labels $Y_i/M_n$, define
\[
 F'_S(\lambda')=-\frac{2\lambda'}\omega
 \log\int\exp\left[-\frac\omega{2\lambda'}
       \widehat R_S(z_w/M_n,Y/M_n)\right]d\mu_{N,\nu}(w).
\]
Set
\[
 b_n=c_H\log(c_Hn/\delta),\qquad
 \theta_n=\frac{c_H(d+1)^2
       [\log(c_Hn(d+1)/\delta)]^{12H+30}}n.
\]
Here sufficiently large constants $c_H$ depend only on the fixed depth.
If $\theta_n b_n^{v_H}5^{v_H-\rho_H}\le1$, choose the unique positive
root
\begin{equation}
 (\lambda')^{v_H}=\theta_n
       [2F'_S(\lambda')+b_n\lambda']^{v_H-\rho_H}.
 \label{eq:calibration}
\end{equation}
Otherwise use the bounded-risk fallback $\lambda'=2\omega$.
The actual raw likelihood parameter is $\lambda=M_n^2\lambda'$.
Its observation variance is $n\lambda\nu/N$.
All observed labels and all network outputs remain unchanged in that
likelihood. Only its scalar variance is calibrated.

\begin{lemma}[Imported raw Gaussian comparison theorem]
\label{lem:gaussian-interface}
Fix a tanh comparison network $q$ of depth $H$, width $m$, and
coefficient cap $B\ge1$, with an input constant and no internal biases.
For a sufficiently large $c_H$, set
\[
 A_n=c_H(d+m+1)^2(B+1)^2
       \log^2\frac{c_Hn(d+m+1)(B+1)}\delta.
\]
Put $u_n=2s_Y(\delta/(8n))^2$ and
\[
 r_{q,n}=(\sqrt{R(q)}+u_n)^2,\qquad
 e_n=c_H\left[A_nM_n^{2/r_H}r_{q,n}^{1-1/r_H}\theta_n^{1/v_H}
       +A_n^{r_H}M_n^2\theta_n^{1/\rho_H}\right]+M_n^2n^{-3}.
\]
With probability at least $1-\delta$, simultaneously for every
$N\ge m$ and every $0<\nu\le1$, the raw posterior under
\eqref{eq:calibration} obeys
\begin{equation}
 \Pi_SR(z)\le R(q)+e_n+4u_n\sqrt{R(q)+e_n}+4u_n^2.
 \label{eq:gaussian-interface}
\end{equation}
The ordinary raw posterior mean satisfies the
same bound. The calibration does not know $q,m,B$ or its risk.
\end{lemma}

The displayed $A_n$ is a convenient upper bound for the sharper
comparison action in~\cite{gaussian}. Its relevant consequence at
fixed $d,H,s_Y,\delta$ is
\begin{equation}
 \Pi_SR(z)\le R(q)+
 \widetilde O_H\left((mB)^2(R(q)+1)^{1-1/r_H}n^{-1/v_H}
              +(mB)^{2r_H}n^{-1/\rho_H}
\right)+o(1),
 \label{eq:gaussian-interface-coarse}
\end{equation}
where harmless additive constants in $mB$ are understood.

For clarity, the unbounded-label transfer does not change the prior.
In the proof replace $Y$ by $Y_{M_n}=\clip_{[-M_n,M_n]}Y$; with
probability at least $1-\delta/2$ every training label already equals
this value. Normalize the observable and labels by $M_n$ and apply
the bounded-label theorem. Dictionary amplitudes and approximation
errors only decrease under this normalization, while all parameter
energies and comparison entropies stay unchanged. The physical
likelihood is exactly the one above. Finally
\[
 \|Y-Y_{M_n}\|_2\le
 2s_Ye^{-M_n^2/(4s_Y^2)}
\]
and Minkowski's inequality transfer posterior-averaged test risk back
to the original labels. This proves the asserted powers and the
negligible-tail statement. The exact transfer formula, if needed at
small $n$, replaces $R(q)$ inside the remainder by
$(\sqrt{R(q)}+u_n)^2$ and adds
$4u_n\sqrt{R(q)+e_n}+4u_n^2$, where
$u_n=2s_Y(\delta/(8n))^2$ and $e_n$ is the displayed remainder.

\section{Compilers whose actual outputs remain bounded}
\subsection{Boolean inputs}
\begin{lemma}\label{lem:raw-Boolean-tanh}
There is a numerical $b\ge2$ such that the following holds.  Let $g$ be any
ReLU network of width $C\ge1$ and hidden depth $h\ge1$, with arbitrary
finite real coefficients and affine biases.  On $\{-1,1\}^d$, for every
$0<\epsilon<1$ and every $H\ge b(h+1)$ there is a tanh network $q$ of
hidden depth $H$, with a linear raw readout, satisfying
\[
 \sup_x|q(x)-\operatorname{clip}(g(x),[-1,1])|\le\epsilon,
 \qquad |q(x)|\le1+\epsilon,
\]
and
\[
 m,B\le\left[4(d+1)(C+1)^h
              +\lceil\log_2(32/\epsilon)\rceil+2\right]^b.
\]
The coefficient and width bounds do not grow with the additional learner
depth $H$.  On the entire input cube the same raw implementation satisfies
$|q|\le4$, because its signed-bit readout has coefficient sum below two
and constant offset of magnitude below two.
\end{lemma}
\begin{proof}
Use the coefficient-free unweighted-majority circuit compiler of~\cite{boolean}.  Its depth
is $O(h+1)$, its size including wires is polynomial in the bracketed
quantity, and a linear readout of its output bits, with absolute
coefficient sum less than four, approximates the clipped target to
$\epsilon/2$.  For a threshold gate with $r$ inputs and threshold $j$,
the signed-bit affine input is
\[
 \ell(\sigma)=\sum_{i=1}^r\sigma_i-(2j-r-1).
\]
It has margin at least one on exact signs.  If the circuit size is $S$,
choose $\delta=\min\{\epsilon/16,1/[8(S+1)]\}$ and
$k=\log(2/\delta)$.  Replacing each gate by $\tanh(k\ell)$ preserves
signed-bit error at most $\delta$ at every layer: incoming error changes
$\ell$ by at most $1/8$, and $2e^{-7k/4}\le\delta$.
Copying signed bits with $\tanh(k\sigma)$ restores the same accuracy, so
arbitrary further depth causes no accumulation.

An exact constant $c_0=\tanh(1)$ is initialized from the input
constant and carried by the weight $1/c_0$. Every later affine offset
is an edge from this wire; the learner has no native internal biases.
Converting the Boolean readout to signed bits gives a nonconstant
coefficient sum less than two, so its raw reconstruction error is at
most $2\delta$.  The triangle inequality proves the claimed raw output
error before any clipping.  The width is $S+1$, gate coefficients are
at most $k$, and bias/constant-wire coefficients at most $(S+1)k/c_0$.
Enlarging a numerical exponent $b$ gives all bounds.
\end{proof}

\subsection{A raw continuous tanh compiler}
Assume $P_X$ has density at most $\kappa$ relative to Lebesgue measure on
$[-1,1]^d$, and write $\bar\kappa=\max\{1,\kappa\}$.
\begin{lemma}\label{lem:raw-continuous-tanh}
Let $g$ be a width-$C$, hidden-depth-$h$ ReLU network with arbitrary
coefficients and biases.  For $0<\zeta<1$, put
$Q_0=\bar\kappa(C+1)^h/\zeta$ and $r_d=\max\{d,2\}$.
At every prescribed hidden depth $H\ge2$ there is a tanh network $q$
with raw linear readout such that
\[
 \E|q(X)-\operatorname{clip}(g(X),[-1,1])|\le c_d\zeta,
 \qquad \sup_{[-1,1]^d}|q|\le2^d+1,
\]
and
\[
 m\le c_d Q_0^d,\qquad
 B\le c_d Q_0^{r_d}\log(c_dQ_0/\zeta).
\]
The constants are independent of $H$ and of the original coefficients.
\end{lemma}
\begin{proof}
The clipped target has variation at most $2(C+1)^h$ on each coordinate
line.  A randomly shifted grid of mesh
$\Delta=c_d^{-1}\zeta/[\bar\kappa(C+1)^h]$ therefore has a deterministic
shift for which the boxwise-constant approximation $p$ has
$\E|p(X)-\operatorname{clip}(g(X))|\le\zeta/8$.
This is the coordinate bounded-variation translation inequality followed
by averaging the grid shift.  There are at most
$M_0\le c_dQ_0^d$ relevant boxes.  Choose a boundary-strip width
$w=c_d^{-1}\zeta\Delta/\bar\kappa\le\Delta/4$, so the total strip
probability is at most $\zeta/32$ and $1/w\le c_dQ_0^2$.

For each box $B_0=\prod_j[a_j,b_j)$, first form signed membership values
$\tanh(k(x_j-a_j))$ and $\tanh(k(b_j-x_j))$, with
$k=(2w)^{-1}\log(32d)$.  Their signed-AND input is their sum minus
$2d-1$.  Applying a second tanh at gain
$\ell=c\log(32M_0/\zeta)$ gives a signed indicator $v_{B_0}$;
write $I_{B_0}=(1+v_{B_0})/2$.  Away from all boundary strips the raw
readout $q=\sum_{B_0}p(B_0)I_{B_0}$ differs from $p$ by at most
$\zeta/32$.

A global raw envelope is also available.  If $x$ lies outside the box
expanded by $w$ in each coordinate, at least one membership value is
at most $-1+1/(16d)$.  Since all other values are at most one, the AND
input is at most $-1+1/(16d)$.  Thus
$I_{B_0}(x)\le\zeta/(16M_0)$ after enlarging $c$.  At most $2^d$
expanded grid boxes contain a given point because $w\le\Delta/4$.
Consequently $\sum_{B_0}I_{B_0}(x)\le2^d+\zeta/16$ everywhere, which
proves the stated raw envelope.  The bounded bad-set probability and
the grid's $L^1$ error now give the displayed $L^1$ bound, and also a
squared-error bound of the same order, with a constant
depending only on $d$.

The first-layer coefficients are $O_d(1/w)$, second-layer coefficients
are $O_d(\log(M_0/\zeta))$, and the constant readout coefficient is at
most $M_0/(2c_0)$ when implemented by the exact wire $c_0=\tanh(1)$.
All affine offsets, including the constant readout term, use this
exact wire. These give the displayed cap. To reach any larger $H$, copy the signed
box indicators with $\tanh(\ell v)$.  Outside a box's expanded region
its negative sign remains accurate, and away from boundary strips all
signs remain accurate.  Errors are restored at each copy, not accumulated.
The global envelope and approximation bounds therefore persist at every
padded depth.
\end{proof}


\subsection{Slab regularity, including sphere inputs}
Assume that for fixed $\kappa\ge1$ and $0<\alpha\le1$,
\begin{equation}
 P_X\{|u\cdot X-t|\le r\}\le\kappa r^\alpha
 \quad(\|u\|_2=1,\ t\in\R,\ r>0).
 \label{eq:slab}
\end{equation}
This does not require an ambient Lebesgue density.

\begin{lemma}[A fixed-depth raw compiler under slab regularity]
\label{lem:slab-box}
For every width-$C$, hidden-depth-$h$ ReLU teacher with arbitrary finite
coefficients, every $0<\zeta<1/4$, and every fixed learner depth $H\ge2$,
there is a raw tanh network $q$ with an input constant and no internal
biases such that
\[
 \E|q(X)-\clip(g(X))|\le c_d\zeta,
 \qquad \sup_{[-1,1]^d}|q|\le2^d+1.
\]
Writing
\[
 R=\frac{\kappa(C+1)^{dh+1}}{\zeta^2},\qquad
 a_{d,\alpha}=\max\{d/\alpha,\ 1/\alpha+1/\alpha^2\},
\]
its width and coefficient cap satisfy
\[
 m\le c_{d,h,\alpha}R^{d/\alpha},\qquad
 B\le c_{d,h,\alpha}R^{a_{d,\alpha}}
                    \log(c_{d,h,\alpha}R/\zeta).
\]
The constants do not depend on original teacher coefficients or on
additional signed-bit padding layers.
\end{lemma}
\begin{proof}
There are at most $A=(C+1)^{dh}$ full-dimensional affine activation
cells: each layer cuts an existing cell by at most $C$ hyperplanes,
creating at most $(C+1)^d$ subcells. Include all their activation
boundary planes, at most $hCA$ of them. On each final cell the output
is affine. Include also its proper level hyperplanes at levels
$-1+2j/L$, $0\le j\le L$, where $L=\lceil8/\zeta\rceil$.
Constant formulas contribute no level planes. The full collection has
at most
\[
 F=(hC+L+1)A\le c_h(C+1)^{dh+1}/\zeta
\]
planes. Degenerate identically zero gates add no boundary.

Choose an axis-aligned grid of mesh
$\Delta=c_d(\zeta/(\kappa F))^{1/\alpha}$, with a sufficiently
small fixed constant, and assign each box the value of $\clip(g)$ at
its lower corner. Corners may lie slightly outside the input cube;
the same affine-region count applies there. If the segment from $x$
to its box corner meets one of the listed planes, then $x$ lies within
$\sqrt d\Delta$ of that plane. By \eqref{eq:slab}, this exceptional
set has probability at most $\kappa F(\sqrt d\Delta)^\alpha\le\zeta/8$.
On its complement the clipped affine output stays within one interval
between adjacent output levels, so its variation is at most
$2/L\le\zeta/4$. Hence the boxwise-constant approximation has $L^1$
error at most $c\zeta$.

There are at most $M_0\le c_d\Delta^{-d}$ relevant boxes. Choose
\[
 w=\min\left\{\Delta/4,
       (\zeta\Delta/(c_d\kappa))^{1/\alpha}\right\}.
\]
The union of strips of width $w$ around coordinate grid boundaries
has probability at most $c_d\kappa\Delta^{-1}w^\alpha\le\zeta/16$.
The soft tanh box indicators from the preceding proof apply without
change. Their raw sum has global magnitude at most $2^d+1$;
their first-layer coefficients are $O_d(1/w)$, and their remaining
coefficients are bounded by
$O_d(M_0+\log(M_0/\zeta))$. They approximate the boxwise constant
function in $L^1$ to error $c_d\zeta$.

Finally $\Delta^{-1}\le c_{d,h,\alpha}R^{1/\alpha}$ and
\[
 w^{-1}\le c_{d,h,\alpha}
       \left[R^{1/\alpha}
              +(\kappa/\zeta)^{1/\alpha}R^{1/\alpha^2}\right]
 \le c_{d,h,\alpha}R^{1/\alpha+1/\alpha^2}.
\]
These inequalities give the displayed width and coefficient bounds.
The signed copying gates preserve both the raw envelope and the
accuracy at every larger prescribed depth.
\end{proof}

Uniform measure on the unit sphere in dimension $d\ge3$ satisfies
\eqref{eq:slab} with $\alpha=1$ and $\kappa=4\sqrt d$. On the circle,
$\alpha=1/2$ and $\kappa=4$ suffice. Indeed the first-coordinate density
on the sphere is proportional to $(1-t^2)^{(d-3)/2}$; it is at most
$\sqrt d$ for $d\ge3$, while the circle's endpoint singularity has
integral $O(\sqrt r)$ over an interval of length $2r$. A bounded density
relative to surface measure multiplies $\kappa$ by that density bound.

\section{Removing the original teacher coefficient magnitudes}
\label{sec:width-oracles}
Let $\mathcal G_{C,h}$ be all raw scalar-output ReLU networks with
hidden width at most $C$, hidden depth $h$, arbitrary finite real
coefficients and affine biases. No bound on their output or coefficient
magnitudes is imposed. The zero function belongs to this class.

\begin{lemma}[Risk transfer through a proof-only bounded compiler]
\label{lem:compiler-risk}
Suppose $f=\clip_{[-M,M]}g$, $M\ge1$, and a deterministic network $q$
satisfies $|q|\le D M$ and
$\E|q-f|\le D M\zeta$, where $D\ge1$ is fixed. Then
\[
 R(q)\le R(g)+C_D M^2\zeta
                   +C_D\E[Y^2\mathbf1_{\{|Y|>M\}}].
\]
\end{lemma}
\begin{proof}
The pointwise projection inequality
\[
 (\clip_{[-M,M]}a-y)^2\le(a-y)^2+(|y|-M)_+^2
\]
follows by maximizing the quadratic excess over $a$ outside the clipping
interval. On $|Y|\le M$, the loss difference between $q$ and $f$ is at
most $(D+3)M|q-f|$. On its complement, it is at most
$C_D M(M+|Y|)\le C_DY^2$. Integrate and add the projection inequality.
\end{proof}

Use $M=M_n$ from Section~\ref{sec:statistical-interface}. Compile
$g/M_n$ to a raw tanh network approximating its clipped output, and
multiply its final readout by $M_n$. Thus the compiled width stays $m$
and its coefficient cap grows by only the logarithmic factor $M_n$.
The lemma gives
\[
 R(q)\le R(g)+\widetilde O(\zeta)+o(n^{-2})
\]
in every continuous case, and with $\zeta$ replaced by the uniform
accuracy $\epsilon$ in the Boolean case. This is only a proof step:
the actual learner uses its unrestricted raw output.

\begin{theorem}[Coefficient-free Gaussian raw oracles]
\label{thm:width-only}
Fix $d,h,H,s_Y,\delta$ and use the Gaussian prior and raw likelihood
specified in Sections~\ref{sec:model}--\ref{sec:statistical-interface}.
The following statements hold, with probability at least $1-\delta$,
simultaneously over every learner width large enough to contain the
stated compiler and every $0<\nu\le1$.

\begin{enumerate}
\item On Boolean inputs $\{-1,1\}^d$ with any input distribution, assume
$H\ge b(h+1)$, where $b$ is the numerical compiler exponent.
At compiler accuracy $\epsilon=n^{-2}$, both width and coefficient
cap before the final $M_n$ rescaling are at most
\[
 m,B\le[4(d+1)(C+1)^h+c\log(n+2)]^b.
\]
The ordinary raw posterior mean and posterior-averaged risk are bounded
by $\inf_{g\in\mathcal G_{C,h}}R(g)+o(1)$ whenever
$C\le n^u$ and $4bhu<1/v_H$.

\item On a continuous cube law with Lebesgue density at most $\kappa$,
assume $H\ge2$. Put $s_d=d+\max\{d,2\}$. For any positive $u,t$ with
\begin{equation}
 s_d(hu+t)<1/(2v_H),
 \label{eq:density-growth}
\end{equation}
take $C\le n^u$ and compiler accuracy $\zeta=n^{-t}$.
Then the same oracle holds, with an excess bounded by
\[
 \widetilde O\left(n^{-t}
       +n^{-1/v_H+2s_d(hu+t)}
       +n^{-1/\rho_H+2r_Hs_d(hu+t)}\right).
\]

\item Under slab regularity \eqref{eq:slab}, including the sphere laws
listed above, assume $H\ge2$ and set
\[
 s_{d,\alpha}=d/\alpha+
             \max\{d/\alpha,1/\alpha+1/\alpha^2\}.
\]
For any positive $u,t$ with
\begin{equation}
 s_{d,\alpha}[(dh+1)u+2t]<1/(2v_H),
 \label{eq:slab-growth}
\end{equation}
put $a=s_{d,\alpha}[(dh+1)u+2t]$.
At $C\le n^u$, $\zeta=n^{-t}$, the same oracle has excess
\[
 \widetilde O\left(n^{-t}+n^{-1/v_H+2a}
                              +n^{-1/\rho_H+2r_Ha}\right).
\]
\end{enumerate}
All constants may depend on the fixed quantities and the input
regularity constants, but never on original teacher coefficient
magnitudes or learner width. The prior, likelihood and predictions
are not clipped or support-restricted.
\end{theorem}
\begin{proof}
Choose a deterministic $n^{-2}$-minimizer of population risk over
$\mathcal G_{C,h}$. Its risk is at most $\E Y^2+n^{-2}$, because the
zero function is available. No attainment of the infimum is needed.
Compile this function after scaling by $M_n$, and use
Lemma~\ref{lem:compiler-risk}. Its compiled risk remains bounded by a
fixed constant plus a vanishing term. Apply
Lemma~\ref{lem:gaussian-interface} to this one deterministic compiled
network. Its event is common to all eligible widths and prior scales.
The calibration has no dependence on the selected teacher or compiler.

In the Boolean case, $mB\le\widetilde O(C^{2bh})$, so the two
statistical terms have powers $n^{-1/v_H+4bhu}$ and
$n^{-1/\rho_H+4bhu r_H}$, both negative under the stated condition.
In the density case, $mB\le\widetilde O(n^{s_d(hu+t)})$.
In the slab case, $mB\le\widetilde O(n^a)$.
Substitution in \eqref{eq:gaussian-interface-coarse} gives the
claimed exponents. The final readout rescaling, label-tail transfer
and confidence factors contribute only logarithms.
\end{proof}

\begin{corollary}[Polynomial input-dimension cost on Boolean inputs]
\label{cor:boolean-samples}
At fixed $h,H,s_Y$, with $H\ge b(h+1)$, a sufficient Boolean-input
sample bound for excess at most $\varepsilon\in(0,1)$ is
\[
 n\ge\widetilde O_{h,H,s_Y}\left(
 (d+1)^2[(d+1)(C+1)^h]^{4b v_H}\varepsilon^{-v_H}\right).
\]
Logarithms include confidence and the displayed resources. The learner
width need only exceed
$[4(d+1)(C+1)^h+c\log(n+2)]^b$, up to a numerical constant. The
conclusion holds simultaneously over every larger finite width and
all $0<\nu\le1$. In particular, the sufficient sample complexity is
polynomial jointly in teacher width and input dimension at fixed
teacher and learner depths, with no original coefficient bound.
\end{corollary}
\begin{proof}
The Boolean compiler gives $mB\le\widetilde O(K^{2b})$ for
$K=(d+1)(C+1)^h$. Hence the comparison action is at most
$\widetilde O_H(K^{4b})$, and $\theta_n$ has dimension factor
$(d+1)^2/n$. The displayed budget makes the first statistical remainder
at most $c\varepsilon$. It makes the second at most
$c\varepsilon^{r_H}\le c\varepsilon$, because $r_H\ge2$.
Compiler and label-tail errors are absorbed by enlarging the constant.
\end{proof}

\begin{corollary}[A literal polynomial-width comparison class]
\label{cor:literal-poly-width}
Each of the following is a sufficient concrete choice. For all
sufficiently large $n$, one event of probability at least $1-\delta$
covers every $N\ge n$ and every $0<\nu\le1$:
\[
 R(\bar z_S)\le\Pi_SR(z)
       \le\inf_{g\in\mathcal G_{\lfloor n^u\rfloor,h}}R(g)
                                   +\widetilde O(n^{-\eta}).
\]
One may take
\[
\begin{array}{c|c|c}
\text{input law}&u&\eta\\ \hline
\text{Boolean}&(16bhv_H)^{-1}&(2v_H)^{-1}\\
\text{bounded density}&(8hs_dv_H)^{-1}&(8s_dv_H)^{-1}\\
\text{slab regularity}&[16s_{d,\alpha}(dh+1)v_H]^{-1}
                                      &[32s_{d,\alpha}v_H]^{-1}.
\end{array}
\]
The required learner depths are those in Theorem~\ref{thm:width-only}.
The constants in the sample threshold may be large. These exponents
are explicit sufficient choices, not claims of optimality.
\end{corollary}
\begin{proof}
Use $t=\eta$ in the continuous cases and substitute in the theorem.
The resulting compiler width is $o(n)$ in each row, including every
logarithmic factor, so $N\ge n$ contains it for all large enough $n$.
In the Boolean row the statistical exponent is at least
$3/(4v_H)$; the displayed $1/(2v_H)$ leaves slack. The other rows'
statistical exponents are faster than their displayed compiler errors.
\end{proof}

\section{An arbitrary-input-law oracle for bounded monotone activations}
\label{sec:all-d-product-midpoint}
The random-subset midpoint construction also gives a simpler oracle in
arbitrary input dimension. It uses the bounded ridge-feature class rather
than a spatial smoothness assumption on the input distribution.

Fix a bounded, nonconstant, nondecreasing, globally Lipschitz activation
$\phi:\mathbb R\to\mathbb R$. Tanh, logistic sigmoid, and hardtanh are
examples. Constants below may depend on its fixed supremum $K$ and
Lipschitz constant. Let $X$ have any probability law on $[-1,1]^d$ and
let $|Y|\le1$ and $d\ge1$. Put $D=d+1$ and $\widetilde x=(x,1)$.
Let $a_0=d$ for the original first-layer convention; the alternative
$a_0=d+1$ is also covered. The learner is
\[
 f_w(x)=N^{-1}\sum_{j=1}^N G_j\phi(\theta_j\cdot\widetilde x),
 \qquad
 \theta_j\sim N(0,\nu I_D/a_0),\quad G_j\sim N(0,\nu),
 \quad0<\nu\le1,
\]
with all coordinates independent. Thus the actual readout coefficients
$G_j/N$ have variance $\nu/N^2$. Define
\[
 Q_G=N^{-1}\sum_jG_j^2,\qquad \omega=N/\nu,\qquad
 d\Pi_\lambda=Z_\lambda^{-1}
 e^{-\omega\widehat R(f_w)/(2\lambda)}d\mu.
\]
Here $R(f)=\mathbb E(f(X)-Y)^2$ and
$\widehat R(f)=n^{-1}\sum_i(f(X_i)-Y_i)^2$.
The physical Gaussian observation variance is $n\lambda\nu/N$.
Both the likelihood and the returned mean use the unrestricted raw output.

Let a deterministic comparison network be
\[
 g(x)=\sum_{i=1}^C c_i\phi(\vartheta_i\cdot\widetilde x),
 \qquad \|g\|_\infty\le1.
\]
The supremum is over the input cube. If $V=\sum_i|c_i|>0$, set
$p_i=|c_i|/V$, omitting zero coefficients, and define
\[
 J_g=a_0\sum_i p_i\|\vartheta_i\|_2^2,
 \qquad
 A_n(g)=1+V^2+J_g+D\log(e n(V+1)).                       \tag{AD1}
\]
For $g=0$ use $V=J_g=0$. The term $J_g$ is measured in the actual
fan-in-normalized incoming Gaussian energy. In particular, if every
teacher coefficient has absolute value at most $B\ge1$, then
\[
 A_n(g)\le C_\phi\{1+C^2B^2+D^2B^2+D\log(e n(CB+1))\}. \tag{AD2}
\]
An incoming row-norm cap $B_{\rm in}$ instead gives
$J_g\le a_0 B_{\rm in}^2$, independently of the teacher width.

\begin{lemma}[Capacity of the bounded ridge hull]
\label{lem:all-d-ridge-hull}
Let $\mathcal F_q$ be the class of all finite sums
$\sum_j v_j\phi(\theta_j\cdot\widetilde x)$ with
$\sum_j|v_j|\le\sqrt q$. On every list of $n$ inputs,
\[
 \operatorname{Rad}_n(\mathcal F_q)
 \le C_\phi\sqrt{Dq/n},\qquad
 \sup_{f\in\mathcal F_q}\|f\|_\infty\le K\sqrt q.        \tag{AD3}
\]
Every learner output belongs to $\mathcal F_{Q_G}$, without any bound
on incoming weights.
\end{lemma}
\begin{proof}
For a monotone activation, each subgraph test
$t<\phi(\theta\cdot\widetilde x)$ is a constant test or an affine
halfspace test in $(x,\phi^{-1}(t))$, with a generalized inverse at
flat portions. Strict or weak boundary choices give the same VC bound.
The ridge class therefore has VC-subgraph dimension at most $D+1$.
The standard VC-subgraph covering bound for a bounded real-valued class
is $N(\epsilon K,\mathcal H,L^2(P_n))\le(A/\epsilon)^{cD}$,
with universal $A,c$. Dudley's entropy integral gives worst-case
empirical Rademacher complexity $C K\sqrt{D/n}$. Adding negatives
only changes the constant. A signed convex hull does not increase
Rademacher complexity beyond its radius, proving (AD3).
Finally $N^{-1}\sum_j|G_j|\le\sqrt{Q_G}$.
\end{proof}

\begin{lemma}[A common relative-loss event]
\label{lem:all-d-relative}
Fix a deterministic $u$ with $\|u\|_\infty\le2$. There is a constant
$C_\phi$ such that, for $n\ge2$, putting
\[
 b_n=\frac{C_\phi(D+1)\log^4(e n/\delta)}n,              \tag{AD4}
\]
with probability at least $1-\delta$, simultaneously for all
$q\le n^{22}$ and $f\in\mathcal F_q$,
\begin{align}
 R(f)-R(u)\le{}&\widehat R(f)-\widehat R(u)
 +\tfrac16\|f-u\|_n^2+c\lambda(q+1)+C_{\phi,c}b_n/\lambda,
                                                                  \tag{AD5}
\end{align}
whenever $n^{-2}\le\lambda\le1$ and $\lambda\ge C_{\phi,c}b_n$.
The positive constant $c$ may be chosen arbitrarily small by enlarging
the constants in the likelihood floor and final remainder.
\end{lemma}
\begin{proof}
Write $h=f-u$. Apply the standard smooth-loss Rademacher inequality
for nonnegative squared loss (Srebro, Sridharan, and Tewari, 2010,
Theorem~1~\cite{sst}) to the class
$\mathcal F_q-u$, with response zero. Its Rademacher complexity is
at most $C_\phi\sqrt{D(q+1)/n}$ and its squared envelope is at most
$C_\phi(q+1)$. For any fixed small $\eta>0$, Young's inequality gives
\[
 (P_X-P_n)h^2\le\eta P_nh^2+C_{\phi,\eta}b_n(q+1).      \tag{AD6}
\]
Only this one-sided norm comparison is needed. The invoked smooth-loss
bound is the usual inequality
\begin{align*}
 P h^2-P_nh^2\le{}&C\sqrt{P_nh^2}
 \left[\log^{3/2}(2n)\operatorname{Rad}_n(\mathcal H)
       +\sqrt{B\log(1/\zeta)/n}\right]\\
 &+C\log^3(2n)\operatorname{Rad}_n(\mathcal H)^2
 +CB\log(1/\zeta)/n.
\end{align*}
where $|h|^2\le B$. Its hypotheses are satisfied by the bounded class
in (AD3).

For the other term, the multiplier $u(X)-Y$ is bounded by three.
Symmetrization, contraction, (AD3), and bounded-difference concentration
give
\[
 |(P-P_n)((u-Y)h)|\le C_\phi\sqrt{b_n(q+1)}.             \tag{AD7}
\]
This argument does not require the multiplier to be independent of $X$:
contraction is applied conditionally to the observed pairs $(X_i,Y_i)$.
Apply both inequalities at dyadic radii through $2n^{22}$, including
radius one, and allocate confidence equally to these $O(\log n)$
events. This explains the safe logarithmic factor in (AD4), and gives
one event for all intervening radii.

Use the exact identity
\[
 R(f)-R(u)-\widehat R(f)+\widehat R(u)
 =(P_X-P_n)h^2+2(P-P_n)((u-Y)h).
\]
Take $\eta=1/6$. The term $C b_nq$ is absorbed into a prescribed
small multiple of $\lambda(q+1)$ when $\lambda\ge Cb_n$.
Young's inequality bounds $C\sqrt{b_n(q+1)}$ by another such multiple
plus $C b_n/\lambda$.
\end{proof}

\begin{lemma}[Variation-weighted comparison packs]
\label{lem:all-d-comparison}
There is a one-pack comparison law $q_g$ such that its mean function
$g_\sigma$ is independent of $\nu$,
\[
 \|g_\sigma-g\|_\infty\le C_\phi n^{-20},\qquad
 \mathbb E_{q_g}G^2\le V^2+1,
 \qquad
 \nu\operatorname{KL}(q_g\Vert\mu_0)\le C_\phi A_n(g).  \tag{AD8}
\]
Consequently, whenever $N\lambda\ge C_\phi$,
\[
 \Pi_\lambda[\widehat R(f)-\widehat R(g_\sigma)]
 +\tfrac13\Pi_\lambda\|f-g_\sigma\|_n^2
 +c_\phi\lambda\Pi_\lambda Q_G
 \le C_\phi\lambda A_n(g).                             \tag{AD9}
\]
\end{lemma}
\begin{proof}
For $V>0$, choose index $i$ with probability $p_i$, then independently
\[
 \theta\sim N(\vartheta_i,\sigma^2I_D/a_0),\qquad
 G\sim N(V\operatorname{sign}c_i,\nu),\qquad
 \sigma=\frac{n^{-20}}{V+1}.
\]
Since $\|\widetilde x\|_2\le\sqrt D$, the incoming perturbation
changes a ridge in mean by at most
$\sqrt{D/a_0}\operatorname{Lip}(\phi)\sigma
 \le\sqrt2\operatorname{Lip}(\phi)\sigma$.
Thus the mean error is at most $C_\phi V\sigma$.
The readout second moment is exactly $V^2+\nu$.
Convexity and the Gaussian entropy formula give
\[
 \nu\operatorname{KL}(q_g\Vert\mu_0)
 \le\tfrac12\{J_g+V^2+D\sigma^2-D\nu
                     -D\nu\log(\sigma^2/\nu)\},
\]
which is bounded by $C_\phi A_n$ because $\nu\log\nu\le0$.
For $g=0$ simply take the native centered pack law.

Apply the random-subset replacement construction of
Lemma~\ref{lem:tanh-subset-midpoint}: keeping a uniformly
chosen $k=\lfloor N/2\rfloor$ posterior packs and replacing the others
by independent $q_g$ draws contracts relative entropy by $p=k/N$.
Its conditional variance is bounded by
\[
 K^2\left\{\frac{p(1-p)}{N-1}Q_G
             +\frac{1-p}{N}\mathbb E_{q_g}G^2\right\}.
\]
Gibbs optimality, the squared-loss midpoint identity, and division by
$1-p$ give
\begin{align*}
 &\Pi_\lambda[\widehat R(f)-\widehat R(g_\sigma)]
 +p\Pi_\lambda\|f-g_\sigma\|_n^2
 +\frac{2\lambda\nu}{N}\operatorname{KL}(\Pi_\lambda\Vert\mu)\\
 &\hspace{8mm}\le2\lambda\nu\operatorname{KL}(q_g\Vert\mu_0)
 +K^2\left\{\frac p{N-1}\Pi_\lambda Q_G
                  +\frac{\mathbb E_{q_g}G^2}N\right\}.
\end{align*}
Use only the readout-block Gaussian moment:
\[
 \operatorname{KL}(\Pi_\lambda\Vert\mu)
 \ge\frac\omega4\Pi_\lambda Q_G-\frac N2\log2.
\]
Absorb the $Q_G/N$ term using $N\lambda\ge C_\phi$, and use
$(V^2+1)/N\le C_\phi\lambda A_n$. This proves (AD9).
\end{proof}

\begin{theorem}[A sharp arbitrary-input-law Gaussian oracle]
\label{thm:all-d-sharp}
For every fixed comparison network $g$ above, with probability at least
$1-\delta$, simultaneously for all $0<\nu\le1$, all
$n^{-2}\le\lambda\le1$ satisfying $\lambda\ge C_\phi b_n$, and all
$N\lambda\ge C_\phi$,
\[
 R(\mathbb E_{\Pi_\lambda}f)\le\Pi_\lambda R(f)
 \le R(g)+C_\phi\left\{\lambda A_n(g)+\frac{b_n}{\lambda}\right\}
       +C_\phi n^{-10}.                                \tag{AD10}
\]
There is no density or anti-concentration condition on the input law.
No output truncation, prior restriction, or architecture mixture is used.
\end{theorem}
\begin{proof}
The prior empirical risk is at most $M=1+K^2$, since the readout is
centered and $\mathbb E_\mu f(x)^2\le K^2\nu/N$.
Thus Jensen gives $Z_\lambda\ge\exp[-\omega M/(2\lambda)]$.
The law of $\omega Q_G$ is $\chi^2_N$, so above a fixed constant
$\mu(Q_G>t)\le\exp(-c\omega t)$. Integrating the corresponding
posterior bound shows that the omitted risk, empirical loss and
quadratic-distance integrals on $Q_G>n^{22}$ are at most $C_\phi n^{-10}$,
uniformly for $\lambda\ge n^{-2}$, $N\ge1$ and $0<\nu\le1$.
For the finitely many smaller $n$, increase the deterministic lower
sample-size threshold by an activation-dependent constant.
All these quantities are bounded by $C_\phi(1+Q_G)$.

Use (AD5) with $u=g_\sigma$ and a constant $c$ smaller than half
$c_\phi$ in (AD9). Its norm is at most two for sufficiently large $n$.
Integrate over the posterior, restore the negligible tail, and apply
(AD9). The extra constant $c\lambda$ is absorbed by $C\lambda A_n$
since $A_n\ge1$. The remaining distance and readout-energy terms are nonnegative
and can be dropped. Finally $R(g_\sigma)\le R(g)+C_\phi n^{-20}$.
Jensen proves the mean-risk statement. The sample event is fixed by
$g_\sigma$ and the width/scale-independent ridge classes; $g_\sigma$
does not depend on $\nu$, so the advertised simultaneity is valid.
\end{proof}

\begin{corollary}[An explicit polynomial]
\label{cor:all-d-sharp-polynomial}
In the nontrivial range $A_n(g)b_n\le c_\phi$, choose
\[
 \lambda=C_\phi\sqrt{b_n/A_n(g)},\qquad
 N\ge C_\phi\sqrt{A_n(g)/b_n}.
\]
Then
\[
 R(\mathbb E_{\Pi_\lambda}f)\le\Pi_\lambda R(f)
 \le R(g)+C_\phi\sqrt{A_n(g)b_n}+C_\phi n^{-10}.         \tag{AD11}
\]
The scale uses a declared deterministic upper bound for $A_n$, not the
unknown risk or a prior over teacher sizes. At fixed $d,B$ this is
$\widetilde O(C/\sqrt n)$ against width-$C$ teachers, and a sufficient
sample size for excess $\varepsilon$ is
\[
 n\ge\widetilde O\left(
 (d+1)[C^2B^2+(d+1)^2B^2]\varepsilon^{-2}\right).        \tag{AD12}
\]
For bounded readout variation and bounded incoming row norm, the resource
in (AD1) is independent of the number of teacher neurons. The comparator is still required to have
bounded raw output. The same risk bound holds against the infimum over
a declared such class by selecting a deterministic population-risk
near-minimizer; no common event over all data-selected comparisons is
asserted. In the nontrivial range, every $N\ge C_\phi n$ satisfies the
width requirement, after enlarging constants and logarithmic factors.
\end{corollary}
\begin{proof}
Balance the two terms of (AD10). The nontrivial-range condition makes
$\lambda\ge C_\phi b_n$, after the constants are chosen in order, and
also $\lambda\le1$ when $b_n$ is sufficiently small. Equation (AD2)
gives (AD12). Finally $A_nb_n\le c$ implies
$\sqrt{A_n/b_n}\le C/b_n\le Cn$.
\end{proof}

\paragraph{An elementary bound without an output cap.}
Theorem~\ref{thm:uncapped-all-d} below removes the output cap without
worsening (AD11). The following simpler envelope-based argument gives
a coarser statement for any finite comparison network.
Let $F=\|g\|_\infty$, so $F\le K V$. In the nontrivial range, the same
argument gives
\[
 R(\mathbb E_{\Pi_\lambda}f)\le\Pi_\lambda R(f)
 \le R(g)+C_\phi\left\{\lambda A_n(g)
             +(F+1)^2b_n/\lambda\right\}+C_\phi(F+1)n^{-10}.
                                                               \tag{AD13}
\]
One may assume $F\le n^2$ to keep the same polynomial energy cutoff;
the nontrivial optimized range below automatically has this property.
Indeed the comparison mean is bounded by $F+1$. The squared-norm
argument now has envelope $C_\phi(q+(F+1)^2)$, and its remainder is
$C_\phi b_n(q+(F+1)^2)$. The bounded multiplier argument gives
$C_\phi(F+1)\sqrt{b_n(q+(F+1)^2)}$. Young's inequality bounds these by
a small multiple of $\lambda q$, plus
$C_\phi\lambda(F+1)^2+C_\phi(F+1)^2b_n/\lambda$.
Since $(F+1)^2\le C_\phi A_n$, the new constant penalty is absorbed
in $C_\phi\lambda A_n$; the midpoint comparison is unchanged.
The loss-smoothing error is $C_\phi(F+1)n^{-20}$ and the Gaussian
tail remains negligible at the stated polynomial cutoff.

Choosing $\lambda=C_\phi(F+1)\sqrt{b_n/A_n}$ gives excess
$\widetilde O_\phi((F+1)\sqrt{(d+1)A_n/n})$, provided this is
smaller than a sufficiently small fixed constant and
$N\lambda\ge C_\phi$. With entrywise coefficient cap $B$ and fixed
$d,B$, this is $\widetilde O(C^2/\sqrt n)$ without a comparison-output
restriction. It is weaker than (AD11), whose bounded-output premise
is what gives the better width polynomial.

An explicit teacher readout bias can be included as one extra constant
neuron: choose a fixed $t_0$ with $\phi(t_0)\ne0$, take incoming row
$(0,\ldots,0,t_0)$, and coefficient $b_{\rm out}/\phi(t_0)$.
It is counted in the variation and incoming-energy resources above.

\paragraph{Marginally sub-Gaussian labels.}
The bounded-label hypothesis can be replaced by the marginal condition
$\mathbb E\exp(Y^2/K_Y^2)\le2$ for a fixed $K_Y<\infty$.
This does not require a bounded conditional mean or a conditionally
sub-Gaussian residual. Keep the bounded-output comparator of (AD1), and set
\[
 B_n=\max\{1,K_Y\sqrt{40\log(e n/\delta)}\},\qquad
 b_n^{\rm sg}=C_\phi(B_n+1)^2(D+1)\log^4(e n/\delta)/n.
\]
Equations (AD10)--(AD11) hold with $b_n^{\rm sg}$ in place of $b_n$
and an $O_{\phi,K_Y}(n^{-10})$ negligible term. The comparison pack and
resource $A_n$ are unchanged. In particular the optimized excess is
\[
 \widetilde O_{\phi,K_Y}\bigl(\sqrt{(d+1)A_n/n}\bigr),               \tag{AD14}
\]
with the same condition $N\lambda\ge C_\phi$. As before, the scale must
satisfy $n^{-2}\le\lambda\le1$ and
$\lambda\ge C_\phi b_n^{\rm sg}$; the optimized statement concerns its
nontrivial range. No observed label is modified in the actual algorithm.

Here are the tail and risk details. In the proof alone define
$Y^{B_n}=\operatorname{clip}(Y,-B_n,B_n)$. With probability at least
$1-\delta/4$, every observed label equals its clipped value. The pairs
$(X,Y^{B_n})$ retain exactly the original input law. Apply the preceding
common statistical event to this clipped-label distribution. Its
multiplier estimate gains a factor $B_n+1$, which changes $b_n$ to
$b_n^{\rm sg}$. The midpoint identity and comparison-pack entropy do
not require bounded labels. On the no-exceedance event the actual
posterior, including its normalizer, is exactly the posterior appearing
in that analysis, and its native empirical risk is at most $B_n^2+K^2$.
The Gaussian readout-tail proof consequently gives
\[
 \Pi_\lambda Q_G\le C_\phi\{1+(B_n^2+1)/\lambda\}.
\]
The eligibility condition $b_n^{\rm sg}\lesssim1$ implies $B_n^2\lesssim n$;
thus $Q_G>n^{22}$ still has negligible posterior integrals uniformly
for $\lambda\ge n^{-2}$.

Most importantly, only a relative-risk transfer is needed:
\begin{align*}
 &[R_Y(f)-R_Y(g)]-[R_{Y^{B_n}}(f)-R_{Y^{B_n}}(g)]\\
 &\hspace{15mm}=-2\mathbb E[(Y-Y^{B_n})(f(X)-g(X))].
\end{align*}
The $Y^2$ terms cancel exactly. The marginal tail assumption gives
$\mathbb E|Y-Y^{B_n}|\le C K_Y(e n/\delta)^{-20}$.
Since $\|f\|_\infty\le K\sqrt{Q_G}$ and $\|g\|_\infty\le1$,
posterior averaging and the displayed energy bound make the absolute
transfer error $O_{\phi,K_Y}(n^{-10})$. A confidence allocation between
this no-exceedance event and the clipped-label statistical event proves
(AD14), with exact leading comparator coefficient one and a common
event over the eligible widths and scales.

\paragraph{A hardtanh learner competing with shallow ReLU teachers.}
The hardtanh instance also gives an exact, same-depth comparison with
coefficient-bounded ReLU teachers. Write
\begin{gather*}
 g(x)=b_{\rm out}+\sum_{i=1}^C c_i\operatorname{ReLU}(a_i\cdot x+b_i),\\
 |b_{\rm out}|,|c_i|,|(a_i)_j|,|b_i|\le B,
 \qquad B\ge1,\qquad M=(d+1)B.
\end{gather*}
On the input cube, $|a_i\cdot x+b_i|\le M$, and the exact identity
\[
 \operatorname{ReLU}(t)
 =\frac M2\{\operatorname{hardtanh}(2t/M-1)+1\},\qquad |t|\le M,
\]
converts $g$ to a hardtanh teacher of width at most $C+1$.
The extra neuron is the constant $\operatorname{hardtanh}(1)=1$.
Its readout variation obeys
\[
 V'\le M\sum_i|c_i|+|b_{\rm out}|
 \le(d+1)CB^2+B.
\]
Each transformed incoming row has Euclidean norm at most three,
including the constant row, so $J'\le9a_0$.
The comparison function itself is unchanged. Therefore (AD11) applies
when $\|g\|_\infty\le1$, giving fixed-$d,B$ sample complexity
$\widetilde O(C^2\varepsilon^{-2})$ under any input law; (AD13) gives
the coarser no-output-cap statement. Marginally sub-Gaussian labels are
covered by (AD14) when the comparison output is bounded.
The actual learner here has hardtanh activation. This is not a claim
that the raw ReLU learner has the bounded-ridge capacity bound (AD3).


\section{Removing the comparison-output restriction}
\label{sec:risk-weighted-multiplier}
A variance-sensitive multiplier estimate removes the output cap from the
preceding arbitrary-input-law oracle without worsening its rate.
The fixed comparison is selected using population risk, before sampling;
it need not have a uniformly small output.

Keep the Gaussian model, activation, resource $A_n(g)$ in (AD1), and
input cube of the preceding section, but now allow any finite
comparison network $g$, with no output restriction. Write
$V=\sum_i|c_i|$ and $F=\|g\|_\infty\le K V$.
For the proof it suffices to consider $R(g)\le R(0)+1\le2$ under
bounded labels: a comparison with larger risk can be replaced in the
proof by the zero network, at the same likelihood scale.

\begin{lemma}[Risk-weighted ridge multipliers]
\label{lem:risk-weighted-ridge}
Let $u$ be a fixed measurable function with $\|u\|_\infty\le F_u$,
let $|Y|\le B$, and put
$m(X,Y)=u(X)-Y$ and $s^2=\mathbb E m^2$.
For $t=\log(2/\delta)$, with probability at least $1-\delta$,
\[
 \sup_\theta\bigl|(P-P_n)[m\phi(\theta\cdot\widetilde X)]\bigr|
 \le C_\phi\left\{
       s\sqrt{(D+t)/n}+(F_u+B)t/n\right\}=:\eta_n.       \tag{WM1}
\]
The inequality holds simultaneously over every incoming weight vector,
without a norm bound or any independence between $m$ and $X$.
\end{lemma}
\begin{proof}
Condition on the observed pairs and write $m_i=m(X_i,Y_i)$.
The weighted Rademacher process has increment metric
\[
 d_m(h,h')^2=n^{-2}\sum_i m_i^2(h(X_i)-h'(X_i))^2.
\]
If $\sum_i m_i^2>0$, normalize these weights to a probability measure
$P_{n,m}=\sum_i m_i^2\delta_{X_i}/\sum_i m_i^2$.
The VC-subgraph covering bound for the bounded ridge class is valid for
this probability measure too. Dudley's integral therefore gives
\[
 \mathbb E_\epsilon\sup_\theta
 \left|n^{-1}\sum_i\epsilon_i m_i\phi(\theta\cdot\widetilde X_i)\right|
 \le C_\phi\sqrt D\,\frac{\sqrt{\sum_i m_i^2}}n.
\]
For zero total weight the process is identically zero. Taking expectation,
using Jensen, and symmetrizing shows that the expected centered supremum
is at most $C_\phi s\sqrt{D/n}$.

For the class with both signs, each summand has envelope
$U=K(F_u+B)$ and variance at most $K^2s^2$.
Bousquet's inequality~\cite[Theorem~2.3]{bousquet} bounds its centered supremum by its expectation
plus
\[
 C\sqrt{(K^2s^2+U\mathbb E Z)t/n}+CUt/n.
\]
The mixed term $\sqrt{U(\mathbb E Z)t/n}$ is at most
$C\mathbb E Z+CUt/n$, by Young's inequality.
This proves (WM1). The weighted covering argument is conditional on the
observed pairs and uses only $m_i^2$, so dependence on $X$ causes no issue.
\end{proof}

\begin{theorem}[Sharp uncapped-comparison oracle]
\label{thm:uncapped-all-d}
Assume $|Y|\le1$, and fix any comparison network $g$.
Let $b_n$ be (AD4), with a sufficiently large
activation-dependent constant. In the nontrivial range
$A_n(g)b_n\le c_\phi$, with probability at least $1-\delta$,
simultaneously for all $0<\nu\le1$, all
$n^{-2}\le\lambda\le1$ with $\lambda\ge C_\phi b_n$, and all
$N\lambda\ge C_\phi$,
\[
 R(\mathbb E_{\Pi_\lambda}f)\le\Pi_\lambda R(f)
 \le R(g)+C_\phi\{\lambda A_n(g)+b_n/\lambda\}
       +C_\phi n^{-10}.                                \tag{WM2}
\]
In particular, $\lambda=C_\phi\sqrt{b_n/A_n(g)}$ yields
\[
 R(\mathbb E_{\Pi_\lambda}f)\le\Pi_\lambda R(f)
 \le R(g)+\widetilde O_\phi\left(\sqrt{(d+1)A_n(g)/n}\right),         \tag{WM3}
\]
with the same learner-width condition as (AD11). There is no output
cap on either the comparison network or the learner.
\end{theorem}
\begin{proof}
It suffices to prove the bound when $R(g)\le R(0)+1$. Otherwise
apply that case to zero at the same $\lambda$: $A_n(0)\le A_n(g)$
and $R(0)\le R(g)$, so its conclusion is stronger than (WM2).
This is a proof choice and does not require the learner to know either risk.

Use the variation-weighted comparison pack from (AD8), with
$\sigma=n^{-20}/(V+1)$, and write $u=g_\sigma$.
Its uniform error is $O_\phi(n^{-20})$. By the triangle inequality in
population $L^2$, its residual risk satisfies
\[
 \sqrt{R(u)}\le\sqrt{R(g)}+C_\phi n^{-20}.
\]
Thus $R(u)$ is bounded by an activation-dependent constant,
without using the output envelope $F$. Moreover $u$ is a signed
Gaussian mixture of ridge features with total variation exactly $V$
(or at most $V$ after combining coincident features).

Apply (WM1) to $m=u-Y$. Its event is uniform over every ridge, so for
any learner function $f\in\mathcal F_q$,
\[
 |(P-P_n)[m(f-u)]|\le(\sqrt q+V)\eta_n.                 \tag{WM4}
\]
For the $u$ term this follows by integrating the common ridge bound
against its signed mixing measure; boundedness justifies Fubini.
For $f$ it follows by its finite signed-hull representation.
Since $\|u\|_\infty\le F+C_\phi n^{-20}$,
\[
 \eta_n^2\le C_\phi\{b_n+(F+1)^2b_n^2\}
 \le C_\phi b_n.                                      \tag{WM5}
\]
The last step uses $(F+1)^2\le C_\phi A_n$ and $A_nb_n\le c_\phi$.
Young's inequality in (WM4), with a chosen small constant, gives
\[
 2|(P-P_n)[m(f-u)]|
 \le c\lambda q+C\lambda V^2+C_\phi b_n/\lambda.        \tag{WM6}
\]

The smooth-loss norm comparison (AD6) does not need $u$ to be bounded
by two. With its actual envelope it gives
\[
 (P_X-P_n)(f-u)^2
 \le\tfrac16P_n(f-u)^2
       +C_\phi b_n\{q+(F+1)^2\}.                      \tag{WM7}
\]
Its $q$ term is absorbed into a small multiple of $\lambda q$ by
$\lambda\ge C_\phi b_n$; its constant term is at most
$C_\phi\lambda A_n$. Combine (WM6)--(WM7) with the relative-risk
identity used in (AD5), integrate, and use the exact midpoint
comparison (AD9). That comparison and its pack entropy do not require
an output cap. Its empirical-distance and readout-energy penalties
absorb the chosen small terms, proving (WM2).

For completeness, the same finite energy cutoff is valid here.
The condition $A_nb_n\le c$ implies $F^2\le C_\phi A_n\lesssim n$.
The discarded distance and relative-risk integrands are bounded by
$C_\phi(1+Q_G+F^2)$, so the original Gaussian readout tail above
$n^{22}$ remains negligible uniformly over the claimed scales.
Finally $R(u)\le R(g)+C_\phi n^{-20}$ by the residual $L^2$ inequality,
and Jensen gives the raw-mean claim. All sample events use only the
fixed $u$ and shared feature classes, preserving simultaneity in $N,\nu$.
\end{proof}

\paragraph{Marginally sub-Gaussian labels and the class infimum.}
Suppose instead $\mathbb E\exp(Y^2/K_Y^2)\le2$.
Fix any uncapped comparison network. As above, comparisons with
$R(g)>R(0)+1$ are replaced by zero only in the proof.
The same conclusion holds with
\[
 b_n^*=C_{\phi,K_Y}(D+1)\log^5(e n/\delta)/n
\]
in place of $b_n$, under the same nontrivial-range and width conditions.
To verify it, use the proof-only label cutoff from (AD14),
$B_n=\max\{1,K_Y\sqrt{40\log(e n/\delta)}\}$.
The clipped-label residual has bounded second moment because
\[
 \|g_\sigma-Y^{B_n}\|_{L^2}
 \le\sqrt{R(g)}+\|Y-Y^{B_n}\|_{L^2}+C_\phi n^{-20}
 \le C_{\phi,K_Y}.
\]
Apply (WM1) with envelope $F+C_\phi+B_n$. Its squared bound is
$C_{\phi,K_Y}\{b_n^*+(F+1)^2(b_n^*)^2\}$, absorbed by
$C_{\phi,K_Y}b_n^*$ when $A_nb_n^*$ is small.
The other proof steps are unchanged. The actual posterior agrees with
the proof's clipped-label posterior on the no-exceedance event, and
\[
 [R_Y(f)-R_Y(g)]-[R_{Y^{B_n}}(f)-R_{Y^{B_n}}(g)]
 =-2\mathbb E[(Y-Y^{B_n})(f(X)-g(X))].
\]
Here $F\lesssim\sqrt{A_n}\lesssim\sqrt n$, and the readout energy
moment from (AD14) bounds the posterior average of $\|f\|_\infty$.
Thus the marginal tail still makes this transfer $O_{\phi,K_Y}(n^{-10})$.
No conditional-mean assumption and no actual label clipping are used.

For any declared coefficient-bounded shallow teacher class containing
zero, select a deterministic population near-minimizer; it satisfies
$R(g)\le R(0)+1$. Equations (WM2)--(WM3) therefore hold with its
risk replaced by the class infimum and a deterministic upper bound on
$A_n$, with exact leading coefficient one. For fixed $d,B$, the error
is $\widetilde O(C/\sqrt n)$ against \emph{all} width-$C$ teachers
with coefficient cap $B$, not just those with bounded output.
In particular, the single schedule $\lambda=n^{-3/4}$ competes with
all width-$n^{1/4}$ such teachers at excess
$\widetilde O(n^{-1/4})$, simultaneously for every $N\ge c_\phi n$
and $0<\nu\le1$, once $n$ is sufficiently large. The schedule does not
use the unknown comparator coefficients or risk.
The same removal of the output cap applies to the exact ReLU-to-hardtanh
compiler in the preceding section.

\begin{corollary}[Growing width and coefficient budgets]
\label{cor:all-d-growing-resources}
Fix $d$, the activation, a marginal sub-Gaussian label envelope, and
nonnegative exponents $a,b$ with $a+b<1/2$. For a class of shallow
teachers of width at most $n^a$ and coefficient cap $n^b$, with no
output cap, a deterministic scale
$\lambda_n=\widetilde\Theta(n^{-1/2-a-b})$ gives
\[
 R(\mathbb E_{\Pi_{\lambda_n}}f)\le\Pi_{\lambda_n}R(f)
 \le\inf_g R(g)+\widetilde O(n^{a+b-1/2})
\]
with probability at least $1-\delta$, simultaneously over all
$0<\nu\le1$ and $N\ge c_\phi n$, for sufficiently large $n$.
The scale uses the declared growth budgets, not the unknown teacher.
For a hardtanh learner against coefficient-capped ReLU teachers, the
corresponding condition is $a+2b<1/2$, with remainder
$\widetilde O(n^{a+2b-1/2})$.
\end{corollary}
\begin{proof}
Equation (AD2) gives $A_n=\widetilde O(n^{2(a+b)})$ at fixed $d$.
Balance the two terms of (WM2). Its nontrivial-range and width conditions
hold because $a+b<1/2$. The hardtanh compiler instead has
$V'=O_d(n^{a+2b})$ and $J'=O_d(1)$, giving the stated substitution.
\end{proof}

There is also one schedule that needs neither growth exponent. At fixed
$d$, put $\ell_n=\log(e n/\delta)$ and
$\lambda_n=\ell_n^7/n$. For sufficiently large $n$, (WM2) gives
\begin{equation}
 R(\mathbb E_{\Pi_{\lambda_n}}f)\le\Pi_{\lambda_n}R(f)
 \le\inf_{g\in\mathcal G_n}R(g)
 +C_{\phi,d,K_Y}\left\{\frac{a_n\ell_n^7}{n}+\ell_n^{-2}\right\}
 +C_{\phi,d,K_Y}n^{-10},                              \tag{WM8}
\end{equation}
for any declared class $\mathcal G_n$ containing zero with
$\sup_{g\in\mathcal G_n}A_n(g)\le a_n$ and
$a_n b_n^*\le c_\phi$. The event is common to every eligible learner
width $N\ge c_\phi n$ and scale $0<\nu\le1$; the algorithm does not
use $a_n$. Indeed $b_n^*/\lambda_n=O_{\phi,d,K_Y}(\ell_n^{-2})$,
while the first term is $\lambda_n a_n$. Thus this same schedule gives
a vanishing remainder for each fixed pair $a,b$ with $a+b<1/2$ above.
This means one learning rule has the stated guarantee for each declared
class, not that all data-selected comparison networks share a single
unproved event.

\section{Finite-variance labels suffice for a confidence oracle}
\label{sec:finite-variance-midpoint}
The same raw Gaussian learner has a high-probability oracle assuming only
$\mathbb E Y^2<\infty$. This extension pays a polynomial confidence factor.
It makes no expected-risk claim and does not clip observed labels.

\begin{lemma}[A second-moment weighted empirical-process bound]
\label{lem:weighted-second-moment}
Let $m(X,Y)$ be any fixed square-integrable multiplier and let
$\mathcal H=\{x\mapsto\phi(\theta\cdot\widetilde x):\theta\in\mathbb R^D\}$
be the bounded monotone ridge class. Then
\[
 \mathbb E\left[\sup_{h\in\mathcal H}
             |(P-P_n)(mh)|^2\right]
 \le\frac{C_\phi D\,\mathbb E m^2}{n}.                 \tag{FV1}
\]
No tail assumption or independence of $m$ and $X$ is needed.
\end{lemma}
\begin{proof}
Let $S'$ be an independent ghost sample. Jensen's inequality and
coordinatewise swapping of the two samples using independent Rademacher
signs give
\begin{align*}
 \mathbb E\sup_h|(P-P_n)(mh)|^2
 &\le\mathbb E_{S,S'}\sup_h|(P_n'-P_n)(mh)|^2\\
 &\le4\mathbb E_{S,\epsilon}\sup_h
           \left|n^{-1}\sum_i\epsilon_i m_i h(X_i)\right|^2.
\end{align*}
For the last inequality, apply the triangle inequality after inserting
the signs, then $(a+b)^2\le2a^2+2b^2$.
All quantities are integrable, since the feature envelope is bounded
and $m$ has a second moment.

Condition on $S$. The weighted VC covering argument in (WM1) bounds
the conditional mean of the Rademacher supremum by
$C_\phi\sqrt{D\sum_i m_i^2}/n$.
Changing one sign changes this supremum by at most $2K|m_i|/n$.
The bounded-difference variance inequality therefore bounds its
conditional variance by $C K^2\sum_i m_i^2/n^2$.
Adding variance and squared mean, then averaging over $S$, proves (FV1).
As usual, continuity of the activation gives a countable determining
parameter subclass, so the empirical suprema are measurable.
\end{proof}

\begin{theorem}[An unclipped finite-variance oracle]
\label{thm:finite-variance-midpoint}
Use the independent Gaussian model and resource $A_n(g)$ of (AD1), with
any input law on the cube and any fixed finite comparison network $g$.
Suppose $\mathbb E Y^2\le M_2<\infty$. Let $0<\delta<1/2$ and put
\[
 \ell_n=\log\frac{e n(1+M_2)}\delta,
 \qquad
 b_n^{\rm fv}=\frac{C_\phi D(M_2+1)\ell_n^4}{n\delta}.  \tag{FV2}
\]
In the nontrivial range $A_n(g)b_n^{\rm fv}\le c_\phi$, with probability
at least $1-\delta$, simultaneously for all $0<\nu\le1$, all
$n^{-2}\le\lambda\le1$ with $\lambda\ge C_\phi b_n^{\rm fv}$,
and all $N\lambda\ge C_\phi$,
\[
 R(\mathbb E_{\Pi_\lambda}f)\le\Pi_\lambda R(f)
 \le R(g)+C_\phi\{\lambda A_n(g)+b_n^{\rm fv}/\lambda\}
       +C_\phi(M_2+1)n^{-10}.                         \tag{FV3}
\]
There is no comparison-output cap. The prior, likelihood, observed
labels, and raw posterior predictions are unchanged.
\end{theorem}
\begin{proof}
As in (WM2), it suffices to prove the result for
$R(g)\le R(0)+1\le M_2+1$. Otherwise use zero at the same likelihood
scale; $A_n(0)\le A_n(g)$ and $R(0)\le R(g)$.
Use the same variation-weighted Gaussian comparison pack and mean
$u=g_\sigma$. Its error is $C_\phi n^{-20}$ uniformly, so
\[
 \mathbb E(u-Y)^2
 \le(\sqrt{R(g)}+C_\phi n^{-20})^2
 \le C_\phi(M_2+1).
\]
Apply (FV1) to $m=u-Y$ and use Markov's inequality. With failure
probability at most $\delta/3$, the common ridge supremum $Z$ satisfies
\[
 Z^2\le C_\phi D(M_2+1)/(n\delta)\le C_\phi b_n^{\rm fv}.
\]
Both $f$ and $u$ belong to signed ridge hulls with variations at most
$\sqrt{Q_G}$ and $V$. Thus
\[
 |(P-P_n)[(u-Y)(f-u)]|\le(\sqrt{Q_G}+V)Z,
\]
and Young's inequality gives the exact analogue of (WM6), with
$b_n^{\rm fv}$ in place of $b_n$. Fubini is justified by
$\mathbb E|u-Y|<\infty$ and the bounded ridge features.

Separately, Markov's inequality gives
\[
 P_nY^2\le3M_2/\delta
\]
with failure probability at most $\delta/3$ (when $M_2=0$, it holds
almost surely). On that event the native empirical risk is at most
$3M_2/\delta+K^2$. Write $L=1+M_2/\delta$.
Jensen's inequality gives
$Z_\lambda\ge\exp[-C_\phi\omega L/(2\lambda)]$.
The native readout energy still has $\omega Q_G\sim\chi^2_N$.
Its Gaussian tail therefore makes all posterior risk, empirical-loss,
and quadratic-distance integrals on
\[
 Q_G>Q_*:=C_\phi n^{22}(1+M_2/\delta)^2
\]
negligible, uniformly for $\lambda\ge n^{-2}$.
Indeed this cutoff dominates $C_\phi L/\lambda$; beyond it the
posterior tail is at most $e^{-c\omega Q_G}$.
The relevant integrands are bounded by a constant times
$Q_G+F^2+M_2+P_nY^2+1$, where $F\le K V$ and
$F^2\le C_\phi A_n\lesssim n$ in the nontrivial range.
Their omitted integrals are bounded by $C_\phi(M_2+1)n^{-10}$.

Apply the squared-norm SST event to the classes $\mathcal F_q-u$ at
dyadic radii through $2Q_*$, assigning total failure $\delta/3$.
Their number is $O_\phi(\ell_n)$, so the logarithmic budget in (FV2)
suffices. This event uses only the inputs; its conclusion is
\[
 (P_X-P_n)(f-u)^2\le\tfrac16P_n(f-u)^2
 +\frac{C_\phi D\ell_n^4}{n}\{Q_G+(F+1)^2\}.
\]
Since $\lambda\ge C_\phi b_n^{\rm fv}$, absorb its energy term into
a small multiple of $\lambda Q_G$ and its constant term into
$C_\phi\lambda A_n$.

The three events have joint probability at least $1-\delta$.
On their intersection, combine the relative-risk identity, the weighted
multiplier bound, and the unchanged exact midpoint objective (AD9).
The same curvature and readout-energy absorption as in (WM2) proves
(FV3). The population smoothing error is bounded by
$C_\phi(\sqrt{M_2+1}+1)n^{-20}$ via the residual $L^2$ triangle
inequality. Jensen proves the statement for the ordinary raw mean.
Each event is fixed independently of the eligible $N,\nu,\lambda$,
which proves the asserted simultaneity.
\end{proof}

\paragraph{Rates.}
Balancing the two terms with
$\lambda=C_\phi\sqrt{b_n^{\rm fv}/A_n(g)}$ gives
\[
 R(\mathbb E_{\Pi_\lambda}f)\le\Pi_\lambda R(f)
 \le R(g)+\widetilde O_\phi\left(
       \sqrt{\frac{(d+1)(M_2+1)A_n(g)}{n\delta}}\right). \tag{FV4}
\]
At fixed $d,B,M_2,\delta$, this is $\widetilde O(C/\sqrt n)$ against
all width-$C$ coefficient-bounded shallow teachers. As before, select a
deterministic population near-minimizer to obtain a class-infimum
oracle with exact leading coefficient one.
In particular, for every fixed coefficient cap $B$ and every finite
second-moment label law, the single schedule $\lambda=n^{-3/4}$ has
confidence-level-$1-\delta$ excess $\widetilde O(n^{-1/4})$ against all
width-$n^{1/4}$ comparison networks, simultaneously at every $N\ge c_\phi n$
and $0<\nu\le1$, once $n$ is sufficiently large.
The constants and the starting sample size may depend on $M_2,\delta$;
this fixed schedule does not require knowing $M_2$.
This is a confidence statement with a polynomial dependence on
$1/\delta$, not an expected-risk guarantee.

At fixed $M_2$ and confidence, the growing-resource corollary and
universal schedule after (WM8) remain valid with the same powers of
$n$. Indeed $b_n^{\rm fv}$ is still $\widetilde O(1/n)$; inserting
$\lambda_n=\log^7(e n/\delta)/n$ in (FV3) gives a remainder tending
to zero whenever $A_n=\widetilde O(n^{2(a+b)})$ with $a+b<1/2$.
The constants and starting sample size may depend on $M_2,\delta$,
but this schedule does not require the second-moment bound as input.
The ReLU-to-hardtanh compiler instead has the condition $a+2b<1/2$.
These remain high-probability statements, not assertions about expected
risk under arbitrary heavy-tailed labels.

\section{The same oracle at every learner width above the teacher}
\label{sec:all-width-finite-variance}
A direct cloned comparison and a variable-tolerance squared-norm estimate
remove the condition $N\lambda\gtrsim1$. The following result applies to
every learner width at least the comparison width. It keeps the same
independent Gaussian prior, raw likelihood and raw prediction.
The activation is bounded, nonconstant, nondecreasing and globally
Lipschitz, as in (AD1); bounded oscillatory activations are not covered
by the ridge-class VC argument.

Let the comparison network have width $C\ge1$ and parameters
$g(x)=\sum_{i=1}^C c_i\phi(\vartheta_i\cdot\widetilde x)$.
There is no output cap. Put
\[
 V=\sum_i|c_i|,\qquad
 E_g=\frac{a_0}{C}\sum_i\|\vartheta_i\|_2^2
                       +2C\sum_i c_i^2,\qquad
 L=\frac{16(M_2+1)}\delta,
\]
\[
 \mathcal A_n(g)=1+E_g+(D+1)
       \log\{e n(V+1)(\sqrt L+1)\},\qquad
 b_n=\frac{C_\phi D[\log(e nL)]^4}{n}.                 \tag{AW1}
\]
Here $\mathbb E Y^2\le M_2<\infty$, $0<\delta<1/2$, and the other
model notation is as in the preceding sections. Since
$V^2\le C\sum_i c_i^2$, both $V^2$ and $\|g\|_\infty^2$ are
controlled by $C_\phi E_g$. If all coefficients are bounded by $B\ge1$,
then $E_g\le a_0D B^2+2C^2B^2$.

\begin{lemma}[A cloned comparison at every $N\ge C$]
\label{lem:all-width-clone}
On a sample satisfying $\widehat R(g)\le L$, for all $N\ge C$ and
$0<\nu\le1$ the actual posterior satisfies
\[
 e_\lambda+\tfrac12\lambda\Pi_\lambda Q_G
 \le C_\phi\lambda\mathcal A_n(g)+C_\phi n^{-30},
 \qquad
 e_\lambda=\Pi_\lambda\widehat R(f)-\widehat R(g).       \tag{AW2}
\]
\end{lemma}
\begin{proof}
Let $k=\lfloor N/C\rfloor\ge N/(2C)$. Replace each teacher neuron by
$k$ clones with the same incoming row and readout coefficient $c_i/k$;
set all remaining parameter centers to zero. This computes $g$ exactly,
even when $\phi(0)\ne0$, since the remaining readout centers vanish.
In pack coordinates $G_j=Nv_j$, the center's normalized energy obeys
\[
 Q_{\rm all}(w_*)=\frac{a_0}N\sum_j\|\theta_j^*\|^2
                         +\frac1N\sum_j(G_j^*)^2
 \le E_g.
\]
Indeed $k/N\le1/C$ and $N/k\le2C$.
Around this center take independent Gaussian perturbations with each
native coordinate variance multiplied by $\epsilon^2$, where
\[
 \epsilon=\frac{n^{-30}}{C_\phi(V+1)(\sqrt L+1)}.
\]
For every input on the cube the incoming perturbation's preactivation
standard deviation is at most $\sqrt2\epsilon$. Minkowski bounds the
mean-readout contribution to the output error by
$C_\phi V\epsilon$. Conditional on incoming perturbations, the centered
readout-noise variance is at most $K^2\epsilon^2\nu/N$.
The uniform root-mean-square error is consequently at most
$r=C_\phi\epsilon(V+1)$. Thus the comparison law $\rho_g$ has
\[
 \rho_g\widehat R(f)
 \le\widehat R(g)+2r\sqrt{\widehat R(g)}+r^2
 \le\widehat R(g)+C_\phi n^{-30}.
\]
There are $(D+1)N$ Gaussian coordinates. The exact Gaussian entropy gives
\[
 \operatorname{KL}(\rho_g\Vert\mu)
 \le\frac\omega2E_g+(D+1)N\log(1/\epsilon)
 \le C_\phi\omega\mathcal A_n(g),
\]
using $N/\omega=\nu\le1$.
Gibbs optimality bounds its free energy by
$\widehat R(g)+C_\phi\lambda\mathcal A_n+C_\phi n^{-30}$.
Finally the readout-only entropy lower bound
$\operatorname{KL}(\Pi_\lambda\Vert\mu)
\ge\omega\Pi_\lambda Q_G/4-(N/2)\log2$
yields (AW2), absorbing the extra $\lambda\nu\log2$ in the resource.
\end{proof}

\begin{theorem}[All-width finite-variance raw oracle]
\label{thm:all-width-raw}
For every fixed finite comparison network $g$, assume
$\mathbb E Y^2\le M_2$ and
$\mathcal A_n(g)L b_n\le c_\phi$.
With probability at least $1-\delta$, simultaneously for every
$N\ge C$, $0<\nu\le1$, and every $n^{-2}\le\lambda\le1$ satisfying
$\lambda\ge C_\phi b_n$,
\[
 R(\mathbb E_{\Pi_\lambda}f)\le\Pi_\lambda R(f)
 \le R(g)+C_\phi\left\{\lambda\mathcal A_n(g)
                             +\frac{L b_n}{\lambda}\right\}
       +C_\phi(M_2+1)n^{-10}.                          \tag{AW3}
\]
There is no condition on $N\lambda$. In particular, choosing
$\lambda=C_\phi\sqrt{L b_n/\mathcal A_n(g)}$ gives
\[
 R(\mathbb E_{\Pi_\lambda}f)\le\Pi_\lambda R(f)
 \le R(g)+\widetilde O_\phi\left(
       \sqrt{\frac{(d+1)(M_2+1)\mathcal A_n(g)}{n\delta}}\right).
                                                               \tag{AW4}
\]
No actual output or label is clipped, and the input distribution on the
cube is arbitrary. This is a confidence guarantee, not an expected-risk
claim.
\end{theorem}
\begin{proof}
It suffices to consider $R(g)\le R(0)+1\le M_2+1$.
If not, use the zero network at the same $\lambda$ in the proof;
its resource is no larger and its risk is smaller.
The finite-variance weighted feature estimate (FV1), now applied directly
to $m=g-Y$, gives with failure at most $\delta/4$
\[
 Z^2:=\sup_\theta|(P-P_n)[(g-Y)\phi(\theta\cdot\widetilde X)]|^2
 \le C_\phi D(M_2+1)/(n\delta)\le C_\phi L b_n.         \tag{AW5}
\]
Since $g$ and $f$ are signed ridge sums,
\[
 |(P-P_n)[(g-Y)(f-g)]|\le(\sqrt{Q_G}+V)Z.
\]
Young's inequality therefore bounds twice this multiplier by
$c\lambda Q_G+C_\phi\lambda\mathcal A_n+C_\phi Lb_n/\lambda$, where
$c>0$ is as small as needed.

Markov's inequality applied to
$P_n[(g-Y)^2+Y^2]$ gives another event, of failure at most $\delta/4$,
on which both $\widehat R(g)\le L$ and $P_nY^2\le L$.
The same evidence and Gaussian tail calculation as in (FV3) controls
posterior integrals above $Q_*=C_\phi n^{22}L^2$, uniformly for
$\lambda\ge n^{-2}$. All relative-risk and empirical-loss integrands
are bounded by $C_\phi(Q_G+\|g\|_\infty^2+M_2+L+1)$.
The nontrivial-range assumption bounds $\|g\|_\infty^2$ by
$C_\phi\mathcal A_n\lesssim n$, so the omitted integrals are negligible.

On a third event, of failure at most $\delta/2$, the SST squared-norm
inequality for the dyadic signed-hull classes through $Q_*$ gives,
for every $0<\eta\le1$,
\[
 (P_X-P_n)(f-g)^2
 \le\eta P_n(f-g)^2+\frac{C_\phi b_n}{\eta}(Q_G+V^2+1).
                                                               \tag{AW6}
\]
This is the same nonnegative smooth-loss inequality as (AD6), with
Young's inequality performed at tolerance $\eta$. It holds for every
such tolerance on the same event; no additional confidence allocation
is needed. The hull entropy and its squared envelope are controlled
by $Q_G+V^2+1$, and the number of energy shells is $O_\phi(\log(e nL))$.

Choose $\eta=K_\phi b_n/\lambda$, with $K_\phi$ large enough that
the $Q_G$ coefficient in (AW6) is at most $\lambda/16$.
The stipulated likelihood floor ensures $\eta\le1/8$.
The remaining part of that term is at most
$C_\phi\lambda\mathcal A_n$.
Combining (AW5)--(AW6) with the relative-risk identity, integrating,
and writing $e=e_\lambda$, gives
\begin{align*}
 \Pi_\lambda R(f)-R(g)
 &\le e+\eta\Pi_\lambda P_n(f-g)^2
          +\tfrac18\lambda\Pi_\lambda Q_G
          +C_\phi\lambda\mathcal A_n+C_\phi Lb_n/\lambda
          +C_\phi(M_2+1)n^{-10}\\
 &\le(1+2\eta)e+4\eta\widehat R(g)
          +\tfrac18\lambda\Pi_\lambda Q_G
          +C_\phi\lambda\mathcal A_n+C_\phi Lb_n/\lambda
          +C_\phi(M_2+1)n^{-10}.
\end{align*}
The second line uses
$P_n(f-g)^2\le2\widehat R(f)+2\widehat R(g)$.
Keep the entire expression $(1+2\eta)e$: it is valid even when $e<0$.
By (AW2), its negative readout-energy contribution absorbs the displayed
$\lambda\Pi Q_G/8$. Also
$4\eta\widehat R(g)\le4K_\phi Lb_n/\lambda$.
Since $\eta\le1/8$, all remaining comparison terms are bounded by
$C_\phi\lambda\mathcal A_n$. This proves (AW3); Jensen gives the raw
mean. The three events are common to all claimed widths and scales.
Balance the two terms for (AW4); the nontrivial-range assumption ensures
that this scale meets its stated likelihood bounds.
\end{proof}

\paragraph{One fixed schedule and a growing comparison class.}
For fixed $d,B,M_2,\delta$, the resource is
$\widetilde O(C^2)$ for every width-$C$ teacher with entrywise coefficient
cap $B$. Taking $C=\lfloor n^{1/4}\rfloor$ and $\lambda=n^{-3/4}$ gives
\[
 R(\mathbb E_{\Pi_\lambda}f)\le\Pi_\lambda R(f)
 \le\inf_{g\text{ of width }\le n^{1/4},\,|\text{coefficients}|\le B}
                  R(g)+\widetilde O(n^{-1/4}),           \tag{AW7}
\]
with probability at least $1-\delta$, simultaneously for every
$N\ge\lfloor n^{1/4}\rfloor$ and $0<\nu\le1$, for sufficiently large
$n$. The schedule uses neither $M_2$ nor the unknown comparator risk;
constants and the starting sample size can depend on them.
The infimum follows by a deterministic population near-minimizer and
a common deterministic resource upper bound.
For widths $C\le n^a$ and coefficient caps $B\le n^b$, the same
substitution gives a vanishing power rate whenever $a+b<1/2$,
using $\lambda=n^{-1/2-a-b}$ and excess
$\widetilde O(n^{a+b-1/2})$ at every $N\ge C$.
The earlier ReLU-to-hardtanh compiler retains its separately stated
wide-learner guarantee. Its variation resource is not substituted for
the cloned quadratic resource $E_g$ in this all-width theorem.

\section{Empirical midpoint comparison in every input dimension}
Let $x\in[-1,1]^d$, $d\ge1$, $\bar x=(x,1)$, and
\[
 z_w(x)=\sum_{j=1}^N v_j(w_j\cdot\bar x)_+,
 \qquad w_j\sim N(0,(\nu/d)I_{d+1}),\qquad
 v_j\sim N(0,\nu/N^2),
\]
independently, where $0<\nu\le1$. Define
\[
 Q=\frac dN\sum_j\|w_j\|_2^2+N\sum_jv_j^2,
 \qquad P=\sqrt d\sum_j|v_j|\|w_j\|_2\le Q/2,
 \qquad\omega=N/\nu.
\]
The actual posterior is proportional to
$\exp[-\omega\widehat R(z)/(2\lambda)]d\mu$.
Thus its Gaussian observation variance is $n\lambda\nu/N$.
Assume the observed labels are bounded by one.

\begin{lemma}[Minimum-mass training-vector compression]
For any finite sample of $n$ inputs and any two finite shallow ReLU
functions $f,g$ of normalized path masses $P_f,P_g$, their midpoint
training vector has an implementation using at most $n$ neurons,
with path mass at most $(P_f+P_g)/2$. The choice is Borel in the
original parameters and sample inputs.
\end{lemma}
\begin{proof}
Use normalized features
$\phi_\theta(x)=d^{-1/2}(\theta\cdot\bar x)_+$,
$\theta\in\mathbb S^d$, with the readout coefficient signs absorbed
into the signed features. Let $A$ contain all the signed training
feature columns of $f$ and $g$. The linear program
\[
 \min\{\boldsymbol1^\top\alpha:A\alpha=(f(S)+g(S))/2,
                                  \ \alpha\ge0\}
\]
is feasible with cost at most $(P_f+P_g)/2$ and has compact sublevel
sets. A minimum-support optimum has independent supported columns:
a dependence would permit motion until a coefficient vanished without
increasing cost. It therefore uses at most $\operatorname{rank}A\le n$
atoms. Enumerate independent column subsets, solve by their Gram
inverses, retain the nonnegative feasible solutions, and choose the
least cost with lexicographic tie-breaking. This is a Borel rule.
\end{proof}

\begin{lemma}[A dimension-explicit empirical midpoint inequality]
\label{lem:general-shallow-midpoint}
Let $n\ge2$, $N\ge n$, $n^{-2}\le\lambda\le1$.
Fix any finite comparator $g$ with $P_g\le n^4$.
Then the actual posterior satisfies
\begin{align}
 &\Pi_\lambda\widehat R(z)-\widehat R(g)
 +\tfrac12\Pi_\lambda P_n(z-g)^2+\lambda\Pi_\lambda P\\
 &\hspace{2cm}\le2\lambda P_g+
 C\lambda\nu(d+2)\log((n+1)(d+1))+Cn^{-10}. \tag{GM1}
\end{align}
The comparator need not fit in width $N$: only its compressed midpoint
training vector is used by the proof. The learner itself is unchanged.
\end{lemma}
\begin{proof}
For each posterior function $f$, compress its midpoint training vector
with $g$ as above and pad to width $N$. Balance a normalized signed
atom of amplitude $c$ by
\[
 \bar w_j=\sqrt{N|c|/d}\,\theta_j,
 \qquad \bar v_j=\operatorname{sign}(c)\sqrt{|c|/N}.
\]
The center has energy at most $P_f+P_g$ and agrees exactly with
$(f+g)/2$ at all training inputs.

Put a Gaussian law around this center with native covariance
multiplied by $\epsilon^2$, where
$\epsilon=((n+1)(d+1))^{-30}$, and mix these laws using the actual
posterior distribution of $f$. Its KL is at most
\[
 \frac\omega2(\Pi P_f+P_g)+(d+2)N\log(1/\epsilon).
\]
This follows from the Gaussian relative-entropy formula and
convexity of relative entropy under mixing. On the other hand,
\[
 \log\mathbb E_\mu e^{3\omega Q/8}=(d+2)N\log2,
\]
so the entropy inequality and $Q\ge2P$ give
\[
 \operatorname{KL}(\Pi\Vert\mu)
 \ge\tfrac34\omega\Pi P-(d+2)N\log2.
\]
The difference between the mixture's KL and the actual posterior KL
is therefore bounded by
\[
 -\tfrac14\omega\Pi P+\tfrac12\omega P_g
       +C(d+2)N\log((n+1)(d+1)). \tag{GM2}
\]

The Gaussian perturbation of a balanced center of path mass $P_0$
has uniform root-mean-square function error at most
\[
 C\epsilon\sqrt{\nu P_0}+C\epsilon^2\nu. \tag{GM3}
\]
For example, the incoming perturbation has directional standard
deviation at most $\epsilon\sqrt{2\nu}$, and
$\sum_j|\bar v_j|\le\sqrt{P_0}$. The independent readout perturbation
has variance at most $C\epsilon^2\nu P_0/N$ when applied to the center
features. The product term has variance at most
$C\epsilon^4\nu^2/N$. These estimates use
$\|\bar x\|_2^2/d\le2$ and are uniform in input dimension and width.

We record a deterministic posterior moment bound to justify the
uniform error estimate after mixing:
\[
 \Pi Q^2\le C\{(d+2)\nu+\lambda^{-1}\}^2. \tag{GM4}
\]
Indeed the native centered output has variance at most $\nu^2/N$,
so Jensen gives $Z\ge e^{-\omega/\lambda}$.
Also $\mathbb E_\mu e^{\omega Q/4}=2^{(d+2)N/2}$.
Combining these bounds gives an exponential posterior tail for $Q$
above $C((d+2)\nu+\lambda^{-1})$; integrating proves (GM4).
Since $P\le Q/2$, $P_g\le n^4$, $\lambda\ge n^{-2}$ and
$|g|\le\sqrt2 P_g$, (GM3) bounds the loss perturbation by
$C\epsilon\Pi(P+P_g+1)^{3/2}+C\epsilon^2\Pi(P+P_g+1)$.
By (GM4) this is at most $Cn^{-10}$. In particular no bounded-output
restriction is imposed on the comparator.
The exact squared-loss midpoint identity thus gives
\[
 \rho\widehat R\le\tfrac12\Pi\widehat R+\tfrac12\widehat R(g)
                  -\tfrac14\Pi P_n(z-g)^2+Cn^{-10}.
\]
Finally, the actual posterior minimizes
$\rho\widehat R+(2\lambda/\omega)\operatorname{KL}(\rho\Vert\mu)$.
Use this comparison and (GM2), multiply by two, and use
$N/\omega=\nu$. This is (GM1).
\end{proof}

\begin{corollary}[Only the observed label second moment is needed]
\label{cor:general-midpoint-label-moment}
In the preceding lemma replace the bounded-label assumption by
$\widehat{Y^2}\le L$, where $L\ge1$. Then the same inequality holds
with entropy remainder
\[
 C\lambda\nu(d+2)\log((n+1)(d+L+1))+Cn^{-10}. \tag{GM5}
\]
No labels are clipped in either the posterior or this comparison.
The statement is deterministic for every sample satisfying the
observed second-moment bound, so $L$ may also be a number computed
from that sample.
\end{corollary}
\begin{proof}
Use exactly the same midpoint centers and let
$\epsilon=((n+1)(d+L+1))^{-30}$ in the proof-only comparison laws.
The evidence bound becomes
$Z\ge\exp[-\omega(L+1)/(2\lambda)]$.
The same Gaussian energy moment therefore gives
\[
 \Pi Q^2\le C\{(d+2)\nu+(L+1)/\lambda\}^2.
\]
The midpoint identity is valid for arbitrary real labels. To bound
the perturbation in empirical loss, apply Cauchy--Schwarz over the
observations before posterior averaging. Its absolute value is at
most
\[
 C\epsilon\Pi[(P+P_g+\sqrt L+1)\sqrt{P+P_g+1}]
       +C\epsilon^2\Pi(P+P_g+1).
\]
The displayed energy moment, $P_g\le n^4$ and $\lambda\ge n^{-2}$
make this at most $Cn^{-10}$. The KL calculations are unchanged
except for the enlarged logarithm of the covariance precision.
This proves (GM5).
\end{proof}

\section{A sharp shallow Gaussian oracle on arbitrary input laws}
\label{sec:general-shallow-sharp}
This section uses the Gaussian model, normalized path mass $P$, and
empirical midpoint inequality (GM1) of the preceding section. Inputs
have any probability law on $[-1,1]^d$; initially $|Y|\le1$.
There is no density, smoothness, or input-separation assumption.
All constants below are numerical. Put
\[
 b_n=\frac{c_0\log^4(\mathrm e(n+1)(d+1)/\delta)}n,
 \qquad \lambda_n=c_1\{\sqrt{b_n}+(d+2)b_n\},                 \tag{GS1}
\]
where $c_0,c_1$ are sufficiently large fixed numerical constants.
The nontrivial finite-sample theorem assumes $\lambda_n\le1$.
For each fixed $d,\delta$, this holds for all sufficiently large $n$.
The likelihood schedule knows neither the comparison coefficients nor
its path mass or risk. Its physical observation variance remains
$n\lambda_n\nu/N$.

\begin{lemma}[Dimension-free atomic and multiplier bounds]
\label{lem:general-shallow-statistics}
Let
\[
 \mathcal H_d=\{x\mapsto d^{-1/2}(\theta\cdot\bar x)_+:
                              \|\theta\|_2=1\}.
\]
Its elements are bounded by $\sqrt2$, and the absolute convex hull
with coefficient mass at most $p$ has empirical Rademacher complexity
at most $Cp/\sqrt n$, for every input sample.
For any deterministic finite shallow ReLU function $g$ satisfying
$R(g)\le2$, with probability at least $1-\delta/2$,
\[
 \sup_{u\in\mathcal H_d}
 |(\mathbb P_n-\mathbb P)[(g(X)-Y)u(X)]|
 \le C\left\{\sqrt{\frac{\log(4/\delta)}n}
       +(P_g+1)\frac{\log(4/\delta)}n\right\}.              \tag{GS2}
\]
Here $\mathbb P$ denotes the joint input-label law; $P_g$ is path
mass and is unrelated to the probability notation.
\end{lemma}
\begin{proof}
ReLU contraction and $\|\bar x\|_2/\sqrt d\le\sqrt2$ give
\[
 \frac1n\mathbb E_\epsilon\sup_{\|\theta\|=1}
 \left|\sum_i\epsilon_i d^{-1/2}(\theta\cdot\bar X_i)_+\right|
 \le\frac Cn\mathbb E_\epsilon
           \left\|\sum_i\epsilon_i\bar X_i/\sqrt d\right\|_2
 \le C/\sqrt n.
\]
Absolute convexification multiplies this by $p$.

For (GS2), write $a_i=g(X_i)-Y_i$. The weighted contraction inequality
has the same proof, since each scalar map
$s\mapsto a_i(s/a_i)_+$, for $a_i\ne0$, is $1$-Lipschitz and vanishes
at zero; zero $a_i$ contributes zero. Conditional on the data, its
bound is $C\sqrt{\sum_i a_i^2}/n$. Symmetrization and Jensen therefore
bound the expected supremum of the centered empirical process by
$C\sqrt{R(g)/n}\le C/\sqrt n$.
Each summand has variance at most $2R(g)\le4$ and absolute envelope
at most $C(P_g+1)$, because $|g|\le\sqrt2P_g$.
The bounded empirical-process Bernstein inequality (Bousquet's
inequality~\cite{bousquet}) gives, with $t=\log(4/\delta)$, the bound
\[
 C\left[n^{-1/2}+\sqrt{t/n}
       +\sqrt{(P_g+1)t/n^{3/2}}+(P_g+1)t/n\right].
\]
The mixed square-root term is at most a constant times
$\sqrt{t/n}+(P_g+1)t/n$, proving (GS2).
Continuous parameterization on the compact sphere supplies a
countable determining subclass; no measurability issue arises.
\end{proof}

\begin{lemma}[A relative square-norm event]
With probability at least $1-\delta/2$, uniformly over all finite
shallow ReLU functions $h$ of path mass $P_h\le 4n^5$,
\[
 \mathbb E h(X)^2-\mathbb P_n h(X)^2
 \le\tfrac14\mathbb P_n h(X)^2+C b_n(P_h+1)^2.              \tag{GS3}
\]
\end{lemma}
\begin{proof}
Apply the smooth nonnegative-loss Rademacher bound of Srebro,
Sridharan, and Tewari~\cite[Theorem~1]{sst} to square loss with zero labels on path balls
of radii $1,2,4,\ldots,2^{\lceil\log_2(4n^5)\rceil}$.
On a radius-$p$ ball the loss is bounded by $2p^2$, is $2$-smooth,
and the preceding lemma gives worst-case Rademacher complexity
$Cp/\sqrt n$. Allocate confidence equally among these
$O(\log n)$ balls. The resulting upper deviation is bounded by
\[
 C\sqrt{\mathbb P_nh^2}\,p\sqrt{b_n}+C p^2b_n.
\]
Young's inequality and the least dyadic radius containing $h$ prove
(GS3). This event depends only on a function class, so it is common
to every learner width and Gaussian scale.
\end{proof}

\begin{lemma}[A uniform posterior energy cutoff]
For $N\ge n$, $0<\nu\le1$, $n^{-2}\le\lambda\le1$, there is a
numerical constant $A$ such that, for
\[
 q_*=A\{d+2+\lambda^{-1}\},
\]
every sample with $|Y_i|\le1$ satisfies
\[
 \Pi_\lambda[(1+Q^2)\mathbf1_{\{Q>q_*\}}]\le n^{-40}.     \tag{GS4}
\]
\end{lemma}
\begin{proof}
The prior is centered and has pointwise output variance at most
$\nu^2/N$, so its empirical risk is at most two. Consequently
$Z\ge\exp(-\omega/\lambda)$. With $D=(d+2)N$,
$\mathbb E_\mu e^{\omega Q/4}=2^{D/2}$.
The elementary bound
\[
 (1+Q^2)\mathbf1_{Q>q}
 \le C(1+\omega^{-2})e^{\omega Q/4-\omega q/8}
\]
therefore gives
\[
 \Pi_\lambda[(1+Q^2)\mathbf1_{Q>q}]
 \le C(1+\omega^{-2})
       \exp\{\omega/\lambda+D\log2/2-\omega q/8\}.
\]
Use $D/\omega=\nu(d+2)\le d+2$, $\omega\ge n$, and enlarge $A$.
The exponent is then sufficiently negative to imply (GS4),
uniformly in all the stated parameters.
\end{proof}

\begin{theorem}[A coefficient-one path oracle]
\label{thm:general-shallow-sharp}
Suppose (GS1) satisfies $\lambda_n\le1$. For any deterministic finite
shallow ReLU comparator $g$, at every $N\ge n$ and $0<\nu\le1$,
with probability at least $1-\delta$,
\begin{align}
 R(\bar z_{\lambda_n})\le\Pi_{\lambda_n}R(z)
 \le R(g)+C\lambda_n\left[
 P_g+1+(d+2)\log((n+1)(d+1))\right].                     \tag{GS5}
\end{align}
The event can be chosen common to all the eligible learner widths
and Gaussian scales. The prior and prediction are unchanged; the
likelihood scale is deterministic and independent of the comparator.
\end{theorem}
\begin{proof}
First assume $R(g)\le2$ and $\lambda_n(P_g+1)\le c$ for a sufficiently
small fixed constant $c$. In particular $P_g\le n^4$.
Intersect (GS2) and (GS3), and write $h=z-g$ and $P=P(z)$.
Since $h$ has path mass at most $P+P_g$, expand the squared-loss
difference and use these events to obtain
\begin{align*}
 R(z)-R(g)\le{}&\widehat R(z)-\widehat R(g)
       +\tfrac14\mathbb P_n(z-g)^2+C b_n(P+P_g+1)^2\\
       &+C(P+P_g)\{\sqrt{b_n}+(P_g+1)b_n\}.
\end{align*}
For $Q\le q_*$, $P\le q_*/2$. The chosen scale satisfies
\[
 b_nq_*\le c'\lambda_n,\qquad
 \sqrt{b_n}+(P_g+1)b_n\le c'\lambda_n,
\]
with $c'$ as small as needed by increasing $c_1$.
Indeed $\lambda_n\ge c_1\sqrt{b_n}$ and
$\lambda_n\ge c_1(d+2)b_n$; also
$b_n(P_g+1)\le c b_n/\lambda_n$.
The comparator-only terms obey
$b_n(P_g+1)^2\le C\lambda_n(P_g+1)$.
Hence on this energy set
\begin{align*}
 R(z)-R(g)\le{}&\widehat R(z)-\widehat R(g)
       +\tfrac14\mathbb P_n(z-g)^2
       +\tfrac12\lambda_n P+C\lambda_n(P_g+1).
\end{align*}
All path masses here are at most $4n^5$ after increasing a harmless
minimum sample constant; $\lambda_n\le1$ implies $d\le Cn$.

Integrate on $\{Q\le q_*\}$ and use (GM1).
Replacing the restricted empirical integral by its full integral
costs at most $\widehat R(g)\Pi(Q>q_*)$ because the other summands
are nonnegative. The population tail costs at most
$C\Pi[(1+Q^2)\mathbf1_{Q>q_*}]$.
Here $\widehat R(g)\le C(P_g+1)^2$ is polynomial in $n$; (GS4)
makes both corrections negligible. Discard the remaining positive
empirical-distance and path terms in (GM1). This proves (GS5).

For completeness, apply exactly the same proof to $g=0$ (without
requiring $\lambda_n\le c$): its multiplier has bounded envelope,
and all constants are absorbed by (GS1). This gives
\[
 \Pi R(z)\le R(0)+C\lambda_n[1+(d+2)\log((n+1)(d+1))].
\]
If $R(g)>R(0)$ this already implies (GS5).
If $\lambda_n(P_g+1)>c$, use $R(0)\le1$ and enlarge the constant
on $\lambda_n(P_g+1)$ to dominate the possible difference between
$R(0)$ and $R(g)$. This handles all remaining comparators, including
ones with $P_g>n^4$. Jensen gives the posterior-mean conclusion.
\end{proof}

\paragraph{Width and coefficient budgets.}
For an ordinary affine-output shallow teacher of width at most $C$
whose coefficients and biases have magnitude at most $B\ge1$, one
constant neuron implements its output bias and
$P_g\le2d(C+1)B^2$. Thus at fixed $d,B$, (GS5) has remainder
$\widetilde O((C+1)/\sqrt n)$, with no density assumption and no
cap on teacher or learner outputs. More explicitly the remainder is
\[
 \widetilde O\left(
 \left[n^{-1/2}+\frac d n\right]
                  [d(C+1)B^2+d]\right).                 \tag{GS6}
\]
The algorithm does not use $C$ or $B$. For every fixed declared
budget class, choose a deterministic population-risk minimizer before
applying (GS5); one exists by continuity on its compact finite
parameter set. Alternatively, a single deterministic $n^{-10}$-near
minimizer suffices, with its tolerance absorbed in the remainder.
For simultaneous integer budgets, keep the same scale (GS1), apply
the multiplier event to the at most $n^4$ pairs $1\le C,B\le n^2$
with failure allocation $\delta/[2(n+1)^4]$, and enlarge the fixed
constants already used in $b_n$. This changes only its logarithmic
constant, not the scale between classes. Larger budgets have
$\lambda_n d(C+1)B^2\ge c$, so the zero-comparator event gives their
stated resource bound directly. No prior over these classes is used.

\section{Unchanged labels with only a finite second moment}
\label{sec:general-shallow-finitevariance}
The sharp Gaussian shallow oracle also holds under the sole label
assumption $\mathbb E Y^2<\infty$. This is a high-probability result,
not an expectation bound. No labels or network outputs are clipped.

\begin{lemma}[A square-integrable multiplier]
Let $a(X,Y)$ be any fixed square-integrable measurable function.
For the normalized ReLU atoms of the preceding section,
\[
 \mathbb E_S\left[
 \sup_{u\in\mathcal H_d}
 |(\mathbb P_n-\mathbb P)[a(X,Y)u(X)]|^2\right]
 \le \frac{C\mathbb E a(X,Y)^2}{n}.                      \tag{GF1}
\]
\end{lemma}
\begin{proof}
An independent ghost sample and Jensen's inequality symmetrize
the squared supremum, with a numerical factor, to
\[
 \mathbb E_{S,\epsilon}
 \left[\sup_{u\in\mathcal H_d}
       \left|n^{-1}\sum_i\epsilon_i a_i u(X_i)\right|^2\right].
\]
Conditional on the sample, the mean of this signed supremum is at
most $C\sqrt{\sum_i a_i^2}/n$, by the weighted contraction argument
above. Changing one Rademacher sign changes the supremum by at most
$2\sqrt2|a_i|/n$. The bounded-differences variance inequality therefore
bounds its conditional variance by $C\sum_i a_i^2/n^2$.
Adding the squared conditional mean, then averaging, proves (GF1).
Truncation of $a$ followed by the same $L^2$ estimate justifies every
step for general square-integrable multipliers.
\end{proof}

\begin{theorem}[A finite-variance, teacher-independent likelihood schedule]
\label{thm:general-shallow-finitevariance}
Let inputs have any law on $[-1,1]^d$, and suppose
$\mathbb E Y^2\le M_2<\infty$. Fix $0<\delta<1/2$ and set
\begin{align*}
 L&=\frac{8(M_2+1)}\delta,\\
 b&=\frac{c_0\log^4(\mathrm e(n+1)(d+L+1)/\delta)}n,\\
 \lambda&=c_1\{\sqrt{Lb}+(d+2)b\}.
\end{align*}
Assume $\lambda\le1$, with sufficiently large fixed numerical
$c_0,c_1$. For any deterministic finite shallow ReLU comparison $g$,
with probability at least $1-\delta$,
\[
 R(\bar z_\lambda)\le\Pi_\lambda R(z)
 \le R(g)+C\lambda
       [P_g+1+(d+2)\log((n+1)(d+L+1))].                 \tag{GF2}
\]
This event is common to every $N\ge n$ and $0<\nu\le1$.
The prior is the original independent Gaussian prior; the likelihood
uses the actual raw observations, and its physical variance is
$n\lambda\nu/N$. The deterministic schedule depends on the declared
second-moment bound and confidence, but not on the comparator.
\end{theorem}
\begin{proof}
Markov's inequality gives $\mathbb P_nY^2\le L$ except with
probability at most $\delta/8$. On this event the moment version
(GM5) of the empirical midpoint inequality applies, and
\[
 Z\ge\exp[-\omega(L+1)/(2\lambda)],\qquad
 \Pi Q^2\le C[(d+2)\nu+(L+1)/\lambda]^2.
\]
The proof of (GS4), with this evidence bound, gives
\[
 \Pi[(1+Q^2)\mathbf1_{Q>q_*}]\le n^{-40},\qquad
 q_*=A\{d+2+(L+1)/\lambda\}.                           \tag{GF3}
\]
This is a deterministic assertion on the empirical-second-moment
event, simultaneous in width and prior scale.

First consider $R(g)\le R(0)+1\le M_2+1$ and
$\lambda(P_g+1)\le cL$. Apply (GF1) to $a=g-Y$ and Markov's
inequality. Except with probability $\delta/4$, its multiplier
supremum is at most $C\sqrt{(M_2+1)/(n\delta)}\le C\sqrt{Lb}$.
Intersect this with the square-norm event (GS3), with its confidence
allocation reduced by a fixed factor. Expanding the loss difference
gives, uniformly on $Q\le q_*$,
\begin{align*}
 R(z)-R(g)\le{}&\widehat R(z)-\widehat R(g)
       +\tfrac14\mathbb P_n(z-g)^2\\
       &+Cb(P+P_g+1)^2+C(P+P_g)\sqrt{Lb}.
\end{align*}
Now $bq_*\le c'\lambda$ follows from
$\lambda\ge c_1\sqrt{Lb}$ and $\lambda\ge c_1(d+2)b$.
Also $b(P_g+1)^2\le C\lambda(P_g+1)$ in the stated comparison
range, since $\lambda^2\ge c_1^2Lb$.
Enlarging $c_1$ consequently bounds the last line by
$\lambda P/2+C\lambda(P_g+1)$.

Integrate, apply (GM5), and use (GF3) for the discarded tail.
The empirical comparator risk is at most $C(P_g+1)^2+2L$.
Since $\lambda\le1$ forces $L,d\le Cn$ and the nontrivial branch
forces $P_g\le Cn$, these tail corrections are negligible.
This proves (GF2) in that branch.

Apply the same argument to zero to get
\[
 \Pi R(z)\le R(0)+C\lambda
              [1+(d+2)\log((n+1)(d+L+1))].
\]
If $R(g)>R(0)$ this already suffices. If
$\lambda(P_g+1)>cL$, the excess resource term can dominate
$R(0)\le M_2\le L$, after enlarging its numerical constant.
Thus all comparisons are covered. Jensen proves the mean bound.
\end{proof}

\paragraph{No second-moment bound is needed asymptotically.}
Fix input dimension and a confidence $0<\delta<1/2$. The one
deterministic schedule
\[
 \lambda_n=\frac{\log^6(n+2)}{\sqrt n}
                  +(d+2)\frac{\log^6(n+2)}n              \tag{GF4}
\]
eventually satisfies every lower bound in the proof for each fixed
finite $M_2$. Its onset may depend on that unknown label law and on
$\delta$; the quantitative finite-sample theorem above uses a
declared moment bound. At fixed coefficient cap $B$ and dimension,
the raw posterior at (GF4) therefore competes with width
$C_n=\lfloor n^{1/4}\rfloor$ shallow ReLU comparisons with
high-probability gap $\widetilde O(n^{-1/4})$, at every $N\ge n$
and every mean-field-or-stronger scale. More generally any prescribed
$C_nB_n^2=o(\sqrt n/\log^7 n)$ gives a vanishing bound. This is a
bounded-coefficient comparison, not a coefficient-free theorem on
arbitrary input laws.

\paragraph{Fixed classes and optional simultaneous budgets.}
For each declared finite coefficient/width class choose a
deterministic population minimizer, which exists by continuity on
its compact parameter set. This gives its oracle immediately.
The fixed-class result is sufficient for the growing class in (GF4).
One can also keep one event over polynomially many nontrivial budget
classes: apply (GF1) only to the common multiplier $a=Y$, and handle
each bounded $g$ multiplier separately by Bousquet's inequality.
Since $\mathbb E g^2\le C(M_2+1)$ for the chosen near-minimizers,
the resulting simultaneous multiplier coefficient is
$C\{\sqrt{(M_2+1)b}+(P_g+1)b+\sqrt{(M_2+1)/(n\delta)}\}$.
It obeys the same absorption inequalities. A polynomial finite union
changes only the logarithms already included in $b$; larger budgets
are again handled by zero. No union of Markov events over all classes
is needed.

\begin{corollary}[Universal risk consistency in probability]
\label{cor:gaussian-universal-consistency}
Fix the input dimension and any probability law on the input cube,
and assume only $\mathbb E Y^2<\infty$. For any prescribed Gaussian scales $0<\nu_n\le1$ and finite widths
$N_n\to\infty$, however slowly, the ordinary raw posterior mean is
risk-consistent in probability under either of two choices: a shallow
ReLU learner at the universal $\log^8$ schedule of
Section~\ref{sec:all-width-relu}, or a bounded, nonconstant,
nondecreasing Lipschitz activation at $\lambda_n=n^{-3/4}$.
The posterior-average risk has both conclusions. Neither schedule
needs the unknown label second moment.
\end{corollary}
\begin{proof}
Finite shallow ridge networks with any of these continuous
nonpolynomial activations are uniformly dense in continuous functions
on the cube~\cite[Theorem~1]{leshno}. Continuous functions are dense in
$L^2(P_X)$, and the regression function is square-integrable.
Consequently, for every $\varepsilon>0$, there is a fixed finite
comparison network $g$ with
$R(g)\le R_*+\varepsilon$, where
$R_*=\mathbb E\operatorname{Var}(Y\mid X)$.
For this fixed $g$, the ReLU remainder in (GA3) at the universal
schedule of Section~\ref{sec:all-width-relu} tends to zero;
in the bounded-activation case $\mathcal A_n(g)=O(\log n)$ and (AW3)
at $\lambda_n=n^{-3/4}$ also has a vanishing remainder. In both cases the
width requirement is eventually met because the fixed comparison
has finitely many neurons and $N_n\to\infty$.
Both statements hold for every fixed confidence level and uniformly
over the eligible prescribed widths and scales. Since test risk is
at least $R_*$, first choose the approximation error small and then
let the confidence failure probability tend to zero. This proves
convergence in probability. It does not assert expected-risk
consistency for arbitrary finite-variance labels. This is a standard
approximation consequence of the quantitative oracles.
\end{proof}

\section{Sub-Gaussian labels and an expected-risk corollary}
\label{sec:general-shallow-subgaussian-expectation}
The finite-variance confidence theorem does not imply an
expected-risk theorem. For marginally sub-Gaussian labels, however,
the sharper logarithmic-confidence bounded-label argument gives both.
Assume $\mathbb E\exp(Y^2/K_Y^2)\le2$, with a fixed known envelope.
Let $\lambda'_n$ be (GS1) with confidence $\delta/4$, and put
\[
 M_n=\max\{1,K_Y\sqrt{c\log((n+2)/\delta)}\},\qquad
 \lambda_n=M_n^2\lambda'_n,                              \tag{GG1}
\]
for a sufficiently large fixed numerical $c$.
The actual likelihood still uses the original labels and unbounded
raw output, with variance $n\lambda_n\nu/N$.

\begin{corollary}[Marginal sub-Gaussian confidence and expectation]
At every $N\ge n$, $0<\nu\le1$, and every sample size with
$\lambda'_n\le1$, the sharp path and coefficient-budget oracles
(GS5)--(GS6) hold with the same polynomial powers and additional
logarithmic factors, for the raw posterior-average and mean risks.
Their confidence is at least $1-\delta$.
Choosing $\delta=n^{-20}$ in (GG1) also gives an expected-risk
oracle with the same powers. In particular, at fixed dimension and
coefficient cap, every declared width-$C$ comparison class has
\[
 \mathbb E_S R(\bar z_{S})\le
 \mathbb E_S\Pi_S R(z)
 \le \inf_{g\in\mathcal G_{C,B}}R(g)
          +\widetilde O((C+1)B^2/\sqrt n).
\]
The prior, prediction convention, and all observations are unchanged.
This expectation statement is specific to the present sub-Gaussian
label hypothesis; it is not asserted under only a second moment.
\end{corollary}
\begin{proof}
Use the proof-only label $Y_M=\operatorname{clip}(Y,[-M_n,M_n])$.
Except with probability at most $\delta/4$, every observed label
already equals its clipped version. Analyze the observable $z/M_n$
and label $Y_M/M_n$, keeping exactly the original parameter prior.
The likelihood then has normalized scale $\lambda'_n$.
Its prior empirical risk is at most two, and its parameter energy
is still the same chi-square quadratic form.

For a physical comparator $g$, the minimum-mass LP compresses the
physical midpoint $(z+g)/2$ as before, with center energy at most
$P(z)+P_g$. Dividing loss by $M_n^2$ leaves (GM1) valid with
scale $\lambda'_n$ and the same physical path penalty.
The normalized function has path mass $P(z)/M_n\le P(z)$;
the normalized comparator has mass $P_g/M_n\le P_g$.
Consequently the atom, relative-square, and multiplier estimates
used in (GS5) are all valid with these larger physical path bounds.
This proves the normalized oracle. Multiplication by $M_n^2$
preserves its comparison coefficient one and adds only logarithms.

This also transfers population risks. On the clipped-label sample,
\[
 \Pi Q^2\le C[d+1+(\lambda'_n)^{-1}]^2,\qquad
 |z(x)|\le Q/\sqrt2.
\]
Therefore the difference between the posterior's risks under $Y$
and $Y_M$ is at most
\[
 C[d+1+(\lambda'_n)^{-1}]
       \|Y-Y_M\|_{L^2}+\mathbb E|Y^2-Y_M^2|.
\]
It is negligible for sufficiently large $c$ in (GG1).
Comparisons with $R(g)>R(0)$ use the zero comparator at the same
schedule; otherwise $\|g\|_{L^2}\le2\|Y\|_{L^2}$.
For a fixed comparison budget choose a deterministic population minimizer with
$R(g)\le R(0)$; then $\|g\|_{L^2}\le2\|Y\|_{L^2}$.
Its analogous population-risk transfer is negligible by the same
estimate. No input distribution is conditioned on a label event.

It remains to control the exceptional samples for the expectation
claim. For every sample, without any label bound, the Gaussian
evidence and energy estimate give
\[
 \Pi Q^2\le
 C\left[d+1+\frac{1+\mathbb P_nY^2}{\lambda_n}\right]^2.
\]
Thus, writing $M_2=\mathbb EY^2$,
\[
 \Pi R(z)\le C\left[M_2+(d+1)^2+
                \frac{(1+\mathbb P_nY^2)^2}{\lambda_n^2}\right].
\]
The sub-Gaussian assumption gives
$\mathbb E(\mathbb P_nY^2)^4\le\mathbb EY^8<\infty$.
Hence the $L^2$ norm, over the data, of this risk upper bound is at
most $C_{K_Y}[1+d^2+\lambda_n^{-2}]$.
At $\delta=n^{-20}$, Cauchy--Schwarz bounds the expected risk on
the exceptional event by
$C_{K_Y}[1+d^2+\lambda_n^{-2}]n^{-10}$.
In the stated sample-size range, $d\le Cn$ and
$\lambda_n\ge n^{-2}$, so this contribution is $O_{K_Y}(n^{-6})$.
Integrating the confidence bound proves the expected oracle.
\end{proof}

\paragraph{The bounded-monotone companion.}
The same confidence-to-expectation step applies to the paper's
bounded-monotone-activation midpoint oracle of
Theorem~\ref{thm:uncapped-all-d}, with its condition $N\lambda\ge c_\phi$.
This paragraph does not broaden the all-width finite-variance theorem
(AW3) to expectation. Its deterministic
readout-energy bound has the smaller form
$\Pi R(z)\le C_\phi[M_2+1+(1+\mathbb P_nY^2)/\lambda]$.
Thus its already stated logarithmic-confidence sub-Gaussian oracle
also yields an expected-risk version by taking $\delta=n^{-20}$;
the fourth label moment suffices for this final integration step.
The finite-second-moment confidence-only conclusions remain separate.

\section{The shallow ReLU oracle at every width above the comparison width}
\label{sec:all-width-relu}
The preceding ReLU oracle can be strengthened from $N\ge n$ to
every width that can hold the comparison network. The proof uses an
ordinary Gaussian neighborhood of a balanced teacher and a joint
empirical-loss/energy cutoff. It does not require the midpoint LP.
Keep the original independent Gaussian prior and normalized path
mass of the preceding sections. Inputs have any law on $[-1,1]^d$.

\begin{lemma}[A balanced comparison at every eligible width]
Let $g$ be implemented by $k$ shallow ReLU neurons, and let $N\ge k$.
Suppose $\mathbb P_nY^2\le L$, $L\ge1$, $P_g\le n^4$, and
$n^{-2}\le\lambda\le1$. Then
\[
 e+\tfrac12\lambda\Pi Q
 \le2\lambda P_g+C\lambda\nu(d+2)
           \log((n+1)(d+L+1))+Cn^{-10},                 \tag{GA1}
\]
where $e=\Pi\widehat R(z)-\widehat R(g)$ and $\Pi$ is the
actual raw posterior.
\end{lemma}
\begin{proof}
Write $g=\sum_{j=1}^k c_jd^{-1/2}(\theta_j\cdot\bar x)_+$,
with unit directions, so $\sum_j|c_j|=P_g$.
At learner width $N$, implement each term by
\[
 \bar w_j=\sqrt{N|c_j|/d}\,\theta_j,\qquad
 \bar v_j=\operatorname{sign}(c_j)\sqrt{|c_j|/N},
\]
and pad all unused positions with zeros. Positive homogeneity
preserves the function exactly and gives $Q(\bar w)=2P_g$.
Around this center use native Gaussian variances multiplied by
$\epsilon^2$, where $\epsilon=((n+1)(d+L+1))^{-30}$.
The normalized Gaussian relative entropy is at most
\[
 \frac{2\lambda}{\omega}\KL(\rho\|\mu)
 \le2\lambda P_g+2\lambda\nu(d+2)\log(1/\epsilon).
\]
The balanced perturbation estimate (GM3) gives uniform RMS function
error $C\epsilon\sqrt{\nu P_g}+C\epsilon^2\nu$.
Empirical Cauchy--Schwarz and $\mathbb P_nY^2\le L$ consequently give
$\rho\widehat R\le\widehat R(g)+Cn^{-10}$.
On the other hand, the Gaussian energy identity implies
\[
 \KL(\Pi\|\mu)\ge\tfrac14\omega\Pi Q
                       -\tfrac12(d+2)N\log2.
\]
Compare the posterior variational objective to $\rho$ and substitute
these two entropy bounds. This is (GA1).
\end{proof}

\begin{lemma}[A joint cutoff uniform even at narrow widths]
There is a numerical $C$ such that, on $\mathbb P_nY^2\le L$, put
\[
 K=C(d+2)\log(\mathrm e(n+1)(d+L+1)),\qquad
 J=\widehat R(z)+\lambda Q,\qquad
 u=2(L+1)+\lambda K.
\]
For every $N\ge1$, $0<\nu\le1$, and $\lambda>0$,
\[
 \Pi[(1+Q^2)\mathbf1_{\{J>u\}}]\le n^{-40}.             \tag{GA2}
\]
\end{lemma}
\begin{proof}
The centered native output has variance at most $\nu^2/N\le1$,
so $Z\ge\exp[-\omega(L+1)/(2\lambda)]$.
On $J>u$,
\[
 e^{-\omega\widehat R/(2\lambda)}
 \le e^{-\omega u/(4\lambda)}e^{\omega Q/4}.
\]
Writing $D=N(d+2)$, exponential tilting of $\omega Q\sim\chi_D^2$
therefore gives
\begin{align*}
 \Pi[(1+Q^2)\mathbf1_{J>u}]
 &\le C[1+(d+2)^2]
 \exp\left\{\frac{\omega(L+1)}{2\lambda}
             -\frac{\omega u}{4\lambda}+\frac{D\log2}{2}\right\}.
\end{align*}
Indeed the tilted second moment is $4D(D+2)/\omega^2$ and
$D/\omega=\nu(d+2)\le d+2$.
Use $\omega=N/\nu\ge1$ and increase $C$ in $K$.
No lower bound $N\ge n$ is used.
\end{proof}

\begin{theorem}[An all-width, finite-variance raw Gaussian oracle]
\label{thm:all-width-relu}
Let $\mathbb EY^2\le M_2<\infty$ and $0<\delta<1/2$. Set
\begin{align*}
 L&=8(M_2+1)/\delta,\\
 b&=c_0\log^4(\mathrm e(n+1)(d+L+1)/\delta)/n,\\
 K&=c_2(d+2)\log(\mathrm e(n+1)(d+L+1)),\\
 \lambda&=c_1\{\sqrt{(L+1)b}+K b\},
\end{align*}
with sufficiently large fixed numerical constants, and assume
$\lambda\le1$. For every deterministic $k$-neuron shallow ReLU
comparison $g$, with probability at least $1-\delta$,
simultaneously for every $N\ge k$ and $0<\nu\le1$,
\[
 R(\bar z_\lambda)\le\Pi_\lambda R(z)
 \le R(g)+C\lambda(P_g+K+1).                            \tag{GA3}
\]
Labels, outputs and prediction are unchanged. The actual Gaussian
observation variance is $n\lambda\nu/N$. The schedule uses neither
comparison width nor coefficient magnitudes nor approximation risk.
\end{theorem}
\begin{proof}
Markov's inequality gives $\mathbb P_nY^2\le L$ except with
probability $\delta/8$. Lemma (GF1), applied only to the common
multiplier $Y$, gives, except with probability $\delta/8$,
\[
 Z_Y:=\sup_{a\in\mathcal H_d}
       |(\mathbb P_n-\mathbb P)[Ya]|\le C\sqrt{M_2/(n\delta)}.
                                                               \tag{GA4}
\]
Apply the smooth-square bound once to the unit path ball, and
rescale each nonzero $f$ by its own path mass. This gives one event
of failure probability $\delta/4$ on which
\[
 \mathbb E f^2-\mathbb P_nf^2
 \le C\sqrt{\mathbb P_nf^2}\,P_f\sqrt b
                         +CbP_f^2                            \tag{GA5}
\]
for every finite shallow function. The zero function is immediate.
This is a function-class event, independent of learner width and
energy; no energy shells or label moment enter this step.

On the joint-good set $J\le u$ of (GA2),
\[
 Q\le u/\lambda,\qquad
 \mathbb P_n z^2\le2\widehat R(z)+2\mathbb P_nY^2
                       \le C(L+1+\lambda K).
\]
Since $P\le Q/2$, the chosen scale obeys
\[
 \sqrt{(L+1+\lambda K)b}\le c\lambda,
 \qquad b(u/\lambda+1)\le c\lambda,
 \qquad Z_Y\le c\lambda,
\]
with $c$ as small as needed by increasing $c_1$.
Using (GA4)--(GA5), the exact cancellation of the $Y^2$ term gives
\begin{align*}
 R(z)-R(0)
 &\le\widehat R(z)-\widehat R(0)
         +\tfrac14\lambda Q+C\lambda.                    \tag{GA6}
\end{align*}
All path radii used here are polynomial in $n$: $\lambda\le1$
implies $d,L\le Cn$ and $u/\lambda\le Cn$.

Integrate (GA6) on $J\le u$. To replace its empirical integral by
the full one, use that its integrand is bounded below by
$-\widehat R(0)\ge-L$. On the complement,
$R(z)\le C Q^2+2M_2$. Thus (GA2) controls both tail corrections.
Combining with (GA1) yields
\[
 \Pi R(z)-R(0)
 \le\widehat R(g)-\widehat R(0)
                 +C\lambda(P_g+K+1)+Cn^{-10}.             \tag{GA7}
\]
First suppose $R(g)\le R(0)+1$ and $\lambda(P_g+1)\le cL$.
Then $\mathbb Eg^2\le C(M_2+1)$ and $|g|\le\sqrt2P_g$.
Bernstein's inequality for the bounded variable $g(X)^2$ gives,
except with probability $\delta/4$,
\[
 |(\mathbb P_n-\mathbb P)g^2|
 \le CP_g\sqrt{(M_2+1)\log(8/\delta)/n}
                 +CP_g^2\log(8/\delta)/n.
\]
The $Yg$ discrepancy is at most $P_gZ_Y$ by absolute convexity.
Consequently
\[
 \widehat R(g)-\widehat R(0)
 \le R(g)-R(0)+C\lambda P_g,
\]
where the quadratic comparator term is absorbed using
$bP_g^2\le cLbP_g/\lambda\le C\lambda P_g$.
Also $P_g\le n^4$, as required in (GA1).
This proves (GA3) in the nontrivial comparison range.

The same argument with $g=0$ gives a bound by $R(0)+C\lambda(K+1)$
at every $N\ge1$. It covers comparisons with $R(g)>R(0)$ directly,
and covers $\lambda(P_g+1)>cL$ since $R(0)\le M_2\le L$.
Finally Jensen gives the posterior-mean claim.
\end{proof}

\paragraph{Interpretation and growing comparisons.}
An affine-output width-$C$ teacher with coefficient cap $B\ge1$
has an implementation with at most $C+1$ neurons and
$P_g\le2d(C+1)B^2$. Thus at fixed $d,M_2,B,\delta$, (GA3) gives
$\widetilde O((C+1)B^2/\sqrt n)$ at every $N\ge C+1$.
If no output bias is present, $N\ge C$ suffices.
There is no input-regularity assumption.
For each declared compact budget class choose a deterministic
population minimizer; (GA3) then compares to its best risk.
The common multiplier $Y$ is used only once, and polynomially many
chosen minimizers can share the event by a Bernstein union bound
with logarithmic cost. Classes whose displayed resource bound is
already larger than $L$ are covered by zero.

For each fixed dimension, moment-finite label law and confidence,
the universal schedule
$\lambda_n=\log^8(n+2)/\sqrt n+(d+2)\log^8(n+2)/n$
eventually satisfies all the required inequalities without a supplied
moment bound. At fixed $B$, width $\lfloor n^{1/4}\rfloor$
comparisons then have gap $\widetilde O(n^{-1/4})$ already for every
$N\ge\lfloor n^{1/4}\rfloor+1$.
The onset may depend on the unknown law; this is not a uniform
finite-sample bound over all finite-second-moment laws.

\paragraph{Sub-Gaussian confidence and expectation.}
Under a marginal sub-Gaussian envelope, take the proof-only label
cutoff $M_n$ from (GG1) and physical scale $M_n^2\lambda'_n$, with
the normalized scale chosen below.
For the normalized bounded labels, use $L=1$ deterministically and
bounded-process Bernstein for the common $Y$ multiplier, giving
$Z_Y\le C\sqrt b$. Thus one may choose
$\lambda'_n=C(\sqrt b+Kb)$ with only logarithmic confidence
dependence; the finite-variance $\delta^{-1}$ Markov factor is not
used. The normalized-output argument is
unchanged: the actual prior and the physical balanced teacher center
are retained. Thus the same all-width oracle has only logarithmic
confidence dependence and the same polynomial sample rates.
The deterministic raw energy moment bound used in the preceding
expectation section is uniform for every $N\ge1$; taking
$\delta=n^{-20}$ therefore also gives the all-width expected-risk
oracle under sub-Gaussian labels. Only the finite-variance version
remains confidence-only.

\section{A sharper one-dimensional result with a shallow ReLU learner}
\label{sec:univariate-Gaussian}
In one input dimension, a one-hidden-layer ReLU learner gives a much
better width exponent while retaining independent Gaussian weights,
raw likelihood, and raw prediction. This section changes the learner
activation to ReLU and uses exactly one hidden layer:
\[
 z_w(x)=\sum_{j=1}^Nv_j\operatorname{ReLU}(a_jx+b_j),\qquad
 a_j,b_j\sim N(0,\nu),\quad v_j\sim N(0,\nu/N^2),\quad0<\nu\le1.
\]
All coordinates are independent. Assume $X\in[-1,1]$ and
\begin{equation}
 P_X(I)\le\kappa |I|^\alpha
 \quad\hbox{for every interval }I\subseteq[-1,1],
 \qquad\kappa\ge1,\quad0<\alpha\le1.
 \label{eq:one-dimensional-regularity}
\end{equation}
A bounded input density gives $\alpha=1$.

\begin{lemma}[Shallow Gaussian path-variation oracle]
\label{lem:shallow-path-interface}
For a deterministic shallow ReLU implementation $q$, put
\[
 P(q)=\sum_j|v_j|\sqrt{a_j^2+b_j^2}.
\]
There is an observable evidence calibration independent of this
implementation such that, with probability at least $1-\delta$,
simultaneously for all $N$ above its width and all $0<\nu\le1$,
\[
 \Pi_SR(z)\le R(q)+\widetilde O\left(
 (P(q)+1)\sqrt{(R(q)+1)/n}+(P(q)+1)^2/n\right).
 \label{eq:one-dimensional-path-oracle}
\]
Labels may have the same marginal sub-Gaussian envelope as before.
The ordinary raw mean obeys the same bound. Logarithms include
$n,P(q),1/\delta$; constants depend on the fixed label envelope.
\end{lemma}
\begin{proof}
This is the path-variation specialization of the Gaussian free-energy
oracle~\cite{adaptive}; we give the normalization and comparison
calculation to specify its application to the present prior.
Set $\omega=N/\nu$ and
\[
 Q=N^{-1}\sum_j(a_j^2+b_j^2)+N\sum_jv_j^2.
\]
Then $\omega Q\sim\chi^2_{3N}$ and $3N/\omega\le3$.
Cauchy--Schwarz gives $P(w)\le Q/2$. The input vector $(x,1)$ has
norm at most $\sqrt2$, so the output envelope is at most $Q/\sqrt2$,
and the energy-ball Rademacher complexity is at most $cQ/\sqrt n$.
The latter follows from ReLU contraction and absolute convexification
of Euclidean unit linear features. The native pointwise variance is
at most one, so its empirical prior risk under bounded unit labels
is at most two.

Positive rescaling of each comparison neuron preserves its function
and balances its two energy contributions. Zero-padding then gives an
exact width-$N$ mean implementation with $Q=2P(q)$ at every eligible
$N$. Shrink all native Gaussian coordinate variances by $\eta^2$.
For the resulting comparison law,
\[
 (2/\omega)\KL(\rho\|\mu)
       \le2P(q)+6\log(1/\eta),
\]
and independence of the first-layer and readout perturbations gives
\[
 \bigl(\E_\rho\|z_w-q\|_\infty^2\bigr)^{1/2}
       \le4\eta\sqrt{P(q)}+2\eta^2.
\]
Taking $\eta=[64n(P(q)+1)^2]^{-1}$ makes the comparison empirical
squared-loss increment at most $c/n$ for bounded labels.
The free-energy certificate with capacity radius linear in $Q$ has
powers $(v,\rho)=(2,1)$ and comparison resource
$\widetilde O(P(q)+1)$. Its scalar root is
\[
 (\lambda')^2=\theta_n[2F'_S(\lambda')+b_n\lambda'],
 \qquad \theta_n=\widetilde O(1/n),\quad b_n=O(\log(n/\delta)),
\]
with the same bounded-risk fallback for small samples. Apply it to
the normalized observable and labels from
Section~\ref{sec:statistical-interface}, and set
$\lambda=M_n^2\lambda'$. The physical observation variance is still
$n\lambda\nu/N$. The normalized-observable transfer changes only
logarithmic factors. This proves \eqref{eq:one-dimensional-path-oracle}.
\end{proof}

\begin{lemma}[A bounded spline with controlled path variation]
\label{lem:spline-path}
If a continuous $f:[-1,1]\to[-1,1]$ has at most $K$ affine
breakpoints, then for every $0<t\le1$ there is a bounded continuous
piecewise-affine $q$ with at most $2K$ breakpoints such that
\[
 \E|q(X)-f(X)|\le6\kappa Kt^\alpha,\qquad
 \operatorname{Lip}(q)\le2/t,
 \qquad P(q)\le c(K+1)/t.
\]
It has a one-hidden-layer ReLU implementation of width at most $2K+2$.
\end{lemma}
\begin{proof}
Group consecutive breakpoints into maximal clusters whose successive
gaps are at most $3t$. For a cluster running from $a$ to $b$, replace
$f$ on $[\max(-1,a-t),\min(1,b+t)]$ by the chord between its endpoint
values. These replacement intervals are disjoint, with gaps greater
than $t$. A cluster containing $k$ breakpoints has interval length at
most $3kt$. The function remains bounded by one, and differs from $f$
only on these intervals. Equation~\eqref{eq:one-dimensional-regularity}
therefore gives error at most
$2\kappa\sum(3kt)^\alpha\le6\kappa Kt^\alpha$.

Each replacement interval has length at least $t$, so its chord slope
is at most $2/t$. Complementary interior intervals have length greater
than $t$ and are affine. At the boundary, any short complementary
interval belongs to an original affine segment extending at least
$t$ to the first original breakpoint; its slope has the same bound.
Thus all slopes have magnitude at most $2/t$.

If the new breakpoints are $u_1<\cdots<u_r$ and the successive slopes
are $s_0,\ldots,s_r$, use the exact representation
\[
 q(x)=q(-1)\operatorname{ReLU}(1)
     +s_0\operatorname{ReLU}(x+1)
     +\sum_{j=1}^r(s_j-s_{j-1})\operatorname{ReLU}(x-u_j).
\]
Here $r\le2K$, the incoming row norms are at most $\sqrt2$, and each
slope jump is at most $4/t$. This proves the path-variation bound.
The case $K=0$ is a bounded affine function and follows directly.
\end{proof}

\begin{theorem}[A raw $C n^{-1/4}$ oracle under a bounded density]
\label{thm:univariate-Gaussian}
Let $\mathcal G_C$ be all one-hidden-layer univariate ReLU networks of
width at most $C$, with arbitrary finite coefficients and biases.
Under \eqref{eq:one-dimensional-regularity} and the marginal
sub-Gaussian label assumption, the shallow Gaussian learner above
has, with probability at least $1-\delta$, simultaneously for all
$N\ge6C+6$ and all $0<\nu\le1$,
\begin{align*}
 R(\bar z_S)\le\Pi_SR(z)
 \le\inf_{g\in\mathcal G_C}R(g)+\widetilde O\bigl(
 &(C+1)\kappa^{1/(\alpha+1)}n^{-\alpha/[2(\alpha+1)]}\\
 &+(C+1)^2\kappa^{2/(\alpha+1)}n^{-\alpha/(\alpha+1)}\bigr).
\end{align*}
In particular, for a bounded density ($\alpha=1$), the excess is
$\widetilde O(C\sqrt\kappa\,n^{-1/4}+C^2\kappa n^{-1/2})$.
The same single calibration competes with every declared comparison
width; it does not use $C,\kappa,\alpha$ or the original coefficients.
A sufficient fixed-accuracy sample bound is
\[
 n\ge\widetilde O\left(
 \kappa^{2/\alpha}
       ((C+1)/\varepsilon)^{2+2/\alpha}\right).
\]
\end{theorem}
\begin{proof}
Choose a deterministic $n^{-2}$-minimizer $g$ of the comparison risk.
The normalized clipped function
$f=\clip(g/M_n)$ has at most $K=3C+2$ breakpoints: each of the
original $C+1$ affine intervals acquires at most two clipping crossings.
Apply Lemma~\ref{lem:spline-path} at
\[
 t=(\kappa\sqrt n)^{-1/(\alpha+1)}\le1,
\]
and multiply the result by $M_n$. The compiled raw network has width
at most $6C+6$, supremum at most $M_n$, and path variation at most
$cM_n(C+1)/t$. Lemma~\ref{lem:compiler-risk} gives a risk increment
$\widetilde O(\kappa(C+1)t^\alpha)$, plus a negligible label-tail term.
Its risk is also at most $O(M_n^2+\E Y^2)$ independently of $C$.
Substitute this path bound into
\eqref{eq:one-dimensional-path-oracle}. The additional terms are
\[
 \widetilde O\left(\frac{C+1}{t\sqrt n}
                         +\frac{(C+1)^2}{t^2n}\right).
\]
The displayed choice of $t$ balances the two leading contributions
and proves the oracle. Solving its first term for sample size gives
the final bound; its squared term is then smaller. Countably many
declared comparison widths can share one event through a summable
confidence allocation, changing only logarithms.
\end{proof}

\section{An exact Laplace representation of the Gaussian function prior}
\label{sec:Gaussian-Laplace}
The one-dimensional Gaussian model admits a simple exact
reparametrization. This changes neither its prior on functions nor
its likelihood or posterior.

\begin{lemma}[Polar representation]
\label{lem:Gaussian-Laplace}
For independent coordinates
\[
 a_j,b_j\sim N(0,\nu),\qquad v_j\sim N(0,\nu/N^2),
\]
write $(a_j,b_j)=r_j(\cos\theta_j,\sin\theta_j)$ and $c_j=v_jr_j$.
Then the pairs $(c_j,\theta_j)$ are independent across neurons,
$\theta_j$ is uniform on $[0,2\pi)$, and $c_j$ is independent of
$\theta_j$ with Laplace density
\[
 \frac{\omega}{2}e^{-\omega|c_j|},
 \qquad \omega=N/\nu.
\]
Consequently the original Gaussian function prior is exactly
\[
 z(x)=\sum_{j=1}^N
        c_j\operatorname{ReLU}(x\cos\theta_j+\sin\theta_j).
\]
Conditional on the directions, the actual posterior coefficient
density is
\[
 \Pi(dc\mid\theta,S)\ \propto\
 \exp\left[-\omega\left\{
       \sum_{j=1}^N|c_j|+\frac{\widehat R_S(z)}{2\lambda}
                         \right\}\right]\,dc.
\]
Thus the conditional coefficient posterior is a Bayesian lasso;
the directions retain their full Bayesian posterior as well.
\end{lemma}
\begin{proof}
Rotational symmetry gives independent uniform angle and
$r_j^2$ exponential with mean $2\nu$. Conditioning on $r_j$,
\[
 \mathbb E e^{itc_j}
 =\mathbb E\exp\left(-\frac{\nu t^2r_j^2}{2N^2}\right)
 =\frac1{1+\nu^2t^2/N^2},
\]
the characteristic function of a centered Laplace variable with
scale $\nu/N$. Independence of the angle is preserved by the
product $v_jr_j$. Positive homogeneity gives the function identity.
The likelihood is a function only of these transformed variables,
so its posterior pushforward is exact; multiplying the product
Laplace density by the original Gaussian likelihood gives the last
display.
\end{proof}

The path variation of this representation is exactly
$\sum_j|c_j|$, whose prior expectation is $\nu$, independently of
width. Its second-derivative variation on $[-1,1]$ is at most that
quantity. This identity explains why the localized analysis can
use a variation budget. It does not assert that the full posterior
is a fixed-feature lasso, since its directions are also inferred,
and it is not an additional prior restriction.

\section{A subquadratic width polynomial under a two-sided density bound}
\label{sec:localized-results}
The preceding theorem needs only an upper density bound. If the input
law has density bounded away from zero as well, localization gives a
stronger guarantee for the same shallow ReLU learner. Throughout this
section assume
\[
 0<p_-\le p_X(x)\le p_+<\infty\quad(-1\le x\le1).
\]
Initially $|Y|\le1$. Write $f_0(x)=\E[Y\mid X=x]$,
$R_0=\E(Y-f_0(X))^2$, and $A(g)=R(g)-R_0$.
The comparison class $\mathcal G_C$ contains all shallow width-$C$
ReLU functions with arbitrary finite coefficients. All conclusions
below control the posterior-average raw risk and therefore the
ordinary raw posterior mean. The prior and observations are unchanged.
Only the likelihood scale differs from the evidence calibration used
in the preceding sections.

\begin{theorem}[A localized width-only Gaussian oracle]
\label{thm:localized-width}
Let $1\le C\le n$ be a declared comparison-width budget. There is a
deterministic likelihood scale depending on $n,C,\delta$ and the stated
density bounds, but not on the comparison coefficients or its risk,
such that, with probability at least $1-\delta$, simultaneously for
every $N\ge6C+6$ and $0<\nu\le1$,
\begin{equation}
 R(\bar z_S)\le\Pi_SR(z)
 \le\inf_{g\in\mathcal G_C}R(g)
       +\widetilde O\left((C+1)^{2/5}n^{-3/10}\right).
 \label{eq:localized-budget}
\end{equation}
In the nontrivial regime, put
\[
 e_{n,C}=\widetilde c(C+1)^{2/5}n^{-3/10},\qquad
 \lambda\asymp\left(\frac{e_{n,C}}{C+1}\right)^2,
\]
with sufficiently large fixed logarithmic factors. The physical
Gaussian observation variance is $n\lambda\nu/N$.
A sufficient sample size for excess risk $\varepsilon$ is
\begin{equation}
 n\ge\widetilde O\left((C+1)^{4/3}\varepsilon^{-10/3}\right).
 \label{eq:localized-samples}
\end{equation}
Thus the width exponent at fixed accuracy is $4/3$. Any prescribed
$C_n=n^u$, $u<3/4$, gives a vanishing additive remainder with the
single deterministic scale
$\lambda_n=\widetilde\Theta(n^{-3/5-6u/5})$.
\end{theorem}

\begin{corollary}[Schedules that do not use the comparison width]
\label{cor:localized-universal}
One sufficiently large fixed power $k$ in
$\lambda=n^{-3/2}\log^k(n/\delta)$ gives, on a common event for all
comparison widths satisfying $(C+1)\sqrt\lambda\le1$, all
$N\ge6C+6$, and all $0<\nu\le1$,
\begin{align}
 R(\bar z_S)\le\Pi_SR(z)
 \le{}&R_0+\{1+O(1/\log n)\}\inf_{g\in\mathcal G_C}A(g)\\
 &+\widetilde O((C+1)n^{-3/4}+n^{-1/3})+O(1/\log n).
 \label{eq:localized-universal}
\end{align}
Hence every fixed growth $C_n=n^u$, $u<3/4$, is allowed by one
schedule without knowing $C$. Its additive logarithmic term is
weaker in accuracy than the declared-budget polynomial bound.
A different universal schedule,
$\lambda=\widetilde\Theta(n^{-3/5})$, gives the coefficient-one bound
\[
 R(\bar z_S)\le\Pi_SR(z)\le\inf_{g\in\mathcal G_C}R(g)
       +\widetilde O\left((C+1)n^{-3/10}
                          +(C+1)^2n^{-9/10}\right).
\]
\end{corollary}

The bounds are asymptotic statements in the displayed nontrivial
regimes; the proof gives explicit sufficiently large logarithmic
powers and sample conditions. A fixed $\lambda=1$ handles the finitely
many excluded small sample sizes by the same Gaussian energy bound.
There is no mixture over comparison widths or likelihood scales.
The declared-budget schedule in Theorem~\ref{thm:localized-width} uses
$C$; the schedules in Corollary~\ref{cor:localized-universal} do not.

\paragraph{Proof structure.}
The spline compiler in Lemma~\ref{lem:spline-path} has bounded output,
risk increment $O(Ct)$ and path variation $O((C+1)/t)$. The new step
controls empirical quadratic norms through the total variation of the
second derivative, rather than a global output envelope.
Section~\ref{sec:sharp-design} proves this norm comparison with
remainder $\widetilde O(\eta^{-5}n^{-3}P^2)$.
Section~\ref{sec:localized-proof} combines it with localized multiplier
bounds and the exact Gaussian posterior energy identity. Substitution
of the stated schedules proves the theorem and corollary. The common
sample events are function-class events independent of $N,\nu$.
For each declared $C$, choose a deterministic approximate risk minimizer;
a summable confidence allocation over the integer budgets makes the
claims simultaneous, changing only logarithms.

\begin{corollary}[Marginally sub-Gaussian labels]
\label{cor:localized-subgaussian}
Suppose instead that $\E\exp(Y^2/s_Y^2)\le2$ for a fixed known upper
bound $s_Y$. The coefficient-one conclusions
\eqref{eq:localized-budget}--\eqref{eq:localized-samples} and the second
conclusion of Corollary~\ref{cor:localized-universal} hold with additional
logarithmic factors and constants depending on $s_Y$. The first
conclusion holds in the possibly weaker form
\[
 R(\bar z_S)\le\Pi_SR(z)
 \le\{1+O(1/\log n)\}\inf_{g\in\mathcal G_C}R(g)
       +\widetilde O((C+1)n^{-3/4}+n^{-1/3})+O(1/\log n).
\]
Observed labels and the returned predictions remain unrestricted.
The physical likelihood scale is $\lambda=M_n^2\lambda'$, where
$M_n=\max\{1,s_Y\sqrt{c\log(n/\delta)}\}$ and $\lambda'$ is the
corresponding bounded-label schedule, with enlarged logarithmic factors.
\end{corollary}
\begin{proof}
Set $Y_M=\clip(Y,[-M_n,M_n])$ only in the analysis. Its joint law
with $X$ has exactly the original input density. With probability
at least $1-\delta/4$, every training label already lies in this
interval, so the actual posterior is exactly the posterior for the
clipped-label sample. Analyze its observable $z/M_n$ and labels
$Y_M/M_n$. The parameter prior stays unchanged. Its path variation is
at most $Q/2$, its prior empirical risk is at most two, and its Gaussian
energy identity is the same as in Section~\ref{sec:localized-proof}.
A normalized bounded spline comparator $q$ is implemented physically
as $M_nq$; its balanced energy is $2M_nP(q)$. The extra factor $M_n$
and the Gaussian comparison precision cost are logarithmic in $n/\delta$.
Thus every step of the localized argument has the same powers of $n,C$,
with enlarged logarithmic factors. Multiplication by $M_n^2$ converts
normalized risks back to physical risks without changing the leading
coefficient. For the universal schedule, enlarge the fixed logarithmic
power further so that its normalized design remainder is at most
$1/(M_n^2\log n)$; its physical contribution is then still
$O(1/\log n)$. The multiplicative approximation slack stays
$O(1/\log n)$ under risk rescaling.

For completeness, this is a transfer of population risks as well as
training posteriors. The deterministic bound
$\Pi_SQ^2\le C(1+(\lambda')^{-2})$ holds on the clipped training sample.
Since $|z(x)|\le Q/\sqrt2$,
\[
 |\Pi_SR_Y(z)-\Pi_SR_{Y_M}(z)|
 \le C(1+(\lambda')^{-1})\{\E(Y-Y_M)^2\}^{1/2}
           +\E|Y^2-Y_M^2|.
\]
The right side is negligible for a sufficiently large fixed $c$ in
$M_n$, uniformly over the learner width and scale. No law is conditioned
on a label event. For an arbitrary raw comparison function $g$, the
pointwise projection inequality
\[
 (\clip(g,[-M_n,M_n])-Y)^2
 \le(g-Y)^2+(|Y|-M_n)_+^2
\]
first bounds its clipped risk by $R(g)$ plus a label-tail term.
Changing that bounded comparator's labels from $Y$ to $Y_M$ costs
at most $2M_n\E|Y-Y_M|+\E|Y^2-Y_M^2|$.
The clipped comparator still has at most $3C+2$ affine breakpoints.
These observations justify the same coefficient-free compiler and
risk infimum. A union with the clipped-label statistical event and
the all-training-labels event completes the proof.
\end{proof}

\section{An empirical norm comparison for functions with bounded slope variation}
\label{sec:sharp-design}
Assume $0<p_-\le p(x)\le p_+<\infty$ on $[-1,1]$.
For each fixed $0<\eta\le\eta_0$, and for continuous $f$ with finite distributional second-derivative
variation $V=\|f''\|_{\rm TV}$, there is a common event on which
\[
 |P_nf^2-P_Xf^2|
 \le C\eta P_Xf^2+
 C\eta^{-5}\left(\frac{\log(Cn/\delta)}n\right)^3V^2,
 \qquad 0<\eta\le\eta_0, \tag{S}
\]
provided $n\eta^2\ge C\log(Cn/\delta)$.
The event is uniform over all such $f$, with no bound on its affine
part. Constants depend only on $p_-,p_+$.

\begin{proof}
Use $m$ equal bins, $\Delta=2/m$, the interpolant $s=I_m f$, and
$r=f-s$. Let $t=C\log(Cn/\delta)$ and choose
$m=\lfloor c n\eta^2/t\rfloor\ge1$. The usual matrix-Chernoff event
for the hat-function space gives
\[
 |(P_n-P_X)s^2|\le\eta P_Xs^2. \tag{S1}
\]
For completeness, the population Gram matrix of the $m+1$ linear hat
functions has all eigenvalues between $c/m$ and $C/m$: integrate the
square of a linear interpolant on each bin, then use the two density
bounds. The vector of hat values has at most two nonzero entries, each
at most one. After population whitening its squared norm is at most
$Cm$. Matrix Chernoff therefore gives (S1) when
$n\ge C m\eta^{-2}\log(Cm/\delta)$, as ensured by the chosen $m$.
The population Gram bound also gives
\[
 \max_j|s(t_j)|\le C\sqrt m\,\|s\|_{L^2(P_X)}. \tag{S2}
\]

For a bin $I=[a,b]$, define the zero-boundary Green function
\[
 G_u(x)=(x-u)_+
       -\frac{b-u}{b-a}(x-a),\quad x\in I,\ u\in I,
\]
and set it to zero outside $I$. Then
$r(x)=\int_I G_u(x)\,\mu(du)$ on $I$, where $\mu=f''$.
The sign of this formula is immaterial. We have
$|G_u|\le\Delta$, and $u\mapsto G_u$ is uniformly $2$-Lipschitz.
Atoms at grid endpoints have zero Green function and do not affect
the residual. Let $\psi_{I,0},\psi_{I,1}$ be the two affine endpoint
hat functions restricted to $I$.

Bernstein's inequality and a net in $u$ give, simultaneously over all
bins, their two endpoint functions and all $u\in I$,
\[
 |(P_n-P_X)(\psi_{I,j}G_u)|
 \le C\{\Delta^{3/2}\sqrt{t/n}+\Delta t/n\}. \tag{S3}
\]
Indeed the summands have magnitude at most $\Delta$ and variance
at most $C\Delta^3$. A net of spacing $\Delta n^{-10}$ has polynomial
size in $n$ in each bin, and the uniform Lipschitz bound controls
interpolation from the net. The negligible net error can be
absorbed in the displayed right side.

Similarly, a two-dimensional net in $u,v$ gives
\[
 |(P_n-P_X)(G_uG_v)|
 \le C\{\Delta^{5/2}\sqrt{t/n}+\Delta^2t/n\}
 \quad(u,v\in I), \tag{S4}
\]
because the magnitude is at most $\Delta^2$, the variance is at
most $C\Delta^5$, and the parameter Lipschitz bound is at most
$C\Delta$. Kernels in different bins have disjoint supports.
The union of (S1), (S3), and (S4) has probability at least $1-\delta$
after choosing constants in $t$ and in the sample-size condition.

Let $V_I=|\mu|(I^\circ)$, so $\sum_I V_I\le V$.
Equations (S2)--(S3) imply
\[
 |(P_n-P_X)(sr)|
 \le C V\|s\|_{L^2(P_X)}
       \left\{m^{-1}\sqrt{t/n}+m^{-1/2}t/n\right\}. \tag{S5}
\]
The two endpoint coefficients on a bin are bounded by the left
side of (S2); summing the signed-measure integrals therefore gives
(S5). Equation (S4) and $\sum_I V_I^2\le V^2$ give
\[
 |(P_n-P_X)r^2|
 \le C V^2\{m^{-5/2}\sqrt{t/n}+m^{-2}t/n\}. \tag{S6}
\]
The deterministic interpolation bound is
\[
 P_Xr^2\le C V^2m^{-3}. \tag{S7}
\]

Put $a=m^{-1}\sqrt{t/n}+m^{-1/2}t/n$.
By Young's inequality, (S5) is bounded by
$\eta P_Xs^2+C\eta^{-1}V^2 a^2$.
Since $mt/n\le c\eta^2$,
\[
 \eta^{-1}a^2\le C\eta m^{-3},\qquad
 m^{-5/2}\sqrt{t/n}+m^{-2}t/n\le C\eta m^{-3}.
\]
Thus (S1), (S5), and (S6) yield
\[
 |(P_n-P_X)f^2|
 \le C\eta P_Xs^2+C\eta V^2m^{-3}.
\]
Finally $P_Xs^2\le2P_Xf^2+2P_Xr^2$, so (S7) proves the same bound
with $f$ in place of $s$. Substituting $m\asymp n\eta^2/t$ gives
the remainder $C\eta^{-5}(t/n)^3V^2$ in (S).
\end{proof}

\section{Localized losses and the Gaussian posterior}
\label{sec:localized-proof}
This section uses the common design-norm event in Section~\ref{sec:sharp-design}.
The statements below use bounded labels $|Y|\le1$ and a density on
$[-1,1]$ bounded above and below by positive constants. Put
$f_0(x)=\mathbb E[Y\mid X=x]$, $A(f)=P_X(f-f_0)^2$,
and $\xi=Y-f_0(X)$. In particular $|\xi|\le2$.

\paragraph{A self-contained entropy bound.}
For $P\ge1$, let $\mathcal V_P$ contain continuous functions with
supremum, Lipschitz constant and total variation of their distributional
second derivative all at most $cP$, where $c$ is an absolute constant.
Then
\[
 \log N_\infty(\varepsilon,\mathcal V_P)
 \le C\sqrt{P/\varepsilon}\log(CP/\varepsilon),
 \qquad 0<\varepsilon\le P. \tag{E}
\]
Here and below an extra universal logarithm may be inserted without
changing the conclusions. To see (E), partition the interval both by
a uniform grid with $k$ bins and by quantiles of the second-derivative
variation measure with at most $k$ internal cuts. The common
refinement has at most $2k+1$ intervals. On each open interval its
length is at most $2/k$ and its variation measure is at most $cP/k$;
atoms at endpoints do not affect linear interpolation. The linear
interpolant therefore differs uniformly from the original function
by at most $CP/k^2$. Choose $k=C\sqrt{P/\varepsilon}$.
Snap the interpolant's knots to an $\varepsilon/(CP)$ grid and
its values to an $\varepsilon/C$ grid. Its Lipschitz constant is
at most $cP$, so snapping and merging colliding knots changes the
uniform error by at most $C\varepsilon$. The number of resulting
descriptions proves (E). The affine part is included by the
supremum and Lipschitz bounds.

\paragraph{Two multiplier estimates.}
Fix a deterministic bounded comparator $g$, with $|g|\le1$ and
path variation $P_g$. For any shallow ReLU output $f$, write
$h=f-g$ and $P=1+P_f+P_g$. Both $h$ and the possible values of $f$
belong to a constant enlargement of $\mathcal V_P$.
For $0<\eta<\eta_0$, where $\eta_0>0$ depends only on the two density bounds, with probability at least $1-\delta$,
simultaneously over all $f$,
\begin{align}
 |P_n\xi h|
 &\le \eta P_nh^2+
 C\eta^{-3/5}P^{2/5}n^{-4/5}\mathcal L^C
 +C\eta^{-1}n^{-1}\mathcal L^C, \tag{M1}\\
 |(P_n-P_X)((f_0-g)h)|
 &\le \eta P_Xh^2+
 C\eta^{-3/5}P^{2/5}n^{-4/5}\mathcal L^C
 +CPn^{-13/8}\mathcal L^C\\
 &\quad+C\eta^{-1/2}P^{1/2}n^{-3/2}\mathcal L^C
 +C\eta^{-1}n^{-1}\mathcal L^C. \tag{M2}
\end{align}
Here $\mathcal L=1+\log(n/\delta)+\log(P+1)$; one can replace its
$P$ dependence on dyadic shells by the corresponding shell logarithm.
Both events depend only on the function classes, not on learner
width or Gaussian scale.

For (M1), condition on the inputs. The process $P_n\xi h$ is
sub-Gaussian for the empirical $L^2$ metric, with variance proxy
$C/n$. On $P_nh^2\le r^2$, Dudley's chaining bound and (E) give
\[
 \mathbb E_\xi\sup |P_n\xi h|
 \le C n^{-1/2}P^{1/4}r^{3/4}\mathcal L^C.
\]
The usual sub-Gaussian chaining tail adds
$Cr\sqrt{t/n}$. Peel in $P$ and $r$ and use Young's inequality.
Terminate the radius peeling at $r=1/n$. For the smaller terminal
class, $|P_n\xi h|\le2r\le2/n$ directly. For the design
multiplier use the localized estimate at $r=1/n$. Thus no infinite
peeling or uncontrolled small-radius limit is needed.

For (M2), localize first to $P_Xh^2\le r^2$. Apply the design-norm
lemma at a fixed tolerance, say $1/4$, and failure $n^{-10}\delta$.
It bounds all empirical radii by
\[
 C r+C P\{\log(Cn/\delta)/n\}^{3/2}.
\]
Conditional Rademacher chaining using (E), with the bounded
multiplier $f_0-g$, gives expected supremum at most
\[
 C n^{-1/2}P^{1/4}r^{3/4}\mathcal L^C
       +CPn^{-13/8}\mathcal L^C.
\]
The exceptional design event contributes at most $CPn^{-10}\delta$.
The exponent $13/8$ follows by inserting the additional empirical
radius $P n^{-3/2}$ in $n^{-1/2}P^{1/4}r^{3/4}$.

The lower density bound supplies a localized envelope:
\[
 \|h\|_\infty\le C\{r+P^{1/3}r^{2/3}\}.
\]
Indeed, if $|h(x_0)|=A$, a one-sided interval of length
$c\min(1,A/P)$ has $|h|\ge A/2$. Hence
$r^2\ge c A^2\min(1,A/P)$, proving the two alternatives.
Bousquet's inequality therefore adds
\[
 Cr\sqrt{t/n}+C(r+P^{1/3}r^{2/3})t/n,
\]
using variance at most $4r^2$. Its square-root term involving the
expected supremum is absorbed by that expectation and the displayed
envelope term. The population localization makes this envelope
deterministic on the whole class. Young's inequality bounds
$P^{1/3}r^{2/3}/n$ by
$\eta r^2+C\eta^{-1/2}P^{1/2}n^{-3/2}$.
Dyadic peeling and the same terminal radius $1/n$ prove (M2).
No independence of $f_0(X)-g(X)$ and $X$ is assumed.

\paragraph{A relative squared-loss inequality.}
Let
\[
 \widehat R(f)=P_n(f-Y)^2,\qquad
 D_n=C\eta^{-5}\{\log(Cn/\delta)/n\}^{3}.
\]
Combine (M1), (M2), and the two directions of the design-norm lemma:
\begin{equation}
 \begin{aligned}
 |(\widehat R(f)-\widehat R(g))-(A(f)-A(g))|
 \le{}& C\eta\{A(f)+A(g)\}
 +C\eta^{-3/5}P^{2/5}n^{-4/5}\mathcal L^C\\
 &+CPn^{-13/8}\mathcal L^C
 +C\eta^{-1/2}P^{1/2}n^{-3/2}\mathcal L^C\\
 &+C\eta^{-1}n^{-1}\mathcal L^C+D_nP^2.
 \end{aligned}\tag{L}
\end{equation}
Indeed the empirical/population quadratic-norm difference for
$h$ is at most $C\eta P_Xh^2+D_nP^2$, and
$P_Xh^2\le2A(f)+2A(g)$. The multiplier (M1) initially involves
$P_nh^2$, which the reverse norm inequality controls by the same
quantity. Constants in the fine-tolerance design lemma can be
absorbed by replacing its input tolerance by a fixed multiple of
$\eta$.

\paragraph{The actual Gaussian posterior.}
Use the original shallow independent Gaussian prior
\[
 a_j,b_j\sim N(0,\nu),\quad v_j\sim N(0,\nu/N^2),
 \quad \omega=N/\nu,\quad 0<\nu\le1,
\]
and define
\[
 Q=N^{-1}\sum(a_j^2+b_j^2)+N\sum v_j^2,\quad
 \Pi_\lambda(dw)\propto
 e^{-\omega\widehat R(z_w)/(2\lambda)}\mu(dw).
\]
Then $P(z_w)\le Q/2$, $\omega Q\sim\chi^2_{3N}$, and the physical
Gaussian observation variance is $n\lambda\nu/N$.
The Gibbs identity and the Gaussian exponential moment imply
\[
 \Pi_\lambda[\widehat R(z)+(\lambda/2)Q]
 \le F_\lambda+3\lambda\log2,\qquad
 F_\lambda=-\frac{2\lambda}{\omega}\log Z_\lambda. \tag{G}
\]
In fact
$\operatorname{KL}(\Pi_\lambda\Vert\mu)
 \ge\omega\Pi_\lambda Q/4-(3N/2)\log2$ and $N/\omega\le1$.

For $0<\lambda\le1$ there is also the deterministic bound
\[
 \Pi_\lambda Q^2\le C(1+\lambda^{-2}). \tag{T}
\]
The prior empirical risk is at most two for every unit-bounded
sample, so Jensen gives $F_\lambda\le2$. A chi-square Chernoff
bound gives $\mu(Q>u)\le e^{-c\omega u}$ above a fixed universal
threshold. Since $Z_\lambda\ge e^{-\omega/\lambda}$,
the posterior tail is at most
$\min(1,e^{\omega/\lambda-c\omega u})$; integrate it and use
$\omega\ge1$ to obtain (T).

For the following uniform statistical bound assume $P_g\le n^4$.
Truncate posterior energy at $Q\le n^{20}/\lambda$; the tail estimate
used in (T) bounds every omitted risk or quadratic-energy integral
by $O(n^{-10})$. On the retained range all shell logarithms are
$O(\log(n/\delta))$ when $n^{-2}\le\lambda\le1$.
For any $P\ge1$ and $a>0$,
\[
 C\eta^{-3/5}P^{2/5}n^{-4/5}\mathcal L^C
 \le a\lambda P+
 C_a\eta^{-1}n^{-4/3}\lambda^{-2/3}\mathcal L^C. \tag{Y}
\]
For the additional sublinear term in (M2), Young gives
\[
 C\eta^{-1/2}P^{1/2}n^{-3/2}\mathcal L^C
 \le a\lambda P+C_a\eta^{-1}\lambda^{-1}n^{-3}\mathcal L^C.
\]
The linear term $Pn^{-13/8}$ is absorbed provided
\[
 \lambda\ge C n^{-13/8}\log^k(n/\delta)
\]
for one sufficiently large fixed $k$. The Gaussian tail truncation
above makes every retained shell logarithm $O(\log(n/\delta))$;
the discarded posterior risk is negligible uniformly in $N,\nu$.

A balanced Gaussian comparison around $g$ gives, for every
$N$ above the width of $g$,
\[
 F_\lambda\le\widehat R(g)
       +C/n+\lambda C[P_g+\log(n(P_g+1))].
\]
Choose $a$ in (Y) small enough that the resulting coefficient
of $Q$ in (L) is at most $\lambda/2$, and integrate (L).
Using (G), (T), and the comparison gives
\begin{align}
 \Pi_\lambda A(z)
 \le &(1+C\eta)A(g)
 +C\lambda[P_g+\log(n(P_g+1))]
 +C\eta^{-1}n^{-4/3}\lambda^{-2/3}\mathcal L_n^C\\
 &+C\eta^{-1}\lambda^{-1}n^{-3}\mathcal L_n^C
 +C\eta^{-5}n^{-3}\mathcal L_n^C
       [\,\lambda^{-2}+(P_g+1)^2\,]+C\eta^{-1}n^{-1}\mathcal L_n^C, \tag{O}
\end{align}
where $\mathcal L_n$ contains only logarithms in
$n,1/\delta,P_g+1$. Jensen gives the same bound for the ordinary
raw posterior mean. All events are shared by every eligible
finite width and every $0<\nu\le1$.

\paragraph{A universal comparison-width guarantee.}
Take $\lambda=n^{-3/2}\log^K(n/\delta)$ with one sufficiently large
fixed $K$, and put $\eta=1/\log n$. Restrict to the nontrivial
comparison widths $(C+1)\sqrt\lambda\le1$.
The bounded spline compiler applied to a clipped width-$C$ univariate
shallow ReLU function has width at most $6C+6$, risk increment
$O(Ct)$ and path variation $P_g\le C'(C+1)/t$.
Use $t=\sqrt\lambda$ in (O). The compiler and comparison penalty
are $\widetilde O((C+1)n^{-3/4})$. The noise remainder is
$\widetilde O(n^{-1/3})$, and the new envelope remainder is
$\widetilde O(n^{-3/2})$.
The posterior design term is
\[
 C\eta^{-5}n^{-3}\mathcal L_n^c\lambda^{-2}
 \le C\log^{c+5-2K}(n/\delta)\le C/\log n
\]
after increasing $K$. Its comparator-energy term is this quantity
times $O((C+1)^2\lambda)$, and is therefore no larger in the
stated range. Thus
\[
 R(\bar z_\lambda)\le\Pi_\lambda R(z)
 \le R_0+\{1+O(1/\log n)\}
       \inf_{g\in\mathcal G_C}\{R(g)-R_0\}
 +\widetilde O((C+1)n^{-3/4}+n^{-1/3})+O(1/\log n).
 \tag{C}
\]
Clipping a proof comparator to $[-1,1]$ cannot increase its squared
risk for unit-bounded labels, and creates at most $3C+2$ breakpoints.
The actual posterior remains raw and unchanged.
The same deterministic schedule works for all qualifying $C$.
A summable confidence allocation over integer comparison widths is
a proof union, not a prior over widths. In particular it permits
every fixed growth $C_n=n^u$ with $u<3/4$.
The logarithmic additive term is explicit: this universal schedule
is not the polynomial-accuracy bound of the declared-budget schedule
below. The conditions on $\eta,\lambda,P_g$ hold for sufficiently
large $n$; a bounded-risk fallback handles smaller samples.

\paragraph{Two coefficient-one schedules.}
The sharper design lemma yields two genuine coefficient-one
consequences, not only the coefficient-$1+o(1)$ interpretation.
All the following displays suppress fixed powers of logarithms;
these can be realized by enlarging the indicated tolerances and
likelihood scales by fixed logarithmic factors.

First take the universal schedule
\[
 \lambda=n^{-3/5},\qquad \eta=c n^{-3/10},
 \qquad t=\sqrt\lambda=n^{-3/10}.
\]
The likelihood schedule does not use $C$. Since $A(g)\le4$ for the
bounded proof comparator, its multiplicative slack is absorbed
additively. Substituting in (O) gives, for every $1\le C\le n$,
\[
 R(\bar z_\lambda)\le\inf_{g\in\mathcal G_C}R(g)
 +\widetilde O\left((C+1)n^{-3/10}
                   +(C+1)^2n^{-9/10}\right). \tag{U}
\]
The multiplier remainder has order $n^{-19/30}$ and the direct
$n^{-1}\eta^{-1}$ remainder has order $n^{-7/10}$; both are smaller
than the first displayed term. The design remainder from the
posterior energy is $n^{-3/10}$, and its comparator-energy term is
$(C+1)^2n^{-9/10}$. The hypotheses $\lambda\ge\widetilde C n^{-13/8}$ and
$n\eta^2\ge C\log(Cn/\delta)$ hold for sufficiently large $n$.
The constants depend on the two density bounds.

Second, if a comparison-width budget $C$ is declared, put
\[
 \varepsilon_{n,C}
 =\widetilde C(C+1)^{2/5}n^{-3/10},\qquad
 \eta=c\varepsilon_{n,C},\quad
 \lambda=c'\left(\frac{\varepsilon_{n,C}}{C+1}\right)^2,\quad
 t=c''\frac{\varepsilon_{n,C}}{C+1}.
\]
Choose the fixed constants and logarithmic powers in this order
so that the inequalities used above hold. For the nontrivial regime
$\varepsilon_{n,C}\le\varepsilon_0$, this yields
\[
 R(\bar z_\lambda)\le\inf_{g\in\mathcal G_C}R(g)
 +\widetilde O\left((C+1)^{2/5}n^{-3/10}\right). \tag{B}
\]
The learner uses $n,C$ and the confidence allocation to choose
one likelihood variance; it does not mix widths, priors or
temperatures. It still uses the same prior and any prescribed
finite learner width $N\ge6C+6$.

For the algebra, the compiler and penalty terms are $O(\varepsilon)$.
The strongest design requirement is
\[
 (C+1)^4 n^{-3}\varepsilon^{-9}\le C\varepsilon,
\]
which is exactly $\varepsilon\gtrsim(C+1)^{2/5}n^{-3/10}$.
The new envelope remainder requires only
$\varepsilon\gtrsim(C+1)^{1/2}n^{-3/4}$ and is weaker.
The noise multiplier asks only for
$\varepsilon\gtrsim(C+1)^{2/5}n^{-2/5}$, also weaker.
The comparator-energy design term is smaller than the posterior
one by a factor $O(\varepsilon^2)$.
In the nontrivial regime $\varepsilon_{n,C}\le\varepsilon_0$,
one has $C+1\lesssim n^{3/4}$ and hence
$\lambda\gtrsim n^{-3/2}$ with logarithmic slack. This exceeds
$n^{-13/8}\log^k(n/\delta)$ for sufficiently large $n$.
The other conditions $P_g\le n^4$ and
$n\eta^2\ge C\log(Cn/\delta)$ also hold.
Thus a sufficient sample size for excess $\varepsilon$ is
\[
 n\ge\widetilde O\left(
 (C+1)^{4/3}\varepsilon^{-10/3}\right). \tag{SC}
\]
Consequently every prescribed comparison width $C_n=n^u$, $u<3/4$,
has a vanishing power-rate remainder, at the deterministic scale
$\lambda_n=\widetilde\Theta(n^{-3/5-6u/5})$.

\section{A sharper oracle when learner width is at least the sample size}
\label{sec:empirical-midpoint}
Averaging a posterior function with a comparison function improves
empirical squared loss by one quarter of their squared disagreement.
A linear program realizes those averaged training values using at
most $n$ neurons and controlled path variation. The Gaussian norm
cost of a neighborhood of that implementation is small enough that
posterior optimality forces an empirical-distance inequality.
The one-dimensional variation estimates transfer this inequality to
test risk with comparison coefficient exactly one. The linear program
is a proof device: the actual independent Gaussian prior, likelihood,
and raw posterior prediction are unchanged.

Use the same shallow model and notation:
\[
 Q_N=N^{-1}\sum(a_j^2+b_j^2)+N\sum v_j^2,\quad
 \omega=N/\nu,\quad P(w)=\sum_j|v_j|\sqrt{a_j^2+b_j^2}\le Q_N/2.
\]
The actual posterior is
$d\Pi_\lambda\propto
 \exp[-\omega\widehat R(z)/(2\lambda)]\,d\mu$.
Labels are initially bounded by one and the input density is bounded
above and below by positive constants.

The finite-sample atomic representation is the standard conic
Carath\'eodory principle for variation-norm neural networks; see
Bach~\cite{bach-convex}, Section 2.2. The following proof also specifies
a measurable choice for its use inside the posterior comparison.

\begin{lemma}[Minimum-mass compression of a training vector]
For every finite sample of $n$ inputs and any two shallow ReLU
functions $f,g$ of path variations $P_f,P_g$, the training vector
$(f(X_i)+g(X_i))_{i=1}^n/2$ has a shallow ReLU implementation
using at most $n$ neurons and path variation at most $(P_f+P_g)/2$.
The implementation can be chosen measurably in the original
parameters and the sample inputs.
\end{lemma}
\begin{proof}
Normalize every nonzero incoming row of $f$ and $g$ to unit norm,
and absorb the sign of its readout coefficient into the signed
unit-ridge feature.
Let the columns of $A$ be the resulting finitely many training
feature vectors, and let $y=(f(X_i)+g(X_i))_{i=1}^n/2$.
The linear program
\[
 \min\left\{\sum_j\alpha_j:A\alpha=y,\ \alpha_j\ge0\right\}
\]
is feasible with cost at most $(P_f+P_g)/2$, using the original
absolute amplitudes divided by two. Its bounded sublevel sets are
compact, so an optimum exists. An optimum with minimal support has
linearly independent supported columns: otherwise move along a
column dependence, choosing the sign that does not increase cost,
until a coefficient vanishes. If the dependence changed cost
strictly it would contradict optimality; if not it would contradict
minimal support. Its support size is therefore at most
$\operatorname{rank}A\le n$.

For measurability, enumerate the finitely many independent column
subsets of size at most $n$. For each, solve the coefficient vector
by its Gram inverse, retaining it only when it is nonnegative and
satisfies the training equations. Choose minimum cost, breaking ties
lexicographically. All tests and formulas are Borel on their domains.
The zero training vector is implemented by the zero network.
\end{proof}

\begin{lemma}[An empirical midpoint comparison law]
\label{lem:empirical-midpoint-law}
Fix a deterministic comparator $g$ of width $k$, $|g|\le1$ and
path variation $P_g\le n^4$. Suppose $N\ge n$ and
$n^{-2}\le\lambda\le1$. There is a law $\rho$ on the same
width-$N$ Gaussian parameter space such that
\begin{align}
 \rho\widehat R(z)
 &\le\tfrac12\Pi_\lambda\widehat R(z)
       +\tfrac12\widehat R(g)
       -\tfrac14\Pi_\lambda P_n(z-g)^2+Cn^{-10}, \tag{EM1}\\
 \operatorname{KL}(\rho\Vert\mu)
       -\operatorname{KL}(\Pi_\lambda\Vert\mu)
 &\le-\tfrac14\omega\Pi_\lambda P
       +\tfrac12\omega P_g+CN\log(n+1). \tag{EM2}
\end{align}
\end{lemma}
\begin{proof}
For each posterior draw $f=z_w$, apply the preceding minimum-mass
compression directly to the midpoint training vector of $f$ and $g$,
and pad its at most $n$ neurons with zeros. Its path variation is
at most $(P_f+P_g)/2$.
Balance each neuron using positive homogeneity: a unit direction
with signed amplitude $c$ is represented by incoming row
$\sqrt{N|c|}\,\theta$ and readout
$\operatorname{sign}(c)\sqrt{|c|/N}$.
Thus the parameter center $\bar w_f$ obeys
\[
 Q_N(\bar w_f)\le P_f+P_g.
\]
At every training input its function is exactly $(f+g)/2$.

Around this center place a Gaussian law with every native
coordinate variance multiplied by $\epsilon^2$, where
$\epsilon=n^{-30}$, and mix these laws using the posterior
distribution of $f$. A single component has relative entropy
\[
 \operatorname{KL}(\rho_f\Vert\mu)
 \le\frac{\omega}{2}(P_f+P_g)+3N\log(1/\epsilon).
\]
Convexity of relative entropy gives the same upper bound after
averaging over $f$.

The Gaussian energy exponential moment is
\[
 \log\mathbb E_\mu e^{3\omega Q_N/8}=3N\log2.
\]
The entropy variational inequality and $Q_N\ge2P$ therefore give
\[
 \operatorname{KL}(\Pi_\lambda\Vert\mu)
 \ge\tfrac34\omega\Pi_\lambda P-3N\log2.
\]
Subtracting proves (EM2).

For clarity, the balanced Gaussian perturbation has uniform
root-mean-square error at most
\[
 C\epsilon\sqrt\nu
 \left(\sum_j|\bar v_j|+N^{-1}\sum_j\|\bar r_j\|_2\right)
 +C\epsilon^2\nu
 \le C\epsilon\sqrt{P_f+P_g+1}+C\epsilon^2.
\]
Indeed balance and support at most $N$ bound each displayed sum
by the square root of the center path variation. The last term
is the product of independent incoming and readout perturbations.
At training inputs the center is exactly $(f+g)/2$.
Since $|g|,|Y_i|\le1$ and $|f|\le\sqrt2P_f$, the expected
squared-loss perturbation is bounded by
\[
 C\epsilon\,\Pi_\lambda
       [(P_f+1)\sqrt{P_f+P_g+1}]
       +C\epsilon^2\Pi_\lambda(P_f+P_g+1).
\]
The deterministic Gaussian posterior bound
$\Pi_\lambda P_f^2\le C/\lambda^2$, together with
$\lambda\ge n^{-2}$ and $P_g\le n^4$, makes this at most
$Cn^{-10}$. The squared-loss midpoint identity now gives (EM1).
\end{proof}

\begin{lemma}[A strong empirical comparison inequality]
\label{lem:empirical-midpoint-objective}
Under the preceding assumptions,
\[
 \Pi_\lambda\widehat R(z)-\widehat R(g)
 +\tfrac12\Pi_\lambda P_n(z-g)^2+\lambda\Pi_\lambda P
 \le2\lambda P_g+C\lambda\log(n+1)+Cn^{-10}. \tag{EM3}
\]
\end{lemma}
\begin{proof}
The actual posterior minimizes
$\rho\widehat R+(2\lambda/\omega)\operatorname{KL}(\rho\Vert\mu)$.
Compare it to the preceding mixture, use (EM1)--(EM2), and
multiply the result by two. Since $N/\omega=\nu\le1$, the entropy
remainder is at most $C\lambda\log(n+1)$.
The comparison law is only a proof device; the learner itself
uses one fixed architecture, prior and likelihood.
\end{proof}

\begin{theorem}[A sharp raw path oracle at every width at least $n$]
\label{thm:wide-sharp-path}
For every deterministic bounded comparator $g$ as above, every
$N\ge n$, every $0<\nu\le1$, and
\[
 \widetilde Cn^{-13/8}\le\lambda\le1,
\]
with probability at least $1-\delta$,
\begin{align}
 R(\bar z_\lambda)\le\Pi_\lambda R(z)
 \le R(g)+\widetilde O\bigl(
 &\lambda(P_g+1)+n^{-4/3}\lambda^{-2/3}
       +n^{-3}\lambda^{-2}\\
 &+n^{-3}(P_g+1)^2+n^{-1}\bigr). \tag{EM4}
\end{align}
The event is common to all the eligible widths and Gaussian scales.
\end{theorem}
\begin{proof}
Take one sufficiently small fixed tolerance in the localized noise
and bounded design-multiplier inequalities of
Section~\ref{sec:localized-proof}, including its localized envelope. Keep their quadratic
terms as norms of $h=z-g$ instead of bounding them by approximation
risks. The sharp design-norm comparison, in both directions,
then gives uniformly over $z$
\begin{align*}
 R(z)-R(g)\le{}&
 \widehat R(z)-\widehat R(g)
 +\tfrac14P_n(z-g)^2+\tfrac12\lambda P(z)
 +C\lambda(P_g+1)\\
 &+\widetilde Cn^{-4/3}\lambda^{-2/3}
 +\widetilde Cn^{-3}(P(z)+P_g+1)^2
 +\widetilde Cn^{-3}\lambda^{-1}+\widetilde C/n.
\end{align*}
Young's inequality absorbs the localized $P^{2/5}$ and $P^{1/2}$
terms into the displayed path penalty. The remaining linear
$Pn^{-13/8}$ term is absorbed by the assumed likelihood floor.
The same finite-shell Gaussian-tail argument controls the logarithms.
Integrate, apply (EM3), and discard the leftover nonnegative
empirical-distance and path-energy terms. Finally use
$\Pi_\lambda P^2\le C/\lambda^2$ and
$\lambda^{-1}\le\lambda^{-2}$. Jensen proves the statement for
the ordinary raw mean. The fixed tolerance is why the leading
comparison coefficient is exactly one.
\end{proof}

\begin{corollary}[Coefficient-free wide-network rates]
\label{cor:wide-sharp-width}
Let $1\le C\le n$ and let the comparison class contain every
univariate shallow ReLU network of width at most $C$, with arbitrary
finite coefficients. At every prescribed
$N\ge n$, the deterministic choice
\[
 \lambda=\widetilde\Theta\left(
 \max\{(C+1)^{-6/7}n^{-8/7},
       (C+1)^{-2/5}n^{-6/5}\}\right)
\]
gives
\[
 R(\bar z_\lambda)\le\Pi_\lambda R(z)
 \le\inf_{\text{width }C}R+
 \widetilde O\left(
 (C+1)^{4/7}n^{-4/7}
 +(C+1)^{4/5}n^{-3/5}\right). \tag{EM5}
\]
The likelihood schedule uses the declared comparison budget but not
its coefficients or risk. The universal schedule
$\lambda=\widetilde\Theta(n^{-8/7})$ instead gives
\[
 R(\bar z_\lambda)\le\inf_{\text{width }C}R+
       \widetilde O((C+1)n^{-4/7}) \tag{EM6}
\]
without using $C$ in the schedule.
A sufficient sample size for excess $\varepsilon$ in (EM5) is
\[
 n\ge\widetilde O\left(
 (C+1)\varepsilon^{-7/4}
 +(C+1)^{4/3}\varepsilon^{-5/3}\right). \tag{EM7}
\]
\end{corollary}
\begin{proof}
Use the bounded spline compiler at $t=\sqrt\lambda$.
It has width at most $6C+6$, bounded output, path variation
at most $c(C+1)/t$, and comparison-risk increment
$O((C+1)t)$. Substitution in (EM4) gives
\[
 \inf_{\text{width }C}R+\widetilde O\left(
 (C+1)\sqrt\lambda+n^{-4/3}\lambda^{-2/3}
 +n^{-3}\lambda^{-2}
 +\frac{(C+1)^2}{n^3\lambda}+n^{-1}\right).
\]
Balance the first term against the second and third and take
the larger of the resulting scales. In the nontrivial regime,
$(C+1)\sqrt\lambda=O(1)$, so the comparator-design term is dominated
by $n^{-3}\lambda^{-2}$. Outside it, use the zero comparator at
the same scale to get a bound absorbed by the displayed constant-size
remainder. The scales are above the $n^{-13/8}$ floor for $C\le n$.
The two balances yield (EM5). For the fixed choice in (EM6),
the noise term is $n^{-4/7}$, the posterior design term is
$n^{-5/7}$, and the comparator design term
$(C+1)^2n^{-13/7}$ is dominated by $(C+1)n^{-4/7}$ for $C\le n$.
Inverting its two optimized remainder terms gives (EM7).
All comparisons can share one event by a summable allocation over
the integer width budgets. No architecture mixture is involved.
\end{proof}


\paragraph{One fixed schedule for a prescribed polynomial width range.}
More generally, fix $0<a<3/2$ and use
$\lambda=\widetilde\Theta(n^{-a})$, independently of the comparison
width. The same proof gives, in the nontrivial range
$(C+1)n^{-a/2}\lesssim1$,
\[
 R(\bar z_\lambda)\le\inf_{\text{width }C}R+
 \widetilde O\left((C+1)n^{-a/2}
       +n^{-4/3+2a/3}+n^{-3+2a}\right).
\]
Thus every fixed width-growth exponent $u<3/4$ is compatible with
one deterministic schedule having a vanishing power-rate remainder:
choose $2u<a<3/2$. This coefficient-one statement concerns the
regime $N\ge n$; it does not replace the separate smaller-learner-width
result.

\begin{corollary}[Marginally sub-Gaussian labels in the wide regime]
\label{cor:wide-sharp-subgaussian}
All conclusions of Corollary~\ref{cor:wide-sharp-width} extend to
$\E\exp(Y^2/s_Y^2)\le2$, with the same powers of $C,n$ and additional
logarithmic factors. Every risk still uses the actual raw output and
original labels. If $\lambda'$ is the corresponding displayed schedule,
the physical scale is $\lambda=M_n^2\lambda'$, with
$M_n=\max\{1,s_Y\sqrt{c\log(n/\delta)}\}$; its observation variance
is $n\lambda\nu/N$. The high-probability event is common to all
$N\ge n$ and all $0<\nu\le1$.
\end{corollary}
\begin{proof}
Use the proof-only clipped label $Y_M$ on the unchanged input law,
as in Corollary~\ref{cor:localized-subgaussian}. Analyze observable
$z/M_n$ and label $Y_M/M_n$ with the original parameter prior.
The minimum-mass linear program compresses the training midpoint
of the physical $z$ and physical comparator $M_nq$ into at most
$n$ neurons, exactly as before. Its path bound becomes
$(P(z)+M_nP(q))/2$. Consequently the midpoint entropy calculation
and its positive empirical-distance term are unchanged, apart from
the logarithmic factor $M_n$ in the comparison resource.
More explicitly, the normalized midpoint objective inequality has
left side
\[
 \Pi\{\widehat R(z/M_n)-\widehat R(q)\}
 +\tfrac12\Pi P_n(z/M_n-q)^2+\lambda'\Pi P(z)
\]
and right side $2\lambda'M_nP(q)+C\lambda'\log(n+1)+Cn^{-10}$.
The normalized path $P(z)/M_n$ is at most $P(z)$, so all localized
process errors still absorb into this physical path penalty.
The prior energy and normalized empirical-risk bounds still give
$\Pi Q_N^2\le C(1+(\lambda')^{-2})$.

All statistical and optimization steps therefore retain their
powers after enlarging the logarithmic factors. Multiplying normalized
risks by $M_n^2$ preserves their leading coefficient one. On the
all-training-labels event the proof posterior equals the actual
posterior. The raw population-risk difference is bounded by
\[
 C(1+(\lambda')^{-1})\{\E(Y-Y_M)^2\}^{1/2}
       +\E|Y^2-Y_M^2|,
\]
which is negligible for the stated polynomial schedules and a
sufficiently large $c$. Finally clip each original comparator only
inside the proof, using the projection inequality and risk transfer
of Corollary~\ref{cor:localized-subgaussian}. This does not change
its number of affine pieces beyond the stated $3C+2$ bound.
A union of the training-label and statistical events finishes.
\end{proof}

\section{A schedule-specific width obstruction on the uniform interval}
\label{sec:univariate-schedule-sharpness}

The following lower bound concerns the likelihood schedules used by the
one-dimensional upper bounds. It holds on one fixed uniform input law
and uses noiseless targets. It is not an obstruction to every possible
temperature choice.

Use one hidden ReLU layer and the ordinary raw output
\[
 z_w(x)=\sum_{j=1}^N v_j(a_jx+b_j)_+,
 \qquad x\in[-1,1].
\]
For $0<\nu\le1$, the independent prior variances are $\nu$ in
$a_j,b_j$ and $\nu/N^2$ in $v_j$. Put
\[
 Q_N(w)=\frac1N\sum_{j=1}^N(a_j^2+b_j^2)
                   +N\sum_{j=1}^Nv_j^2,
 \qquad \omega=\frac N\nu.
\]
For any sample $S=((X_i,Y_i))_{i=1}^n$ and $\lambda>0$, the actual
posterior is
\begin{equation}\label{eq:univariate-sharpness-posterior}
 d\Pi_S(w)\propto
 \exp\!\left[-\frac\omega2
       \left(Q_N(w)+\frac{\widehat R_S(z_w)}\lambda\right)\right]dw,
 \qquad
 \widehat R_S(z)=\frac1n\sum_{i=1}^n(z(X_i)-Y_i)^2.
\end{equation}
Its physical Gaussian observation variance is $n\lambda\nu/N$.

\begin{lemma}[Posterior quadratic energy]
\label{lem:univariate-posterior-energy}
For every finite sample with finite labels, every $N\ge1$,
$0<\nu\le1$, and $\lambda>0$,
\begin{align}
 \int Q_N\,d\Pi_S
 &=3\nu-\frac2\lambda\int
     \frac1n\sum_{i=1}^n z_w(X_i)(z_w(X_i)-Y_i)\,d\Pi_S(w)
       \label{eq:univariate-energy-identity}\\
 &\le 3\nu+\frac{1}{2\lambda n}\sum_{i=1}^nY_i^2.
       \label{eq:univariate-energy-bound}
\end{align}
\end{lemma}
\begin{proof}
The parameter dimension is $D=3N$. For positive scalars $r$,
$Q_N(rw)=r^2Q_N(w)$ and $z_{rw}=r^2z_w$.
Integrate the derivative of $r^D$ times the unnormalized posterior
density on each ray, then integrate over directions. The endpoint
terms vanish at both zero and infinity. The resulting identity is
\[
 0=D-\omega\int Q_N\,d\Pi_S
   -\frac{2\omega}{\lambda}\int
        \frac1n\sum_i z_w(X_i)(z_w(X_i)-Y_i)\,d\Pi_S(w).
\]
This argument does not differentiate ReLU across a gate boundary:
its positive radial homogeneity is exact. For each fixed $N,\nu$
and sample, the positive definite Gaussian quadratic form dominates
all polynomial factors arising in the differentiation; the
likelihood is at most one. This justifies the integrations and also
shows that every posterior moment used here is finite. Since
$D/\omega=3\nu$ and $t(t-y)\ge-y^2/4$, the two claims follow.
\end{proof}

Define the continuous one-period triangular function
\[
 \varphi(t)=
 \begin{cases}
  1-4t,&0\le t\le1/2,\\
  4t-3,&1/2\le t\le1,
 \end{cases}
\]
and extend it periodically. For an integer $m\ge1$, let
\[
 f_m(x)=\varphi\!\left(\frac{m(x+1)}2\right),\qquad -1\le x\le1.
\]
Thus $f_m$ has $m$ full periods, takes values in $[-1,1]$, and is
the conditional mean of the noiseless law $X\sim\operatorname{Unif}[-1,1]$,
$Y=f_m(X)$. Write $P$ for integration against this uniform input law.

\begin{lemma}[A triangular target and its double primitive]
\label{lem:univariate-triangular-target}
The target satisfies $P f_m^2=1/3$ and has an exact raw shallow ReLU
representation with width $2m+1$ and coefficient cap $4m$.
There is a twice continuously differentiable function $F_m$ on
$[-1,1]$ such that
\[
 F_m''=f_m/2,\qquad F_m(\pm1)=F_m'(\pm1)=0,
 \qquad \|F_m\|_\infty=\frac{1}{12m^2}.
\]
Consequently every network $z_w$ above obeys
\begin{equation}\label{eq:univariate-triangle-correlation}
 |P(z_wf_m)|\le\frac{Q_N(w)}{24m^2}.
\end{equation}
\end{lemma}
\begin{proof}
Direct integration gives
\[
 \int_0^1\varphi(t)\,dt=\int_0^1t\varphi(t)\,dt=0,
 \qquad \int_0^1\varphi(t)^2\,dt=\frac13.
\]
For $F_*(t)=\int_0^t(t-u)\varphi(u)\,du$, both $F_*$ and
$F_*'$ vanish at $0,1$. On $[0,1/2]$,
$F_*(t)=t^2/2-2t^3/3$; reflection gives its values on the other half.
It is nonnegative and has maximum $F_*(1/2)=1/24$.
On each period of length $2/m$, set $F_m$ equal to the rescaled
copy $(2/m^2)F_*$. The copies and their first two derivatives match
at period endpoints, giving all asserted properties.

For $x_j=-1+j/m$, the exact representation on $[-1,1]$ is
\[
 f_m(x)=1-2m(x+1)_+
           +4m\sum_{j=1}^{2m-1}(-1)^{j+1}(x-x_j)_+.
\]
The constant one is a neuron with input weight zero and bias one.
All hidden input weights and biases have magnitude at most one,
and all readout coefficients have magnitude at most $4m$.

The distributional second derivative of $z_w$ in the open interval
is a finite signed measure with
\[
 \operatorname{TV}(z_w'')
 \le\sum_j|v_j a_j|
 \le\frac12\sum_j\left(Nv_j^2+\frac{a_j^2}{N}\right)
 \le\frac12Q_N(w).
\]
Terms with $a_j=0$ have no second derivative. Two integrations by
parts, with the displayed boundary values of $F_m$, give
$P(z_wf_m)=\int F_m\,d z_w''$. Multiplying the total variation
bound by $\|F_m\|_\infty$ proves
\eqref{eq:univariate-triangle-correlation}.
\end{proof}

\begin{theorem}[Raw prediction lower bound for a chosen schedule]
\label{thm:univariate-schedule-obstruction}
Let $m,N\ge1$, $0<\nu\le1$, and $\lambda>0$. For every sample
with $X_i\in[-1,1]$ and $|Y_i|\le1$, both the ordinary posterior
mean $\bar z_S=\int z_w\,d\Pi_S(w)$ and the posterior-averaged
squared risk against the noiseless triangular target satisfy
\begin{align}
 P(\bar z_S-f_m)^2
 &\ge\frac13-\frac{3\nu+(2\lambda n)^{-1}\sum_iY_i^2}{12m^2}
 \ge\frac13-\frac{3+(2\lambda)^{-1}}{12m^2},
       \label{eq:univariate-schedule-lower}\\
 \int P(z_w-f_m)^2\,d\Pi_S(w)
 &\ge\frac13-\frac{3+(2\lambda)^{-1}}{12m^2}.
       \label{eq:univariate-schedule-gibbs-lower}
\end{align}
In particular these bounds hold surely, and in expectation, for
IID noiseless training data from the uniform-input target law.
The best raw teacher of width $2m+1$ has risk zero. The bounds are
uniform over every $N\ge2m+1$ and every mean-field-or-stronger scale.
\end{theorem}
\begin{proof}
Discard the nonnegative $P\bar z_S^2$ term and apply Fubini and
\eqref{eq:univariate-triangle-correlation}:
\[
 P(\bar z_S-f_m)^2
 \ge\frac13-2P(\bar z_S f_m)
 \ge\frac13-\frac1{12m^2}\int Q_N\,d\Pi_S.
\]
Lemma~\ref{lem:univariate-posterior-energy} proves the mean bound.
The identical argument before taking the posterior expectation
proves the Gibbs-risk bound. All moments are finite, so the
interchanges are justified. The explicit teacher representation
proves the zero-risk comparison claim.
\end{proof}

\paragraph{Consequences for three universal schedules.}
For a deterministic schedule $\lambda_n>0$, if $m_n\to\infty$
and $\lambda_nm_n^2\to\infty$, the lower bound approaches $1/3$.
More quantitatively, if $\lambda_n\to0$ and
$m_n=\lceil\lambda_n^{-1/2}\rceil$, its limit inferior is at least
$7/24$. Thus:
\begin{itemize}
\item For $\lambda_n=L_n/n^{3/2}$, where $1\le L_n\le
c\log^a(n+2)$ for fixed constants $a,c>0$, a teacher of width
$O(n^{3/4}/\sqrt{L_n})$ already has nonvanishing risk for this learner.
This matches the nearly $n^{3/4}$ comparison-width threshold of
the corresponding localized upper bound, up to logarithmic factors.
\item For $\lambda_n=n^{-3/5}$, a teacher of width $O(n^{3/10})$
has nonvanishing risk. This matches the width exponent of the
corresponding coefficient-one universal schedule.
\item For $\lambda_n=n^{-8/7}$, a teacher of width $O(n^{4/7})$
has nonvanishing risk. This matches the width exponent of the
coefficient-one empirical-midpoint oracle at $N\ge n$.
Fixed logarithmic factors in $\lambda_n$ change this threshold
only by the corresponding inverse square-root logarithmic factor.
\end{itemize}
The input law is the same uniform law throughout; the noiseless
targets and their coefficient magnitudes grow with $n$. No prior
is changed. These conclusions concern the specified schedules.
The lower bound weakens when $\lambda$ is made smaller and does
not exclude a faster theorem at another temperature or with
teacher-dependent calibration.


\section{Interpretation and limitations}
The literal polynomial-width oracle is a qualified version of the
approximation-learning conjecture: a vanishing additive term is present,
the input and comparison depth are fixed, and the learning rule has a
specific observable likelihood calibration. Every training observation
is used. There is no architecture mixture, output saturation, prior
conditioning, or infinitely wide limit.

On Boolean inputs the required learner depth is a fixed constant times
the comparison depth. On the continuous domains considered here two
hidden tanh layers suffice for compilation, although the width exponent
then deteriorates with the original comparison depth and input
dimension. Additional fixed learner depth is allowed. For the general fixed-depth compiler theorems, the assumption
$N\ge n$ is a simple sufficient presentation; the proof only needs
$N$ above the explicit compiled width.

The raw-output restriction matters. It is not enough to apply a theorem
for clipped predictions to an unbounded predictor. The compilers here
have bounded raw outputs themselves, and the statistical theorem
controls the actual raw posterior. Their artificial box indicators and
Boolean circuits occur only in a comparison argument; the Gaussian
learner's architecture and activation stay unchanged.

The comparison is statistical, not computational. Exact Bayesian
evidence and posterior means may be infeasible to compute. The
sufficient exponents are small and the depth-dependent constants may
be large. The results do not give polynomial sample complexity jointly
in unrestricted comparison depth and width, nor a claim for every
fixed observation variance.

\paragraph{Related work.}
Finite atomic representations under neural variation norms are standard;
Bach~\cite{bach-convex}, Section 2.2, gives the at-most-$n$ conic
Carath\'eodory principle on $n$ observed inputs. We use that principle
inside a Gaussian posterior comparison and supply the finite measurable
linear-programming selection. Tinsi and Dalalyan~\cite{tinsi-dalalyan}
analyze shallow networks with product Gaussian priors through
PAC--Bayesian aggregation inequalities and approximation/entropy bounds.
Their random-design example uses mirror averaging (Section 4.2), and
their final exponentially weighted aggregate is justified for fixed
design (Section 4.1). The one-dimensional analysis here directly
controls the ordinary final raw Gaussian posterior under random design,
uniformly over the stated learner widths. Positive-homogeneous balancing and univariate variation descriptions
are also established tools; Savarese and coauthors~\cite{savarese}
study a related function-space norm with unpenalized hidden biases.
Our Gaussian energy includes those biases, so its bounds are derived
directly here. The elementary midpoint identity, balancing, and finite
atomic representation are not claimed as new.

Entropy-regularized mean-field learning and its generalization theory
are well established, including the convex analysis of Nitanda, Wu,
and Suzuki~\cite{nitanda2022}, Gaussian-reference mean-field learning
of Nishikawa and coauthors~\cite{nishikawa2022}, and the generalization
analysis of Aminian, Cohen, and Szpruch~\cite{aminian2023}.
Nitanda~\cite{nitanda2024} also proves a finite-particle stationary
Gibbs objective bound. Section~\ref{sec:all-d-product-midpoint} concerns
the actual finite-width correlated Gaussian posterior with an
unbounded linear readout and gives a high-probability random-design
oracle. Its proof uses product-subset relative entropy and the
smooth-loss complexity theorem of Srebro, Sridharan, and
Tewari~\cite{sst}; these entropy and empirical-process tools are
established methods.

The proofs have received independent AI-agent mathematical reviews.
They are not Lean formalizations. The verification record distinguishes
the imported Gaussian statistical theorem from the new compiler and
oracle assembly checked for this manuscript.

\appendix
\section{A product-prior midpoint argument for shallow tanh networks}
\label{sec:tanh-product-midpoint}
This result uses the unchanged independent Gaussian prior, a raw squared-loss
likelihood, and the ordinary raw posterior mean. It does not use positive
homogeneity. The midpoint construction requires an explicit lower bound
on the learner width.

Let $X\in[-1,1]$, $|Y|\le1$, and suppose the input density is bounded
above and below by fixed positive constants. Write
\[
 f_w(x)=\frac1N\sum_{j=1}^N G_j\tanh(a_jx+b_j),\qquad
 (a_j,b_j,G_j)\stackrel{\rm iid}{\sim}N(0,\nu I_3),\quad0<\nu\le1.
\]
Thus the physical readout weights $v_j=G_j/N$ have variance $\nu/N^2$.
Set $\omega=N/\nu$ and
\[
 Q(w)=N^{-1}\sum_j(a_j^2+b_j^2+G_j^2),\qquad
 Q_G(w)=N^{-1}\sum_jG_j^2.
\]
The actual posterior is
\[
 d\Pi_\lambda(w)=Z_\lambda^{-1}
 \exp\{-\omega\widehat R(f_w)/(2\lambda)\}\,d\mu(w),
 \qquad \widehat R(f)=n^{-1}\sum_i(f(X_i)-Y_i)^2.
\]
Write $R(f)=\mathbb E(f(X)-Y)^2$ for population squared risk.
Its physical observation variance is $n\lambda\nu/N$.

\begin{lemma}[Replacing a random subset of neuron packs]
\label{lem:tanh-subset-midpoint}
Let $q$ be any law on one neuron pack with finite relative entropy
against $\mu_0=N(0,\nu I_3)$ and finite $M_q=\mathbb E_qG^2$.
Put $g_q(x)=\mathbb E_q[G\tanh(ax+b)]$. For $N\ge2$, let
$k=\lfloor N/2\rfloor$ and $p=k/N\in[1/3,1/2]$. Then
\begin{align}
 &\Pi_\lambda[\widehat R(f)-\widehat R(g_q)]
 +p\Pi_\lambda\|f-g_q\|_n^2
 +\frac{2\lambda\nu}{N}\operatorname{KL}(\Pi_\lambda\Vert\mu)
 \nonumber\\
 &\hspace{12mm}\le
 2\lambda\nu\operatorname{KL}(q\Vert\mu_0)
 +\frac{p}{N-1}\Pi_\lambda Q_G+\frac{M_q}{N}.
 \tag{TM1}
\end{align}
\end{lemma}
\begin{proof}
Choose uniformly a $k$-subset $I$ of $\{1,\ldots,N\}$. Keep the
packs in $I$ from a posterior draw and replace the other packs by
independent draws from $q$. Call the resulting law $\rho$.
The product-measure Shearer inequality gives the averaged subset
entropy bound; see, for example, Caputo and Parisi~\cite{entropy-shearer},
Lemma 4.2. The following chain-rule proof supplies the continuous-state
form for these Gaussian packs:

\[
 \mathbb E_I\operatorname{KL}((\Pi_\lambda)_I\Vert\mu_0^{\otimes I})
 \le p\operatorname{KL}(\Pi_\lambda\Vert\mu).
\]
Indeed, expand each marginal entropy by its coordinate chain rule.
Convexity shows that conditioning a coordinate on only the preceding
coordinates belonging to $I$ costs no more, on average, than conditioning
on all preceding coordinates. Each coordinate belongs to $I$ with
probability $p$. Sum these inequalities. Product additivity conditional
on $I$, followed by convexity over $I$, gives
\[
 \operatorname{KL}(\rho\Vert\mu)
 \le p\operatorname{KL}(\Pi_\lambda\Vert\mu)
 +(1-p)N\operatorname{KL}(q\Vert\mu_0).                 \tag{TM2}
\]
Conditional on the original packs, the mean new output is
$pf+(1-p)g_q$. The variance formula for a sample sum without replacement
and $|G\tanh(ax+b)|\le|G|$ give
\[
 \operatorname{Var}(f_{\rm new}(x)\mid w)
 \le \frac{p(1-p)}{N-1}Q_G(w)+\frac{1-p}{N}M_q.
\]
The squared-loss identity therefore implies
\begin{align*}
 \rho\widehat R(f)\le{}&p\Pi_\lambda\widehat R(f)
 +(1-p)\widehat R(g_q)-p(1-p)\Pi_\lambda\|f-g_q\|_n^2\\
 &+\frac{p(1-p)}{N-1}\Pi_\lambda Q_G+\frac{1-p}{N}M_q.
\end{align*}
Use the minimizing property of the actual Gibbs posterior for
$\rho\widehat R+(2\lambda\nu/N)\operatorname{KL}(\rho\Vert\mu)$,
insert (TM2), and divide by $1-p$.
\end{proof}

\begin{lemma}[A bounded-coefficient comparison pack]
\label{lem:tanh-teacher-pack}
Let
\[
 g(x)=\sum_{i=1}^C c_i\tanh(\alpha_i x+\beta_i),\qquad
 |c_i|,|\alpha_i|,|\beta_i|\le B,\quad B\ge1,
 \qquad \|g\|_\infty\le1.
\]
For $0<\sigma<1$, take a uniform index $i\in\{1,\ldots,C\}$,
then independently draw
\[
 (a,b)\sim N((\alpha_i,\beta_i),\sigma^2 I_2),\qquad
 G\sim N(Cc_i,\nu).
\]
Its mean function $g_\sigma$ is independent of $\nu$ and satisfies
\[
 \|g_\sigma-g\|_\infty\le2CB\sigma,\qquad M_q\le C^2B^2+1,
\]
\[
 \nu\operatorname{KL}(q\Vert\mu_0)
 \le C_0\{C^2B^2+1+\log(1/\sigma)\}.                    \tag{TM3}
\]
Its Lipschitz constant and second-derivative variation are at most
$C_0(CB^2+1)$.
\end{lemma}
\begin{proof}
The mean of $G$ gives a Gaussian smoothing of each teacher ridge.
Lipschitz continuity of tanh and Cauchy--Schwarz bound the uniform
error by $2CB\sigma$. The second moment is immediate. Convexity bounds
the mixture entropy by the mean component entropy. The exact Gaussian
formula, multiplied by $\nu$, is
\[
 \tfrac12\{\alpha_i^2+\beta_i^2+C^2c_i^2
 +2\sigma^2-2\nu-2\nu\log(\sigma^2/\nu)\}.
\]
Use $\nu\log\nu\le0$ and $0<\nu\le1$.
Finally $\int_{\mathbb R}|\tanh''u|\,du=2$ gives
$\|[c\tanh(ax+b)]''\|_{\rm TV}\le2|ca|$ on the input interval;
integrate this and the first-derivative bound over the smoothing.
\end{proof}

\begin{proposition}[Strong empirical comparison]
\label{prop:tanh-strong-objective}
Set
\[
 A_n=1+C^2B^2+\log(n(CB+1)),\qquad
 \sigma=\frac{n^{-20}}{CB+1}.
\]
There are absolute constants $c_0,C_0>0$ such that, whenever
$0<\lambda\le1$ and $N\lambda\ge c_0$,
\[
 \Pi_\lambda[\widehat R(f)-\widehat R(g_\sigma)]
 +\tfrac13\Pi_\lambda\|f-g_\sigma\|_n^2
 +c_0^{-1}\lambda\Pi_\lambda Q
 \le C_0\lambda A_n.                                   \tag{TM4}
\]
\end{proposition}
\begin{proof}
Since $\omega Q\sim\chi^2_{3N}$, the entropy inequality gives
\[
 \operatorname{KL}(\Pi_\lambda\Vert\mu)
 \ge\frac{\omega}{4}\Pi_\lambda Q-\frac{3N}{2}\log2.
\]
Insert this and (TM3) in (TM1). The variance term is at most
$C\Pi_\lambda Q/N$, absorbed by half the resulting $\lambda Q/2$
penalty when $N\lambda$ is large enough. The remaining $M_q/N$
term is at most $C\lambda A_n$.
\end{proof}

\paragraph{The bounded-variation statistical interface.}
The univariate localized squared-loss estimates apply to every continuous
function whose supremum, Lipschitz constant, and second-derivative
variation are bounded by a common radius. They apply here because
\[
 \|f_w\|_\infty\le\sqrt{Q_G},\qquad
 \operatorname{Lip}(f_w)\le Q/2,\qquad
 \|f_w''\|_{\rm TV}\le Q.                               \tag{TM5}
\]
Cauchy--Schwarz proves the first bound; $2|Ga|\le G^2+a^2$ and the
preceding tanh integral prove the other two. Thus use radius
$C(1+Q+A_n)$ in the same bounded-variation class.
Specifically, at a fixed small tolerance, the localized multiplier
and design-norm inequalities give, uniformly in $f$,
\begin{align}
 R(f)-R(g_\sigma)\le{}&
 \widehat R(f)-\widehat R(g_\sigma)
 +\tfrac16\|f-g_\sigma\|_n^2+c\lambda Q+C\lambda A_n\nonumber\\
 &+\widetilde Cn^{-4/3}\lambda^{-2/3}
 +\widetilde Cn^{-3}(Q+A_n)^2
 +\widetilde Cn^{-3}\lambda^{-1}+\widetilde Cn^{-1}.       \tag{TM6}
\end{align}
Choose $c$ smaller than the penalty constant in (TM4).
The hypotheses are $A_n\le n^4$ and
$\lambda\ge\widetilde Cn^{-13/8}$.
This is a function-class statement, not a ReLU-specific step.
Its entropy is
$\log N_\infty(\epsilon)\lesssim\sqrt{P/\epsilon}\log(CP/\epsilon)$;
the noise multiplier costs $\widetilde O(P^{2/5}n^{-4/5})$;
the localized design multiplier adds
$\widetilde O(Pn^{-13/8}+P^{1/2}n^{-3/2})$; and the two-sided
design norm comparison costs $\widetilde O(P^2n^{-3})$.
Young's inequality against $\lambda P$ gives (TM6).
These are the localized bounded-variation and design-norm lemmas
of Sections~\ref{sec:sharp-design} and~\ref{sec:localized-proof}. The comparator is bounded by
$1+2n^{-20}$, changing only absolute constants. Keeping the quadratic
distance, rather than replacing it by approximation risks, is what
preserves the exact leading comparator coefficient.

\begin{theorem}[A sharp raw tanh oracle at sufficiently large width]
\label{thm:tanh-midpoint-oracle}
Fix a deterministic comparator $g$ as above. For
$n^{-2}\le\lambda\le1$ with $\lambda\ge\widetilde Cn^{-13/8}$
and $A_n\le n^4$, on an event of probability at least $1-\delta$,
simultaneously for all $0<\nu\le1$ and all $N\ge c_0/\lambda$,
\begin{align}
 R(\mathbb E_{\Pi_\lambda}f)\le\Pi_\lambda R(f)
 \le R(g)+\widetilde O\left(
 \lambda A_n+n^{-4/3}\lambda^{-2/3}
 +n^{-3}\lambda^{-2}+n^{-3}A_n^2+n^{-1}\right).           \tag{TM7}
\end{align}
Constants depend only on the two input-density bounds. The tildes
suppress fixed powers of logarithms in $n,1/\delta,C,B$.
\end{theorem}
\begin{proof}
The native empirical risk is at most two since tanh is bounded and
$\mathbb E_\mu f_w(x)^2\le\nu/N\le1$. Jensen bounds the free energy
by two. Deterministically,
\[
 \Pi_\lambda Q^2\le C(1+\lambda^{-2}).
\]
Indeed $Z_\lambda\ge e^{-\omega/\lambda}$, whereas the chi-square
Chernoff bound gives $\mu(Q>u)\le e^{-c\omega u}$ above a fixed
threshold. Integrate
$\min(1,e^{\omega/\lambda-c\omega u})$.
Truncating at a sufficiently large polynomial multiple of $1/\lambda$
makes all omitted risk and energy terms $O(n^{-10})$ and reduces the
shell logarithms to the stated logarithms. Integrate (TM6), use (TM4),
discard its remaining nonnegative distance and energy terms, and use
the moment bound. Finally $R(g_\sigma)\le R(g)+Cn^{-20}$ and Jensen
proves the raw-mean claim. Statistical events involve only the common
function classes and the fixed, $\nu$-independent function $g_\sigma$.
\end{proof}

\begin{corollary}[An explicit bounded-coefficient polynomial]
\label{cor:tanh-midpoint-poly}
In the nontrivial range $A_n\le c n^{3/2}$, choose sufficiently large
fixed logarithmic factors in
\[
 \lambda=\widetilde\Theta\left(
 \max\{A_n^{-3/5}n^{-4/5},\ A_n^{-1/3}n^{-1}\}\right),
 \qquad N\ge c_0/\lambda.
\]
Then
\[
 R(\mathbb E_{\Pi_\lambda}f)\le\Pi_\lambda R(f)
 \le R(g)+\widetilde O\left(A_n^{2/5}n^{-4/5}
                         +A_n^{2/3}n^{-1}\right).        \tag{TM8}
\]
At fixed coefficient cap $B$, a sufficient sample size for excess
$\varepsilon$ is
\[
 n\ge\widetilde O\left(C\varepsilon^{-5/4}
                         +C^{4/3}\varepsilon^{-1}\right),             \tag{TM9}
\]
with the displayed learner-width condition. The remainder is
$\widetilde O(C^{4/5}n^{-4/5}+C^{4/3}/n)$.
The same bound holds against the infimum risk over this bounded-output,
coefficient-bounded teacher class: choose a deterministic population-risk
near-minimizer before invoking the probability statement. This does not
assert a common event for every data-selected comparator.
\end{corollary}
\begin{proof}
The first scale balances $\lambda A_n$ against
$n^{-4/3}\lambda^{-2/3}$, and the second balances it against
$n^{-3}\lambda^{-2}$. Taking their maximum bounds both inverse-scale
terms by the displayed sum. Also $n^{-3}A_n^2\le A_n^{2/3}/n$ when
$A_n\le n^{3/2}$. The second scale is at least $n^{-3/2}$, above
the $n^{-13/8}$ floor. At fixed $B$, $A_n=\widetilde O(C^2)$.
\end{proof}

\paragraph{Other bounded activations.}
The same result holds for any fixed bounded globally Lipschitz activation
$\phi$ whose weak first derivative has finite total variation on
$\mathbb R$. Constants may depend on its supremum, Lipschitz constant,
and derivative variation. Indeed the subset argument is unchanged,
its replacement variance is at most
$\|\phi\|_\infty^2[p(1-p)Q_G/(N-1)+(1-p)M_q/N]$,
and
\[
 \operatorname{Lip}(f)\le\tfrac12\operatorname{Lip}(\phi)Q,
 \qquad
 \|f''\|_{\rm TV}\le\tfrac12\operatorname{TV}(\phi')Q.
\]
Gaussian teacher smoothing uses only Lipschitz continuity. Thus logistic
sigmoid and hardtanh are also covered, with the same exponents and
with $N\lambda$ exceeding an activation-dependent constant.
No assumption $\phi(0)=0$ is needed. This paragraph does not extend
the midpoint variance bound to unbounded ReLU or GELU activations.

\begin{corollary}[Marginally sub-Gaussian labels]
\label{cor:tanh-midpoint-subgaussian}
Suppose instead that $\E\exp(Y^2/s_Y^2)\le2$. Keep the same teacher
class, including $\|g\|_\infty\le1$ and its coefficient cap $B$.
The powers of $A_n,n$ and the fixed-$B$ sample polynomial in
Corollary~\ref{cor:tanh-midpoint-poly} remain valid up to additional
logarithmic factors. Define
\[
 M_n=\max\{1,s_Y\sqrt{c\log(n/\delta)}\},\qquad
 \lambda=M_n^2\lambda',
\]
where $\lambda'$ is the displayed optimized schedule with sufficiently
large logarithmic factors. The sufficient learner-width condition is
$N\lambda'\ge c_0$, and the physical observation variance is
$nM_n^2\lambda'\nu/N$. The event remains common to all such $N$
and every $0<\nu\le1$. The bounded-activation extension has the
same conclusion with activation-dependent constants, including its
width threshold. Observed labels, the prior and raw predictions
are unchanged.
\end{corollary}
\begin{proof}
Clip labels to $Y_M$ only in the analysis, keeping the original input
law, and study $f/M_n$ against $Y_M/M_n$. Use the same physical
comparison pack as in Lemma~\ref{lem:tanh-teacher-pack}. Its normalized
mean is $g_\sigma/M_n$, still independent of $\nu$. In (TM1), the
replacement variance is divided by $M_n^2$, while the entropy terms
have the same coefficients with $\lambda'$ in place of $\lambda$.
Thus $N\lambda'\ge c_0$ suffices for the strong objective inequality,
with path-energy penalty proportional to $\lambda'\Pi Q$ and
comparison resource $C\lambda'A_n$. The normalized function bounds
in (TM5) only improve. Its prior empirical risk is at most two for tanh, and at most a fixed
activation-dependent constant in the extension, so the same moment estimate
$\Pi Q^2\le C(1+(\lambda')^{-2})$ holds.

Apply the fixed-tolerance statistical proof and multiply normalized
risks by $M_n^2$. This changes only logarithmic factors and preserves
the leading coefficient one. The smoothing error of the bounded
physical teacher is negligible at the chosen $\sigma$. On the event
that all training labels have magnitude at most $M_n$, this is exactly
the actual posterior. For the population transfer use
$|f(x)|^2\le\|\phi\|_\infty^2 Q_G\le\|\phi\|_\infty^2 Q$
(with $\|\phi\|_\infty=1$ for tanh) and Cauchy--Schwarz; the raw risk difference
between $Y$ and $Y_M$ is at most
\[
 C(1+(\lambda')^{-1})\{\E(Y-Y_M)^2\}^{1/2}
       +\E|Y^2-Y_M^2|.
\]
A sufficiently large fixed $c$ makes this negligible for all the
stated polynomial schedules. The bounded comparator's corresponding
risk transfer is smaller. The input density never changes, and the
common-width and common-scale event follows as before.
\end{proof}


\begin{thebibliography}{99}
\bibitem{gaussian} ChatGPT-6 Astra, Claude Opus 5.5, Kaarel H\"anni, and
Dmitry Vaintrob. \emph{Finite Bayesian neural networks as agnostic
approximation learners}. Research manuscript, November 2026. Finite
all-width raw-output Gaussian oracle and evidence calibration.
\bibitem{tails} ChatGPT-6 Astra, Claude Opus 5.5, Kaarel H\"anni, and
Dmitry Vaintrob. \emph{Sub-Gaussian labels without changing the parameter
prior}. Proof module, November 2026. Exact raw-output label-tail transfer.
\bibitem{adaptive} ChatGPT-6 Astra, Claude Opus 5.5, Kaarel H\"anni,
and Dmitry Vaintrob. \emph{Adaptive approximation oracles for finite
Bayesian neural networks}. Research manuscript, November 2026.
Gaussian free-energy certificate and shallow ReLU path-variation oracle.
\bibitem{boolean} ChatGPT-6 Astra, Claude Opus 5.5, Kaarel H\"anni, and
Dmitry Vaintrob. \emph{A width-only approximation guarantee for Bayesian
neural networks on Boolean inputs}. Research manuscript, November 2026.
Coefficient-free finite-domain majority-circuit compiler.
\bibitem{bach-convex} Francis Bach. Breaking the Curse of Dimensionality with Convex Neural Networks. \emph{Journal of Machine Learning Research} 18(19), 1--53, 2017. \url{https://jmlr.org/papers/v18/14-546.html}.
\bibitem{tinsi-dalalyan} Laura Tinsi and Arnak Dalalyan. Risk bounds for aggregated shallow neural networks using Gaussian priors. \emph{Proceedings of Machine Learning Research} 178, 227--253, 2022. \url{https://proceedings.mlr.press/v178/tinsi22a.html}.
\bibitem{savarese} Pedro Savarese, Itay Evron, Daniel Soudry, and Nathan Srebro. How do infinite width bounded norm networks look in function space? \emph{Proceedings of Machine Learning Research} 99, 2667--2690, 2019. \url{https://proceedings.mlr.press/v99/savarese19a.html}.
\bibitem{nitanda2022} Atsushi Nitanda, Denny Wu, and Taiji Suzuki. Convex Analysis of the Mean Field Langevin Dynamics. \emph{Proceedings of Machine Learning Research} 151, 9741--9757, 2022. \url{https://proceedings.mlr.press/v151/nitanda22a.html}.
\bibitem{nishikawa2022} Naoki Nishikawa, Taiji Suzuki, Atsushi Nitanda, and Denny Wu. Two-layer neural network on infinite-dimensional data: global optimization guarantee in the mean-field regime. \emph{Advances in Neural Information Processing Systems} 35, 2022. \url{https://proceedings.neurips.cc/paper_files/paper/2022/file/d2155b1f7eb42350d7bc3013eefe5480-Paper-Conference.pdf}.
\bibitem{aminian2023} Gholamali Aminian, Samuel N. Cohen, and {\L}ukasz Szpruch. Mean-field Analysis of Generalization Errors. Preprint, 2023. \url{https://arxiv.org/abs/2306.11623}.
\bibitem{nitanda2024} Atsushi Nitanda. Improved Particle Approximation Error for Mean Field Neural Networks. \emph{Advances in Neural Information Processing Systems} 37, 2024. \url{https://arxiv.org/abs/2405.15767}.
\bibitem{leshno} Moshe Leshno, Vladimir Ya. Lin, Allan Pinkus, and Shimon Schocken. Multilayer feedforward networks with a nonpolynomial activation function can approximate any function. \emph{Neural Networks} 6(6), 861--867, 1993. \url{https://doi.org/10.1016/S0893-6080(05)80131-5}.
\bibitem{bousquet} Olivier Bousquet. A Bennett concentration inequality and its application to suprema of empirical processes. \emph{Comptes Rendus. Math\'ematique} 334(6), 495--500, 2002. \url{https://doi.org/10.1016/S1631-073X(02)02292-6}.
\bibitem{sst} Nathan Srebro, Karthik Sridharan, and Ambuj Tewari. Smoothness, low noise and fast rates. \emph{Advances in Neural Information Processing Systems} 23, 2010. \url{https://www.cs.cornell.edu/~sridharan/smoothness.pdf}.
\bibitem{entropy-shearer} Pietro Caputo and Daniel Parisi. Block Factorization of the Relative Entropy via Spatial Mixing. \emph{Communications in Mathematical Physics} 388, 793--818, 2021. \url{https://doi.org/10.1007/s00220-021-04237-1}.
\end{thebibliography}
\end{document}
